% El turismo y el desarrollo económico — Espagnol 1ère A — Cours signé OnBuch+
% Programme officiel MINESEC. Compiler avec : tectonic EA13-turismo-desarrollo-economico.tex
\newcommand{\DOCMATIERE}{Espagnol}
\newcommand{\DOCNIVEAU}{1ère A}
\newcommand{\DOCMODULE}{Módulo 4 : Vie économique}
\newcommand{\DOCLECON}{EA13}
\newcommand{\DOCTITRE}{El turismo y el desarrollo económico}
% OnBuch+ — préambule « cours signés » ENRICHI (compilé avec tectonic / XeLaTeX)
% Système visuel « Pop » d'OnBuch+ : fond crème (jamais blanc), bordures encre
% épaisses, ombres « dures » décalées, titres Archivo Black en capitales.
% Version MATHÉMATIQUES : graphes (pgfplots/tikz), boîtes pédagogiques
% (définition, propriété, méthode, attention, le savais-tu, exercice résolu,
% exercice, TP, bilan), page de garde et signature OnBuch+ ancrée en bas.
%
% Macros attendues dans main.tex AVANT \input{preamble} :
%   \DOCMATIERE  \DOCNIVEAU  \DOCMODULE  \DOCLECON (numéro)  \DOCTITRE
\documentclass[11pt]{article}

\usepackage{fontspec}
\usepackage[spanish,french]{babel} % français = langue principale ; l'espagnol via \esp{..} / espagnol
\usepackage[a4paper,margin=2cm,top=3.4cm,bottom=2.8cm,headheight=1.4cm,footskip=1.3cm]{geometry}
\usepackage{amsmath,amssymb,amsfonts}
\usepackage{mathtools} % \xmapsto, \coloneqq... souvent utilisés par les modèles
\usepackage{cancel} % simplifications barrées (bilans de forces)
% \abs{z} (module, valeur absolue) et \norm{u} (norme), utilisés par les modèles
\providecommand{\abs}{}\let\abs\relax\DeclarePairedDelimiter{\abs}{\lvert}{\rvert}
\providecommand{\norm}{}\let\norm\relax\DeclarePairedDelimiter{\norm}{\lVert}{\rVert}
\usepackage[table]{xcolor} % [table] : fournit aussi \rowcolors (alternance de lignes)
\usepackage{pagecolor}
\usepackage{tikz}
\usetikzlibrary{arrows.meta,positioning,calc,shapes.geometric,shapes.misc,shapes.arrows,decorations.pathmorphing,decorations.pathreplacing,decorations.markings,patterns,shadows,mindmap,trees,fit,backgrounds,matrix,intersections,angles,quotes}
\usepackage{pgfplots}
\usepackage[version=4]{mhchem} % équations nucléaires \ce{^{235}_{92}U}
\usepackage[european,siunitx]{circuitikz} % schémas électriques
\pgfplotsset{compat=1.18}
\usepgfplotslibrary{fillbetween,groupplots} % aires entre courbes (calcul intégral)
\usepackage{enumitem}
\usepackage{fancyhdr}
% [most] charge toutes les bibliothèques tcolorbox (listings, minted, raster,
% documentation...) et gonfle inutilement la mémoire TeX sur un document long
% (~300 boîtes) : on ne charge que ce qui est réellement utilisé.
\usepackage[skins,breakable]{tcolorbox}
\usepackage{graphicx}
\usepackage{colortbl}
\usepackage{booktabs}
\usepackage{tabularx}
\usepackage{diagbox} % case d'en-tête diagonale des tableaux à double entrée
% Colonnes extensibles centrée / gauche / droite (C, L, R), souvent utilisées par les modèles
\newcolumntype{C}{>{\centering\arraybackslash}X}
\newcolumntype{L}{>{\raggedright\arraybackslash}X}
\newcolumntype{R}{>{\raggedleft\arraybackslash}X}
\usepackage{array}
\usepackage{multicol}
\usepackage{siunitx}
% parse-numbers=false : accepte aussi \SI{2,5 \times 10^{-3}}{...}
\sisetup{locale=FR,output-decimal-marker={,},group-minimum-digits=4,parse-numbers=false}
\DeclareSIUnit\ohm{\ensuremath{\symup{\Omega}}} % U+2126 absent de la police
\DeclareSIUnit\division{div} % réglages d'oscilloscope (ms/div, V/div)
\DeclareSIUnit\tr{tr} % tours (vitesse de rotation, ex. \SI{1500}{\tr\per\minute})
\DeclareSIUnit\molar{M} % concentration molaire (ex. \SI{3}{\molar}, \SI{1}{\micro\molar})
\DeclareSIUnit\an{an} % année (le modèle confond parfois avec \year, un compteur TeX) — ex. \SI{5}{\kilo\meter\per\an}
\DeclareSIUnit\FCFA{FCFA} % franc CFA (ex. \SI{500}{\FCFA\per\kilo\gram})
\DeclareSIUnit\degreeBrix{\textdegree Brix} % teneur en sucre d'un sirop/jus (réfractométrie)
\DeclareSIUnit\month{mois} % le modèle utilise parfois \month, un compteur TeX, comme unité de durée
\usepackage{qrcode}
% \degree : commande fréquemment utilisée par les modèles pour un angle ou une
% température en dehors de \SI{}{} (non fournie par les paquets chargés).
\providecommand{\degree}{^{\circ}}
% \qed (fin de démonstration), utilisé par les modèles sans amsthm
\providecommand{\qed}{\hfill$\square$}
% \barre : double barre des valeurs interdites dans un tableau de variations (tkz-tab)
\providecommand{\barre}{\ensuremath{\|}}
% environnement proof (amsthm non chargé) utilisé par les modèles
\ifdefined\proof\else\newenvironment{proof}[1][Démonstration]{\par\noindent\textit{#1.}\ }{\qed\par}\fi
\usepackage{lastpage}
\usepackage{hyperref}
\hypersetup{hidelinks}

% ============================================================
% Palette « Pop » OnBuch+ — mêmes hex que la classe P (plus_ui.dart)
% ============================================================
\definecolor{poporange}{HTML}{F59321}
\definecolor{popdark}{HTML}{E07A0C}
\definecolor{popcream}{HTML}{FFF6E8}
\definecolor{popink}{HTML}{1C1714}
\definecolor{poppurple}{HTML}{7A5AE0}
\definecolor{popgreen}{HTML}{2E9E6A}
\definecolor{poppink}{HTML}{E0457B}
\definecolor{popblue}{HTML}{2F7DE1}
\definecolor{popgold}{HTML}{C98A12}
\definecolor{popmuted}{HTML}{8A7F76}
\definecolor{popline}{HTML}{EADFD2}
\definecolor{popgray}{HTML}{8A7F76}
\colorlet{popgrey}{popgray}
% Alias tolérants (le modèle nomme parfois les couleurs différemment)
\colorlet{popwhite}{white}
\colorlet{popblack}{popink}
\colorlet{popred}{poppink}
\colorlet{popyellow}{popgold}
% Teintes claires (fonds de tableaux/encadrés) — le modèle nomme parfois
% les couleurs avec un suffixe L (ex. popblueL) sans les avoir définies.
\colorlet{poporangeL}{poporange!20}
\colorlet{popdarkL}{popdark!20}
\colorlet{poppurpleL}{poppurple!20}
\colorlet{popgreenL}{popgreen!20}
\colorlet{poppinkL}{poppink!20}
\colorlet{popblueL}{popblue!20}
\colorlet{popgoldL}{popgold!20}
\colorlet{popredL}{popred!20}
\colorlet{popgrayL}{popgray!20}
% Teintes claires pour fonds de tableaux / zones
\colorlet{popblueL}{popblue!10!white}
\colorlet{poporangeL}{poporange!14!white}
\colorlet{popgreenL}{popgreen!12!white}
\colorlet{poppurpleL}{poppurple!12!white}
\colorlet{poppinkL}{poppink!12!white}
\colorlet{popgoldL}{popgold!14!white}

\pagecolor{popcream}

% ============================================================
% Typographie
% ============================================================
\defaultfontfeatures{Ligatures=TeX}
\setmainfont{PlusJakartaSans-Medium.ttf}[Path=../../../fonts/,BoldFont=PlusJakartaSans-Bold.ttf]
% Maths en sans-serif (Fira Math) avec chiffres et lettres droites de Plus
% Jakarta, pour que les formules se fondent dans le texte.
\usepackage[math-style=ISO,bold-style=ISO]{unicode-math}
\setmathfont{FiraMath-Regular.otf}[Path=../../../fonts/]
\setmathfont{PlusJakartaSans-Medium.ttf}[Path=../../../fonts/,range={up/num,up/{Latin,latin},\mathup}]
\usepackage{icomma} % virgule décimale sans espace en mode math
% Caractères mathématiques Unicode absents des polices de texte (Plus Jakarta,
% Archivo Black) : rendus par leur équivalent mathématique, en texte comme en
% formule — sinon ils s'affichent en carré vide (ex. « Arithmétique dans ℤ »).
\usepackage{newunicodechar}
\newunicodechar{ℝ}{\ensuremath{\mathbb{R}}}
\newunicodechar{ℤ}{\ensuremath{\mathbb{Z}}}
\newunicodechar{ℂ}{\ensuremath{\mathbb{C}}}
\newunicodechar{ℕ}{\ensuremath{\mathbb{N}}}
\newunicodechar{ℚ}{\ensuremath{\mathbb{Q}}}
\newunicodechar{𝔻}{\ensuremath{\mathbb{D}}}
\newunicodechar{α}{\ensuremath{\alpha}}
\newunicodechar{β}{\ensuremath{\beta}}
\newunicodechar{γ}{\ensuremath{\gamma}}
\newunicodechar{δ}{\ensuremath{\delta}}
\newunicodechar{ε}{\ensuremath{\varepsilon}}
\newunicodechar{θ}{\ensuremath{\theta}}
\newunicodechar{λ}{\ensuremath{\lambda}}
\newunicodechar{μ}{\ensuremath{\mu}}
\newunicodechar{σ}{\ensuremath{\sigma}}
\newunicodechar{τ}{\ensuremath{\tau}}
\newunicodechar{φ}{\ensuremath{\varphi}}
\newunicodechar{ψ}{\ensuremath{\psi}}
\newunicodechar{ω}{\ensuremath{\omega}}
\newunicodechar{Σ}{\ensuremath{\Sigma}}
% majuscules produites par \MakeUppercase dans les titres (α-aminés -> Α)
\newunicodechar{Α}{\ensuremath{\alpha}}
\newunicodechar{Β}{\ensuremath{\beta}}
\newunicodechar{Γ}{\ensuremath{\Gamma}}
\newunicodechar{Δ}{\ensuremath{\Delta}}
\newunicodechar{Ε}{\ensuremath{\varepsilon}}
\newunicodechar{Θ}{\ensuremath{\Theta}}
\newunicodechar{Λ}{\ensuremath{\Lambda}}
\newunicodechar{Μ}{\ensuremath{\mu}}
\newunicodechar{Τ}{\ensuremath{\tau}}
\newunicodechar{Φ}{\ensuremath{\Phi}}
\newunicodechar{Ψ}{\ensuremath{\Psi}}
\newunicodechar{Ω}{\ensuremath{\Omega}}
\newunicodechar{↦}{\ensuremath{\mapsto}}
\newunicodechar{∈}{\ensuremath{\in}}
\newunicodechar{∉}{\ensuremath{\notin}}
\newunicodechar{∧}{\ensuremath{\wedge}}
\newunicodechar{⇔}{\ensuremath{\Leftrightarrow}}
\newunicodechar{⟺}{\ensuremath{\Longleftrightarrow}}
\newunicodechar{⇒}{\ensuremath{\Rightarrow}}
\newunicodechar{⟹}{\ensuremath{\Longrightarrow}}
\newunicodechar{∘}{\ensuremath{\circ}}
\newunicodechar{∪}{\ensuremath{\cup}}
\newunicodechar{∩}{\ensuremath{\cap}}
\newunicodechar{≡}{\ensuremath{\equiv}}
\newunicodechar{⊂}{\ensuremath{\subset}}
\newunicodechar{⊥}{\ensuremath{\perp}}
\newunicodechar{∃}{\ensuremath{\exists}}
\newunicodechar{∀}{\ensuremath{\forall}}
\newunicodechar{∛}{\ensuremath{\sqrt[3]{\,}}}
\newunicodechar{∜}{\ensuremath{\sqrt[4]{\,}}}
\newunicodechar{⃗}{\ensuremath{{}^{\to}}}
\newunicodechar{ⁿ}{\ifmmode^{n}\else\textsuperscript{\ensuremath{n}}\fi}
\newunicodechar{ˣ}{\ifmmode^{x}\else\textsuperscript{\ensuremath{x}}\fi}
\newunicodechar{ᵇ}{\ifmmode^{b}\else\textsuperscript{\ensuremath{b}}\fi}
\newunicodechar{ʳ}{\ifmmode^{r}\else\textsuperscript{\ensuremath{r}}\fi}
\newunicodechar{ᵘ}{\ifmmode^{u}\else\textsuperscript{\ensuremath{u}}\fi}
\newunicodechar{ʸ}{\ifmmode^{y}\else\textsuperscript{\ensuremath{y}}\fi}
\newunicodechar{ᵃ}{\ifmmode^{a}\else\textsuperscript{\ensuremath{a}}\fi}
\newunicodechar{ᵉ}{\ifmmode^{e}\else\textsuperscript{\ensuremath{e}}\fi}
\newunicodechar{ᵏ}{\ifmmode^{k}\else\textsuperscript{\ensuremath{k}}\fi}
\newunicodechar{ᵖ}{\ifmmode^{p}\else\textsuperscript{\ensuremath{p}}\fi}
\newunicodechar{ᴺ}{\ifmmode^{N}\else\textsuperscript{\ensuremath{N}}\fi}
\newunicodechar{ᶜ}{\ifmmode^{c}\else\textsuperscript{\ensuremath{c}}\fi}
\newunicodechar{ʰ}{\ifmmode^{h}\else\textsuperscript{\ensuremath{h}}\fi}
\newunicodechar{ᵠ}{\ifmmode^{q}\else\textsuperscript{\ensuremath{q}}\fi}
\newunicodechar{ₙ}{\ifmmode_{n}\else\textsubscript{\ensuremath{n}}\fi}
\newunicodechar{ₐ}{\ifmmode_{a}\else\textsubscript{\ensuremath{a}}\fi}
\newunicodechar{ᵢ}{\ifmmode_{i}\else\textsubscript{\ensuremath{i}}\fi}
\newunicodechar{ᵣ}{\ifmmode_{r}\else\textsubscript{\ensuremath{r}}\fi}
\newunicodechar{ₖ}{\ifmmode_{k}\else\textsubscript{\ensuremath{k}}\fi}
\newunicodechar{ᵦ}{\ifmmode_{\beta}\else\textsubscript{\ensuremath{\beta}}\fi}
\newunicodechar{ₓ}{\ifmmode_{x}\else\textsubscript{\ensuremath{x}}\fi}
\newfontfamily\popdisplay{ArchivoBlack-Regular.ttf}[Path=../../../fonts/]
\newfontfamily\poplogo{SpaceGrotesk-Bold.ttf}[Path=../../../fonts/]
\newfontfamily\popsemi{PlusJakartaSans-SemiBold.ttf}[Path=../../../fonts/]
\newfontfamily\popxbold{PlusJakartaSans-ExtraBold.ttf}[Path=../../../fonts/]
\newfontfamily\popxboldls{PlusJakartaSans-ExtraBold.ttf}[Path=../../../fonts/,LetterSpace=3.0]

\setlist[enumerate]{leftmargin=1.7em,itemsep=5pt,topsep=4pt}
\setlist[itemize]{leftmargin=1.7em,itemsep=4pt,topsep=4pt,label={\color{poporange}\textbullet}}
\setlength{\parindent}{0pt}
\setlength{\parskip}{5pt}
\renewcommand{\arraystretch}{1.25}

% pgfplots : style Pop par défaut
\pgfplotsset{
  every axis/.append style={axis line style={popink,line width=1pt},tick style={popink},
    grid style={popline,line width=0.5pt},label style={font=\small},tick label style={font=\footnotesize}},
  popcurve/.style={poporange,line width=1.8pt,no markers,smooth},
}
% Même style disponible dans un \draw TikZ simple (hors pgfplots)
\tikzset{popcurve/.style={poporange,line width=1.8pt}}
\tikzset{popgrid/.style={popline,very thin}}
\colorlet{popcurve}{poporange} % le nom du style est parfois utilisé comme couleur

% ============================================================
% Pill / squiggle
% ============================================================
\newcommand{\poppill}[3]{%
  \tikz[baseline=-0.55ex]\node[draw=popink,line width=1.2pt,rounded corners=2.6pt,%
    fill=#2,inner xsep=6.5pt,inner ysep=3pt]{%
    \popxboldls\fontsize{7}{7}\selectfont\color{#3}\MakeUppercase{#1}};%
}
\newcommand{\popsquiggle}[1]{%
  \tikz[baseline=-0.4ex]{
    \draw[#1,line width=2.6pt,line cap=round]
      plot [smooth] coordinates {(0,0.09) (0.34,0.01) (0.68,0.21) (1.02,0.01) (1.36,0.10)};
  }%
}
% Mot-clé surligné (marqueur orange sous le texte)
\newcommand{\cle}[1]{{\popxbold\color{popdark}#1}}

% ============================================================
% En-tête / pied
% ============================================================
\pagestyle{fancy}
\fancyhf{}
\renewcommand{\headrulewidth}{0pt}
\renewcommand{\footrulewidth}{0pt}
\lhead{\raisebox{-0.15ex}{\poplogo\fontsize{13}{13}\selectfont\color{popink}OnBuch\color{poporange}+}}
\rhead{\poppill{\DOCMATIERE\ \textbullet\ \DOCNIVEAU\ \textbullet\ Leçon \DOCLECON}{white}{popink}}
\lfoot{\popxbold\fontsize{7.6}{7.6}\selectfont\color{popmuted}\MakeUppercase{Cours signé OnBuch+}\ \textbullet\ {\popsemi\DOCTITRE}}
\rfoot{\tikz[baseline=-0.6ex]\node[rounded corners=8.5pt,fill=popink,minimum height=17pt,minimum width=17pt,inner xsep=5pt,inner sep=0]{\popxbold\fontsize{8}{8}\selectfont\color{popcream}\thepage\,/\,\pageref*{LastPage}};}

% ============================================================
% Titres
% ============================================================
\newcommand{\chaptitre}[2]{%
  \begin{tcolorbox}[enhanced,colback=poporange,colframe=popink,boxrule=2.6pt,arc=11pt,
      drop shadow=popink,left=15pt,right=15pt,top=12pt,bottom=11pt,before skip=3pt,after skip=18pt]
    \poppill{#2}{popink}{popcream}\par\vspace{9pt}
    {\raggedright\hyphenpenalty=10000\exhyphenpenalty=10000\popdisplay\fontsize{22}{24}\selectfont\color{popcream}\MakeUppercase{#1}\par}
  \end{tcolorbox}
}
% Section numérotée : pastille orange avec le numéro + titre Archivo
\newcounter{popsec}
\newcommand{\coursec}[1]{%
  \stepcounter{popsec}\setcounter{popsub}{0}%
  \par\vspace{16pt}\noindent
  \begin{minipage}{\linewidth}
  \tikz[baseline=-0.7ex]\node[draw=popink,line width=1.6pt,fill=poporange,rounded corners=4pt,minimum size=22pt,inner sep=0]{\popdisplay\fontsize{12}{12}\selectfont\color{popcream}\thepopsec};%
  \hspace{8pt}{\popdisplay\fontsize{15}{17}\selectfont\color{popink}\MakeUppercase{#1}}\par
  \vspace{4pt}\noindent{\color{poporange}\rule{36pt}{3.6pt}}
  \end{minipage}\par\nopagebreak\vspace{8pt}
}
\newcounter{popsub}
\newcommand{\courssub}[1]{%
  \stepcounter{popsub}%
  \par\vspace{9pt}\noindent{\popxbold\fontsize{11.5}{13}\selectfont\color{popink}{\color{poporange}\thepopsec.\thepopsub}\hspace{6pt}#1}\par\nopagebreak\vspace{4pt}
}

% ============================================================
% Boîtes pédagogiques — même recette : fond blanc, bordure épaisse colorée,
% ombre dure encre, badge plein. Titre optionnel [#1].
% ============================================================
\tcbset{popbase/.style={breakable,enhanced,colback=white,boxrule=2.3pt,arc=9pt,
  shadow={3pt}{-3pt}{0pt}{popink},left=14pt,right=14pt,top=11pt,bottom=11pt,before skip=11pt,after skip=15pt,
  fontupper=\popsemi\fontsize{10.6}{14.5}\selectfont\color{popink},
  fontlower=\popsemi\fontsize{10.6}{14.5}\selectfont\color{popink}}}
% Coche dessinée (le glyphe ✓ n'existe pas dans les polices du document)
\newcommand{\popcheck}[1]{\tikz[baseline=-0.5ex]\draw[#1,line width=1.6pt,line cap=round,line join=round](-0.13,0.0)--(-0.04,-0.09)--(0.13,0.1);}
\providecommand{\coche}{\popcheck{popgreen}}
\providecommand{\resultat}[1]{\par\smallskip\noindent\fcolorbox{popgreen}{popgreenL}{\parbox{\dimexpr\linewidth-2\fboxsep-2\fboxrule\relax}{\textbf{Résultat.} #1}}\par\smallskip}
\newcommand{\popboxhead}[3]{\poppill{#1}{#2}{white}\ifx\relax#3\relax\else\hspace{6pt}{\popxbold\fontsize{10.6}{12}\selectfont\color{popink}#3}\fi\par\vspace{7pt}}

\newtcolorbox{aretenir}[1][]{popbase,colframe=popgreen,before upper={\popboxhead{À retenir}{popgreen}{#1}}}
% Texte en espagnol (document de lecture, dialogue, poème, extrait) : cadre
% d'étude. Un titre optionnel (source, auteur) s'affiche dans l'en-tête.
\newtcolorbox{textebox}[1][]{popbase,colframe=popblue,before upper={\popboxhead{Texto}{popblue}{#1}\begin{otherlanguage}{spanish}},after upper={\end{otherlanguage}}}
% Marques ✓ / ✗ (forme correcte / fautive) souvent utilisées par les modèles
\providecommand{\cross}{\textcolor{poppink}{\ensuremath{\times}}}
\providecommand{\xmark}{\cross}\providecommand{\crossmark}{\cross}\providecommand{\cmark}{\ensuremath{\checkmark}}
% Ligne de réponse / justification à compléter (inventée par les modèles)
\providecommand{\justification}[1]{\par\noindent\textit{Justification~:} \dotfill\par}
% \textipa (package tipa non chargé) : la police ne couvre pas l'API -> simple passage
\providecommand{\textipa}[1]{#1}
% Soulignement (ulem non chargé) : \uline, \ul
\providecommand{\uline}[1]{\underline{#1}}\providecommand{\ul}[1]{\underline{#1}}
% Ordinaux français écrits en macro par les modèles : 1\ière, 3\ième
\newcommand{\ière}{\textsuperscript{re}}\newcommand{\ième}{\textsuperscript{e}}
% Mot ou phrase en espagnol à l'intérieur d'un paragraphe français
\newcommand{\esp}[1]{\ifmmode\text{\textit{#1}}\else\begingroup\selectlanguage{spanish}\itshape #1\endgroup\fi}
\newtcolorbox{exemplebox}[1][]{popbase,colframe=poporange,before upper={\popboxhead{Exemple}{poporange}{#1}}}
\newtcolorbox{definition}[1][]{popbase,colframe=popblue,before upper={\popboxhead{Définition}{popblue}{#1}}}
\newtcolorbox{propriete}[1][]{popbase,colframe=poppurple,before upper={\popboxhead{Propriété}{poppurple}{#1}}}
\newtcolorbox{methode}[1][]{popbase,colframe=popgold,before upper={\popboxhead{Méthode}{popgold}{#1}}}
\newtcolorbox{attention}[1][]{popbase,colframe=poppink,before upper={\popboxhead{Attention}{poppink}{#1}}}
\newtcolorbox{experience}[1][]{popbase,colframe=popblue,colback=popblueL,before upper={\popboxhead{Expérience}{popblue}{#1}}}
% « Le savais-tu ? » : carte violette pleine (contexte, culture, Cameroun)
\newtcolorbox{savaistu}[1][]{popbase,colback=poppurple,colframe=popink,
  fontupper=\popsemi\fontsize{10.4}{14.2}\selectfont\color{white},
  before upper={\poppill{Le savais-tu\,?}{white}{poppurple}\ifx\relax#1\relax\else\hspace{6pt}{\popxbold\color{white}#1}\fi\par\vspace{7pt}}}
% Exercice résolu : énoncé en haut, \tcblower puis la solution sur fond crème
\newcounter{popexo}\newcounter{popexr}
\newtcolorbox{exoresolu}[1][]{popbase,colframe=poporange,colbacklower=popcream,
  segmentation style={popink,line width=1.2pt,dashed},
  before upper={\stepcounter{popexr}\popboxhead{Exercice résolu \thepopexr}{poporange}{#1}},
  before lower={\poppill{Solution}{popink}{popcream}\par\vspace{6pt}}}
% Exercice (non résolu en place) — la correction est dans la partie Corrigés
\newtcolorbox{exercice}[1][]{popbase,colframe=popink,
  before upper={\stepcounter{popexo}\popboxhead{Exercice \thepopexo}{popink}{#1}}}
% Corrigé
\newtcolorbox{corrige}[1][]{popbase,colframe=popgreen,colback=popgreenL,
  before upper={\popboxhead{Corrigé}{popgreen}{#1}}}
% Objectifs / prérequis / bilan
\newtcolorbox{objectifs}{popbase,colframe=popink,colback=poporangeL,
  before upper={\popboxhead{Objectifs de la leçon}{popink}{}}}
\newtcolorbox{prerequis}{popbase,colframe=popink,colback=popblueL,
  before upper={\popboxhead{Prérequis}{popblue}{}}}
\newtcolorbox{bilan}{popbase,colframe=popink,colback=white,boxrule=2.8pt,
  before upper={\popboxhead{Fiche bilan}{poporange}{L'essentiel à connaître pour l'examen}}}
% Figure encadrée avec légende
\newtcolorbox{popfigure}[1][]{enhanced,breakable=false,colback=white,colframe=popline,boxrule=1.4pt,arc=8pt,
  left=10pt,right=10pt,top=10pt,bottom=8pt,before skip=10pt,after skip=12pt,halign=center,
  lower separated=false,halign lower=center,
  fontlower=\popsemi\fontsize{9}{11.5}\selectfont\color{popmuted},
  #1}
\newcommand{\legende}[1]{\par\vspace{6pt}{\popsemi\fontsize{9}{11.5}\selectfont\color{popmuted}#1}}

% ============================================================
% Signature OnBuch+
% ============================================================
\newcommand{\poplogoinline}[1]{{\poplogo\fontsize{#1}{#1}\selectfont\color{popink}OnBuch\color{poporange}+}}
\newcommand{\signatureonbuch}{%
  \begin{tcolorbox}[enhanced,colback=popink,colframe=popink,boxrule=2.2pt,arc=10pt,
      drop shadow=poporange,left=14pt,right=14pt,top=10pt,bottom=10pt,before skip=0pt,after skip=0pt]
    \tikz[baseline=-0.7ex]\node[circle,fill=popgreen,draw=popcream,line width=1.3pt,minimum size=22pt,inner sep=0]{\popcheck{white}};%
    \hspace{9pt}\begin{minipage}[c]{0.62\linewidth}
      {\popxbold\fontsize{12}{13}\selectfont\color{popcream}Signé\ \textbullet\ {\poplogo OnBuch\color{poporange}+}}\par\vspace{2pt}
      {\raggedright\popsemi\fontsize{8}{10.5}\selectfont\color{popline}\MakeUppercase{Équipe pédagogique OnBuch+}\\\MakeUppercase{Conforme au programme officiel MINESEC}\par}
    \end{minipage}\hfill
    {\poplogo\fontsize{22}{22}\selectfont\color{popcream}OnBuch\color{poporange}+}
  \end{tcolorbox}%
}

% ============================================================
% Page de garde
% #1 titre · #2 matière · #3 niveau · #4 module · #5 n° leçon · #6 durée ·
% #7 description / objectifs (texte ou itemize)
% ============================================================
\newcommand{\pagegarde}[7]{%
  \thispagestyle{empty}
  \newgeometry{a4paper,margin=1.7cm,top=1.6cm,bottom=1.4cm}%
  \begin{tikzpicture}[remember picture,overlay]
    \fill[poporange!18] ([xshift=-3.2cm,yshift=-3.6cm]current page.north east) circle (4.6cm);
    \fill[poppurple!14] ([xshift=2.4cm,yshift=7.6cm]current page.south west) circle (3.8cm);
  \end{tikzpicture}%
  {\poplogo\fontsize{30}{30}\selectfont\color{popink}OnBuch\color{poporange}+}\hfill
  \raisebox{0.6ex}{\poppill{Cours signé \textbullet\ Premium}{popink}{popcream}}\par
  \vspace{1pt}\popsquiggle{poppurple}\par
  \vspace{8pt}
  \begin{tcolorbox}[enhanced,colback=poporange,colframe=popink,boxrule=2.8pt,arc=13pt,
      drop shadow=popink,left=18pt,right=18pt,top=12pt,bottom=13pt,after skip=10pt]
    \poppill{#2 \textbullet\ #3}{popink}{popcream}\hspace{5pt}\poppill{#4}{white}{popink}\par\vspace{8pt}
    {\popdisplay\fontsize{14}{14}\selectfont\color{popink}LEÇON #5}\par\vspace{4pt}
    {\raggedright\hyphenpenalty=10000\exhyphenpenalty=10000\popdisplay\fontsize{25}{27}\selectfont\color{popcream}\MakeUppercase{#1}\par}
    \vspace{8pt}
    {\popxbold\fontsize{9.5}{9.5}\selectfont\color{popink}\MakeUppercase{Durée conseillée : #6}}
  \end{tcolorbox}
  \begin{tcolorbox}[enhanced,colback=white,colframe=popink,boxrule=2.2pt,arc=10pt,
      drop shadow=popink,left=16pt,right=16pt,top=10pt,bottom=10pt,after skip=8pt]
    \poppill{Dans cette leçon}{poppurple}{white}\par\vspace{5pt}
    \popsemi\fontsize{9.3}{12.6}\selectfont\color{popink}#7
  \end{tcolorbox}
  \noindent\begin{minipage}[t]{0.487\linewidth}
    \begin{tcolorbox}[enhanced,colback=popgreen,colframe=popink,boxrule=2.2pt,arc=10pt,
        drop shadow=popink,left=11pt,right=11pt,top=10pt,bottom=10pt,height=3.1cm,valign=center]
      \tcbox[colback=white,colframe=popink,boxrule=1.3pt,arc=5pt,boxsep=0pt,nobeforeafter,
        left=3pt,right=3pt,top=3pt,bottom=3pt,valign=center]{\qrcode[height=1.9cm]{https://play.google.com/store/apps/details?id=com.luvvix.onbuch&hl=fr}}%
      \hspace{8pt}\begin{minipage}[c]{0.5\linewidth}
        {\popdisplay\fontsize{11}{12.5}\selectfont\color{white}\MakeUppercase{Télécharge OnBuch}}\par\vspace{4pt}
        {\raggedright\popsemi\fontsize{8.4}{11}\selectfont\color{white}Gratuit \textbullet\ Google Play\\Tuteur IA, annales, concours, résultats\par}
      \end{minipage}
    \end{tcolorbox}
  \end{minipage}\hfill
  \begin{minipage}[t]{0.487\linewidth}
    \begin{tcolorbox}[enhanced,colback=poppurple,colframe=popink,boxrule=2.2pt,arc=10pt,
        drop shadow=popink,left=11pt,right=11pt,top=10pt,bottom=10pt,height=3.1cm,valign=center]
      \tikz[baseline=-0.7ex]\node[circle,fill=white,draw=popink,line width=1.3pt,minimum size=1.6cm,inner sep=0]{\popdisplay\fontsize{24}{24}\selectfont\color{poppurple}+};%
      \hspace{8pt}\begin{minipage}[c]{0.6\linewidth}
        {\popdisplay\fontsize{11}{12.5}\selectfont\color{white}\MakeUppercase{Passe à OnBuch+}}\par\vspace{4pt}
        {\raggedright\popsemi\fontsize{8.4}{11}\selectfont\color{white}Tous les cours signés, Tuteur IA illimité, fiches PDF — onglet OnBuch+ de l'app\par}
      \end{minipage}
    \end{tcolorbox}
  \end{minipage}
  \vspace{8pt}
  \popcontenu
  \vfill
  \signatureonbuch
  \restoregeometry
  \newpage
}

% Rangée « Ce cours contient » (page de garde)
\newcommand{\poptile}[3]{% #1 couleur #2 gros texte #3 légende
  \begin{tcolorbox}[enhanced,colback=#1,colframe=popink,boxrule=2pt,arc=9pt,shadow={2.5pt}{-2.5pt}{0pt}{popink},
      left=6pt,right=6pt,top=9pt,bottom=9pt,halign=center,valign=center,height=1.75cm,width=\linewidth,nobeforeafter]
    {\popdisplay\fontsize{12.5}{13}\selectfont\color{white}\MakeUppercase{#2}}\par\vspace{3pt}
    {\centering\popsemi\fontsize{7.8}{9.5}\selectfont\color{white}#3\par}
  \end{tcolorbox}}
\newcommand{\popcontenu}{%
  \par\noindent{\popxboldls\fontsize{7.5}{7.5}\selectfont\color{popmuted}CE COURS SIGNÉ CONTIENT}\par\vspace{7pt}
  \noindent\begin{minipage}[t]{0.235\linewidth}\poptile{popblue}{Cours}{complet, illustré, pas à pas}\end{minipage}\hfill
  \begin{minipage}[t]{0.235\linewidth}\poptile{poporange}{Méthodes}{et exercices résolus}\end{minipage}\hfill
  \begin{minipage}[t]{0.235\linewidth}\poptile{poppink}{Exercices}{type examen + corrigés}\end{minipage}\hfill
  \begin{minipage}[t]{0.235\linewidth}\poptile{popgreen}{Fiche bilan}{l'essentiel à retenir}\end{minipage}\par}

% Sommaire
\newtcolorbox{sommaire}{enhanced,colback=white,colframe=popink,boxrule=2.2pt,arc=10pt,
  drop shadow=popink,left=18pt,right=18pt,top=13pt,bottom=13pt,before skip=6pt,after skip=20pt,
  before upper={\poppill{Sommaire}{poppurple}{white}\par\vspace{9pt}\popsemi\fontsize{10.6}{14.5}\selectfont\color{popink}}}

% Signature de fin de cours — poussée tout en bas de la dernière page
\newcommand{\findecours}{%
  \par\vfill\nopagebreak
  \begin{center}\popsquiggle{poporange}\end{center}\vspace{4pt}
  \signatureonbuch
}
\AtBeginDocument{\def\llbracket{\mathopen{[\mkern-3.5mu[}}\def\rrbracket{\mathclose{]\mkern-3.5mu]}}} % intervalles d'entiers (glyphe ⟦ absent de la police)
% Glyphes absents de Fira Math : redéfinis à partir de glyphes disponibles
\AtBeginDocument{%
  \def\setminus{\mathbin{\backslash}}\let\smallsetminus\setminus
  \def\vdots{\mathord{\vbox{\baselineskip3.2pt\lineskiplimit0pt\kern4pt\hbox{.}\hbox{.}\hbox{.}}}}%
  \def\ddots{\mathinner{\mkern1mu\raise6.5pt\hbox{.}\mkern2mu\raise3.6pt\hbox{.}\mkern2mu\raise0.7pt\hbox{.}\mkern1mu}}%
  \def\triangle{\mathord{\scalebox{0.9}{$\symup{\Delta}$}}}%
  \def\nabla{\mathord{\raisebox{\depth}{\scalebox{1}[-1]{$\symup{\Delta}$}}}}%
  \def\checkmark{\popcheck{popdark}}%
  \def\star{\mathbin{\ast}}\def\bigstar{\mathord{\ast}}%
  \def\diamond{\mathbin{\scalebox{0.6}{$\Diamond$}}}%
  \def\bigcup{\mathop{\mathchoice{\raisebox{-0.3em}{\scalebox{1.7}{$\cup$}}}{\scalebox{1.25}{$\cup$}}{\cup}{\cup}}\displaylimits}%
  \def\bigcap{\mathop{\mathchoice{\raisebox{-0.3em}{\scalebox{1.7}{$\cap$}}}{\scalebox{1.25}{$\cap$}}{\cap}{\cap}}\displaylimits}%
  \def\lesseqqgtr{\mathrel{\lessgtr}}\def\bigoplus{\mathop{\oplus}\displaylimits}%
}
\providecommand{\ssi}{si et seulement si} % abréviation fréquente
\AtBeginDocument{\providecommand{\wideparen}{\overparen}} % arc de cercle

%% FIGFIX-BEGIN v5 — correctifs globaux des figures (docs/onbuch-generation/tools/figfix)
\usepackage{adjustbox}\usepackage{etoolbox}\tcbuselibrary{raster}
% toute figure TikZ/pgfplots plus large que la ligne est réduite à la largeur disponible
\BeforeBeginEnvironment{tikzpicture}{\begin{adjustbox}{max width=\linewidth}}
\AfterEndEnvironment{tikzpicture}{\end{adjustbox}}
% légendes pgfplots lisibles par-dessus les courbes
\pgfplotsset{every axis/.append style={legend style={fill=white,fill opacity=0.92,text opacity=1,draw=black!15,inner sep=3pt}}}
% étiquettes d'arêtes (syntaxe edge["..."]) avec fond blanc
\tikzset{every edge quotes/.append style={fill=white,inner sep=1.2pt}}
%% FIGFIX-END
\begin{document}

\pagegarde{El turismo y el desarrollo económico}{Espagnol}{1ère A}{Módulo 4 : Vie économique}{EA13}{2 h}{Cette leçon mixte du Módulo 4 (Vie économique) étudie le tourisme comme moteur du développement économique au Cameroun et dans le monde hispanophone. Elle combine la lecture d'un texte support, l'enrichissement du lexique thématique et la maîtrise de hay, ser et estar pour décrire et promouvoir un lieu touristique durable.
\vspace{4pt}
\begin{multicols}{2}\small
\begin{itemize}
\item À la fin de cette leçon, je sais lire et comprendre un texte en espagnol sur le tourisme et le développement économique
\item À la fin de cette leçon, je sais utiliser le lexique du tourisme (turista, hotel, paisaje, patrimonio, reserva natural, ecoturismo, guía)
\item À la fin de cette leçon, je sais distinguer et employer correctement hay, ser et estar pour décrire un lieu
\item À la fin de cette leçon, je sais accorder et placer les adjectifs qualificatifs dans une description
\item À la fin de cette leçon, je sais présenter oralement et par écrit un site touristique du Cameroun (Kribi, mont Cameroun, Waza)
\item À la fin de cette leçon, je sais citer des destinations touristiques de pays hispanophones
\end{itemize}\end{multicols}}

\begin{sommaire}
\begin{multicols}{2}
\begin{enumerate}
\item Texte support : El turismo, motor de desarrollo
\item Lexique thématique du tourisme
\item Grammaire : hay, ser, estar pour décrire un lieu
\item Les adjectifs dans la description
\item Méthode du commentaire et commentaire modèle
\item Production guidée : présenter et promouvoir un site
\item Activité d'intégration
\item Méthodes et exercices résolus
\item Exercices
\item Corrigés des exercices
\item Fiche bilan
\end{enumerate}
\end{multicols}
\end{sommaire}

\begin{objectifs}
\begin{itemize}
\item À la fin de cette leçon, je sais lire et comprendre un texte en espagnol sur le tourisme et le développement économique
\item À la fin de cette leçon, je sais utiliser le lexique du tourisme (turista, hotel, paisaje, patrimonio, reserva natural, ecoturismo, guía)
\item À la fin de cette leçon, je sais distinguer et employer correctement hay, ser et estar pour décrire un lieu
\item À la fin de cette leçon, je sais accorder et placer les adjectifs qualificatifs dans une description
\item À la fin de cette leçon, je sais présenter oralement et par écrit un site touristique du Cameroun (Kribi, mont Cameroun, Waza)
\item À la fin de cette leçon, je sais citer des destinations touristiques de pays hispanophones
\item À la fin de cette leçon, je sais promouvoir le tourisme durable avec des arguments en espagnol
\item À la fin de cette leçon, je sais traduire un court texte descriptif du français vers l'espagnol et inversement
\end{itemize}
\end{objectifs}

\begin{prerequis}
\begin{itemize}
\item Présent de l'indicatif des verbes réguliers et irréguliers courants (Seconde)
\item Le verbe haber impersonnel : hay (Seconde)
\item Genre et nombre des noms et des adjectifs (Seconde)
\item Vocabulaire de base des voyages et des vacances (Seconde)
\item Notions de géographie du Cameroun (sites de Kribi, Waza, mont Cameroun)
\end{itemize}
\end{prerequis}

\newpage

\coursec{Texte support : El turismo, motor de desarrollo}

Pendant les vacances, ta famille hésite : partir à Kribi pour ses plages de sable fin, gravir le mont Cameroun, le fameux « char des dieux », ou observer les éléphants au parc national de Waza ? Un ami espagnol te demande alors : \esp{¿Qué se puede visitar en Camerún?} « Que peut-on visiter au Cameroun ? » Pour lui répondre, il faut savoir décrire un lieu, vanter ses richesses et expliquer pourquoi le tourisme fait vivre tout un pays. Cette leçon t'apprend à lire un texte sur le tourisme et à présenter un site touristique en espagnol.

\textbf{Problématique :} comment le tourisme peut-il être un moteur du développement économique, au Cameroun comme dans les pays hispanophones ?

\courssub{Lecture globale du texte}

Avant de lire, observe le titre : \esp{El turismo, motor de desarrollo}. Les mots \esp{turismo} et \esp{motor} sont des \cle{mots transparents} : ils ressemblent au français « tourisme » et « moteur ». Tu peux donc déjà deviner le sujet du texte : le tourisme comme force qui fait avancer l'économie.

\begin{textebox}[El turismo, motor de desarrollo]
El turismo es un motor de desarrollo económico. En Camerún, miles de turistas visitan cada año las playas de Kribi, el monte Camerún y el parque nacional de Waza. Este sector crea empleos: guías, hoteles, restaurantes y transportes. Los turistas gastan dinero en alojamiento, comida y recuerdos, y este dinero entra en el país en forma de divisas.

En España, el turismo representa una parte importante de la riqueza nacional. Millones de visitantes llegan cada año atraídos por el sol, las playas, los museos y el patrimonio histórico de ciudades como Madrid, Barcelona o Sevilla. El turismo da trabajo a muchas familias españolas: camareros, recepcionistas, guías y conductores de autobús viven de esta actividad.

Sin embargo, el turismo también tiene problemas. Demasiados visitantes pueden dañar el paisaje y contaminar las playas. Por eso es necesario proteger el paisaje y el patrimonio: así nace el ecoturismo, un turismo que respeta la naturaleza y las reservas naturales. En Camerún, el ecoturismo permite visitar el parque de Waza sin molestar a los elefantes ni a los leones. En América Latina, países como Costa Rica son famosos por su ecoturismo en la selva.

En conclusión, el turismo es una gran oportunidad para el desarrollo, pero debe ser un turismo responsable y sostenible.
\end{textebox}

\textbf{Traduction du texte.} « Le tourisme est un moteur du développement économique. Au Cameroun, des milliers de touristes visitent chaque année les plages de Kribi, le mont Cameroun et le parc national de Waza. Ce secteur crée des emplois : guides, hôtels, restaurants et transports. Les touristes dépensent de l'argent en hébergement, nourriture et souvenirs, et cet argent entre dans le pays sous forme de devises. En Espagne, le tourisme représente une part importante de la richesse nationale. Des millions de visiteurs arrivent chaque année, attirés par le soleil, les plages, les musées et le patrimoine historique de villes comme Madrid, Barcelone ou Séville. Le tourisme donne du travail à de nombreuses familles espagnoles : serveurs, réceptionnistes, guides et chauffeurs de bus vivent de cette activité. Cependant, le tourisme a aussi des problèmes. Trop de visiteurs peuvent abîmer le paysage et polluer les plages. C'est pourquoi il faut protéger le paysage et le patrimoine : ainsi naît l'écotourisme, un tourisme qui respecte la nature et les réserves naturelles. Au Cameroun, l'écotourisme permet de visiter le parc de Waza sans déranger les éléphants ni les lions. En Amérique latine, des pays comme le Costa Rica sont célèbres pour leur écotourisme dans la forêt tropicale. En conclusion, le tourisme est une grande opportunité pour le développement, mais ce doit être un tourisme responsable et durable. »

\textbf{Glossaire (mots difficiles) :}

\begin{center}
\begin{tabularx}{\linewidth}{|X|X|X|}
\hline
\rowcolor{popblueL} Español & Français & Exemple du texte \\
\hline
\esp{motor} & moteur & \esp{un motor de desarrollo} \\
\hline
\esp{divisas} & devises (monnaies étrangères) & \esp{en forma de divisas} \\
\hline
\esp{empleo} & emploi & \esp{crea empleos} \\
\hline
\esp{cifra} & chiffre & \esp{la cifra de turistas} \\
\hline
\esp{atraer} & attirer & \esp{atraídos por el sol} \\
\hline
\esp{alojamiento} & hébergement & \esp{dinero en alojamiento} \\
\hline
\esp{recuerdos} & souvenirs & \esp{comida y recuerdos} \\
\hline
\esp{selva} & forêt tropicale & \esp{la selva de Costa Rica} \\
\hline
\esp{dañar} & abîmer, endommager & \esp{dañar el paisaje} \\
\hline
\esp{sostenible} & durable & \esp{turismo sostenible} \\
\hline
\end{tabularx}
\end{center}

\begin{methode}[Comment lire un texte en espagnol]
\begin{enumerate}
\item \textbf{Lire le titre et anticiper} : devine le sujet avant de lire. Ici, \esp{motor de desarrollo} annonce un texte économique.
\item \textbf{Lecture globale sans dictionnaire} : lis tout le texte une première fois pour saisir l'idée générale, même si tu ne comprends pas chaque mot.
\item \textbf{Repérer les mots transparents} : \esp{turismo}, \esp{hotel}, \esp{reserva}, \esp{patrimonio}, \esp{restaurante} se devinent facilement. Ils forment le squelette du sens.
\item \textbf{Lecture détaillée avec le glossaire} : relis phrase par phrase en t'aidant du glossaire pour les mots nouveaux (\esp{divisas}, \esp{atraer}, \esp{alojamiento}).
\item \textbf{Vérifier} : résume le texte en une phrase en français pour contrôler ta compréhension.
\end{enumerate}
\end{methode}

\begin{attention}[Pièges de la lecture]
Ne traduis jamais mot à mot : \esp{Los turistas gastan dinero} se traduit « Les touristes dépensent de l'argent », pas « gaspillent » (\esp{gastar} = dépenser, c'est un \cle{faux ami} partiel). Attention aussi : \esp{empleo} signifie « emploi », et non « emploiement » ; \esp{recuerdos} signifie « souvenirs » (objets), pas « se rappeler » ; \esp{divisas} désigne les devises étrangères, pas les « devises » au sens de devis de travaux. Les mots transparents aident, mais vérifie toujours le contexte.
\end{attention}

\begin{savaistu}[Le tourisme en Espagne]
L'Espagne est le deuxième pays le plus visité du monde, juste après la France : elle accueille chaque année plus de 80 millions de touristes, soit plus que sa propre population ! Le tourisme y représente environ 12 \% du produit intérieur brut (PIB) et des millions d'emplois. C'est pourquoi on dit souvent que le tourisme est « la industria sin chimeneas », « l'industrie sans cheminées » : elle produit de la richesse sans usines ni fumée.
\end{savaistu}

\courssub{Repérage des informations : qui ? quoi ? où ?}

La \cle{lecture globale} consiste à répondre à trois questions simples pour identifier le contenu du texte.

\begin{center}
\begin{tabularx}{\linewidth}{|X|X|X|}
\hline
\rowcolor{popblueL} Question & Réponse dans le texte & Citation \\
\hline
¿Quién? (qui ?) & Les touristes, les guides, les familles espagnoles & \esp{miles de turistas visitan} \\
\hline
¿Qué? (quoi ?) & Le tourisme crée des emplois et rapporte des devises & \esp{Este sector crea empleos} \\
\hline
¿Dónde? (où ?) & Au Cameroun (Kribi, Waza), en Espagne, au Costa Rica & \esp{las playas de Kribi} \\
\hline
¿Cuál es el problema? & Trop de touristes abîment le paysage & \esp{pueden dañar el paisaje} \\
\hline
¿Cuál es la solución? & L'écotourisme, un tourisme responsable & \esp{así nace el ecoturismo} \\
\hline
\end{tabularx}
\end{center}

\begin{popfigure}
\begin{center}\tcbox[colback=popblue,colframe=popblue,coltext=white,arc=9pt,boxsep=4pt,left=6pt,right=6pt]{\bfseries\small TURISMO}\end{center}
\vspace{-2pt}
\begin{tcbraster}[raster columns=2,raster equal height=rows,raster column skip=6pt,raster row skip=6pt,enhanced,blankest]
\begin{tcolorbox}[enhanced,colback=popblue,colframe=popblue,coltext=white,boxrule=0.9pt,arc=3pt,left=4pt,right=4pt,top=3pt,bottom=3pt,boxsep=1pt,halign=left]\bfseries\scriptsize empleo guías, hoteles, restaurantes\end{tcolorbox}
\begin{tcolorbox}[enhanced,colback=poporange,colframe=poporange,coltext=white,boxrule=0.9pt,arc=3pt,left=4pt,right=4pt,top=3pt,bottom=3pt,boxsep=1pt,halign=left]\bfseries\scriptsize divisas dinero que entra en el país\end{tcolorbox}
\begin{tcolorbox}[enhanced,colback=popgreen,colframe=popgreen,coltext=white,boxrule=0.9pt,arc=3pt,left=4pt,right=4pt,top=3pt,bottom=3pt,boxsep=1pt,halign=left]\bfseries\scriptsize cultura patrimonio, museos\end{tcolorbox}
\begin{tcolorbox}[enhanced,colback=poppurple,colframe=poppurple,coltext=white,boxrule=0.9pt,arc=3pt,left=4pt,right=4pt,top=3pt,bottom=3pt,boxsep=1pt,halign=left]\bfseries\scriptsize naturaleza playas, parques, reservas\end{tcolorbox}
\end{tcbraster}

\legende{Carte mentale : les quatre visages du tourisme dans le texte.}
\end{popfigure}

\courssub{Lecture sélective : chiffres, lieux, avantages}

La \cle{lecture sélective} consiste à relire le texte en ne cherchant qu'un type d'information précis, comme un détective. Entraîne-toi à relever :

\begin{itemize}
\item \textbf{Les quantités} : \esp{miles de turistas} (des milliers de touristes), \esp{millones de visitantes} (des millions de visiteurs).
\item \textbf{Les lieux cités} : Kribi, le mont Cameroun, le parc national de Waza (Cameroun) ; Madrid, Barcelona, Sevilla (Espagne) ; Costa Rica (Amérique latine).
\item \textbf{Les avantages économiques} : création d'emplois (\esp{crea empleos}), entrée de devises (\esp{divisas}), travail pour les familles (\esp{da trabajo a muchas familias}).
\item \textbf{Les risques} : dégradation du paysage (\esp{dañar el paisaje}), pollution des plages (\esp{contaminar las playas}).
\end{itemize}

\begin{popfigure}
\begin{tikzpicture}[node distance=0.6cm,
caja/.style={rectangle, rounded corners, draw=popblue, fill=popblueL, align=center, font=\small, minimum width=3.4cm, minimum height=1.1cm},
flecha/.style={-{Stealth[length=3mm]}, thick, poporange}]
\node[caja, align=center] (t){turistas\\ visitan el país};
\node[caja, right=of t, align=center] (h){hoteles y guías\\ trabajan};
\node[caja, right=of h, align=center] (i){ingresos\\ empleo y divisas};
\node[caja, right=of i, align=center] (d){desarrollo\\ económico};
\draw[flecha] (t) -- (h);
\draw[flecha] (h) -- (i);
\draw[flecha] (i) -- (d);
\end{tikzpicture}
\legende{La chaîne économique du tourisme : des visiteurs au développement.}
\end{popfigure}

Ce schéma résume l'argument principal du texte : les touristes arrivent, ils font travailler les hôtels et les guides, l'argent entre dans le pays, et le pays se développe. C'est exactement le sens du mot \esp{motor} dans le titre.

\courssub{Questions de compréhension}

\textbf{A. Vrai ou faux ? Justifie ta réponse par une citation du texte.}

\begin{enumerate}
\item \esp{El turismo no crea empleos en Camerún.} (V / F)
\item \esp{En España, el turismo es una parte importante de la riqueza nacional.} (V / F)
\item \esp{El ecoturismo respeta la naturaleza.} (V / F)
\item \esp{Costa Rica es famoso por sus playas contaminadas.} (V / F)
\end{enumerate}

\textbf{B. Questions ouvertes.}

\begin{enumerate}
\item \esp{¿Qué lugares de Camerún visitan los turistas?}
\item \esp{¿Por qué es necesario proteger el paisaje?}
\item \esp{¿Qué es el ecoturismo?}
\end{enumerate}

\begin{exoresolu}[Vrai ou faux : justifier par le texte]
Énoncé : réponds à l'affirmation \esp{El turismo no crea empleos en Camerún} (vrai ou faux) et justifie par une citation.
\tcblower
\textbf{Solution.} C'est \textbf{faux} (\esp{falso}). Le texte dit exactement le contraire : \esp{Este sector crea empleos: guías, hoteles, restaurantes y transportes.} « Ce secteur crée des emplois : guides, hôtels, restaurants et transports. » La phrase complète de réponse serait : \esp{Es falso: el texto dice que este sector crea empleos.} Remarque la méthode : on donne d'abord le jugement (falso), puis la citation exacte entre guillemets, sans ajouter d'opinion personnelle.
\end{exoresolu}

\begin{corrige}[Questions ouvertes]
\begin{enumerate}
\item \esp{Los turistas visitan las playas de Kribi, el monte Camerún y el parque nacional de Waza.} « Les touristes visitent les plages de Kribi, le mont Cameroun et le parc national de Waza. »
\item \esp{Es necesario proteger el paisaje porque demasiados visitantes pueden dañarlo y contaminar las playas.} « Il faut protéger le paysage parce que trop de visiteurs peuvent l'abîmer et polluer les plages. »
\item \esp{El ecoturismo es un turismo que respeta la naturaleza y las reservas naturales.} « L'écotourisme est un tourisme qui respecte la nature et les réserves naturelles. »
\end{enumerate}
\end{corrige}

\courssub{Reformulation : résumer le texte en trois phrases}

Résumer, c'est dire l'essentiel \emph{avec tes propres mots}, sans recopier les phrases du texte. En espagnol, on utilise des connecteurs simples : \esp{primero} (d'abord), \esp{además} (de plus), \esp{sin embargo} (cependant), \esp{en conclusión} (en conclusion).

\begin{exemplebox}[Modèle de résumé en trois phrases]
\esp{El turismo es muy importante para la economía de Camerún y de España porque crea empleos y trae divisas.} « Le tourisme est très important pour l'économie du Cameroun et de l'Espagne parce qu'il crée des emplois et rapporte des devises. »

\esp{Sin embargo, demasiados turistas pueden dañar el paisaje y el patrimonio.} « Cependant, trop de touristes peuvent abîmer le paysage et le patrimoine. »

\esp{En conclusión, el ecoturismo es la solución: un turismo responsable que respeta la naturaleza.} « En conclusion, l'écotourisme est la solution : un tourisme responsable qui respecte la nature. »
\end{exemplebox}

\begin{aretenir}[L'essentiel de la section]
\begin{itemize}
\item Le texte présente le tourisme comme un \cle{motor de desarrollo} : il crée des emplois (\esp{empleo}) et rapporte des devises (\esp{divisas}).
\item Stratégie de lecture : titre $\rightarrow$ lecture globale $\rightarrow$ mots transparents $\rightarrow$ lecture détaillée avec glossaire.
\item Lecture sélective : repérer les chiffres, les lieux, les avantages et les risques.
\item Pour justifier une réponse : jugement + citation exacte du texte.
\item Pour résumer : trois phrases, ses propres mots, des connecteurs (\esp{sin embargo}, \esp{en conclusión}).
\end{itemize}
\end{aretenir}

\begin{exercice}[À toi de jouer]
\begin{enumerate}
\item Relis le texte et souligne en rouge les lieux, en bleu les métiers, en vert les problèmes écologiques.
\item Réponds par écrit : \esp{¿Qué atrae a los turistas en España?} Justifie par une citation.
\item Résume le texte en trois phrases avec tes propres mots, en utilisant \esp{además} et \esp{por eso}.
\end{enumerate}
\end{exercice}

\coursec{Lexique thématique du tourisme}

Dans la section précédente, nous avons lu et commenté le texte \esp{El turismo, motor de desarrollo}. Ce texte contenait de nombreux mots nouveaux : \esp{turista}, \esp{paisaje}, \esp{patrimonio}, \esp{reserva natural}, \esp{ecoturismo}... Pour bien comprendre un document sur le tourisme et, surtout, pour pouvoir présenter toi-même un site touristique du Cameroun ou d'un pays hispanophone, il faut d'abord maîtriser ce vocabulaire. C'est l'objectif de cette section : organiser, comprendre et employer le lexique du tourisme.

\courssub{Observation : un texte pour découvrir le vocabulaire}

Lis d'abord ce texte original. Souligne mentalement tous les mots qui appartiennent au champ lexical du tourisme.

\begin{textebox}[El Camerún, un destino por descubrir]
El Camerún es un país de África Central con un enorme potencial turístico. Cada año, muchos turistas visitan sus playas, sus montañas y sus parques nacionales.

En el sur, las playas de Kribi atraen a los viajeros que buscan sol y mar. Cerca de la ciudad, las cataratas de la Lobé caen directamente en el océano Atlántico: es un paisaje único en el mundo. Los hoteles de la costa ofrecen un alojamiento cómodo, y los guías locales muestran el patrimonio natural de la región a los visitantes.

En el oeste, el monte Camerún, un volcán de más de cuatro mil metros, es ideal para las excursiones y el senderismo. Los turistas deportivos recorren sus senderos durante varios días.

En el norte, el parque nacional de Waza es una gran reserva natural. Allí se pueden observar elefantes, leones, jirafas y muchas aves. Vale la pena visitar este parque con un guía experimentado.

Hoy, el país quiere promover el turismo sostenible. El ecoturismo permite proteger el medio ambiente y respetar la cultura de las poblaciones locales. Conservar la naturaleza es esencial para el futuro: si destruimos los paisajes y los monumentos, los turistas dejarán de venir. Hacer turismo de manera responsable es una obligación para todos.
\end{textebox}

\textbf{Glossaire.}

\begin{center}
\begin{tabularx}{\linewidth}{|X|X|}
\hline
\rowcolor{popblueL} \textbf{Español} & \textbf{Français} \\
\hline
un destino por descubrir & une destination à découvrir \\
\hline
el potencial turístico & le potentiel touristique \\
\hline
buscan sol y mar & cherchent le soleil et la mer \\
\hline
caen directamente & tombent (se jettent) directement \\
\hline
el senderismo & la randonnée \\
\hline
los senderos & les sentiers \\
\hline
las aves & les oiseaux \\
\hline
experimentado & expérimenté \\
\hline
dejarán de venir & cesseront de venir \\
\hline
una obligación & une obligation \\
\hline
\end{tabularx}
\end{center}

\textbf{Questions de compréhension.}

\begin{enumerate}
\item ¿Qué atrae a los viajeros en el sur del Camerún? (Qu'est-ce qui attire les voyageurs dans le sud du Cameroun ?)
\item ¿Qué se puede observar en el parque nacional de Waza? (Que peut-on observer dans le parc national de Waza ?)
\item ¿Qué tipo de turismo quiere promover el país? (Quel type de tourisme le pays veut-il promouvoir ?)
\end{enumerate}

\courssub{Champ 1 : les acteurs du tourisme}

Qui participe au tourisme ? Observe ces phrases tirées du texte :

\esp{Los guías locales muestran el patrimonio a los visitantes.} « Les guides locaux montrent le patrimoine aux visiteurs. »

\esp{Las playas de Kribi atraen a los viajeros.} « Les plages de Kribi attirent les voyageurs. »

\begin{definition}[Los actores del turismo]
\begin{itemize}
\item \cle{el turista / la turista} : le/la touriste, personne qui voyage pour le plaisir.
\item \cle{el viajero / la viajera} : le/la voyageur(-euse), personne qui voyage (sens plus large que \esp{turista}).
\item \cle{el guía / la guía} : le/la guide, personne qui accompagne et informe les touristes.
\item \cle{el recepcionista / la recepcionista} : le/la réceptionniste, personne qui accueille les clients à l'hôtel.
\item \cle{el visitante} : le visiteur.
\end{itemize}
\end{definition}

\begin{attention}[Le genre de « guía » et noms épicènes]
Les mots \esp{guía}, \esp{turista}, \esp{recepcionista}, \esp{estudiante}, \esp{artista} sont \emph{épicènes} : leur forme ne change pas, seul l'article indique le genre (\esp{el guía} / \esp{la guía}, \esp{el turista} / \esp{la turista}). Ne dis jamais \esp{la guío} ni \esp{la turisto}. C'est le même phénomène qu'en français avec « un/une journaliste » ou « un/une pianiste ».
\end{attention}

\courssub{Champ 2 : les lieux et l'hébergement}

\begin{definition}[Los lugares y el alojamiento]
\begin{itemize}
\item \cle{el hotel} : l'hôtel. \esp{Los hoteles de la costa son cómodos.} « Les hôtels de la côte sont confortables. »
\item \cle{el alojamiento} : l'hébergement, le logement (terme général).
\item \cle{el camping} : le camping (terrain). Pour l'activité, on dit \esp{ir de camping} ou \cle{acampar} (camper).
\item \cle{el paisaje} : le paysage. \textbf{Attention : masculin} (\esp{el paisaje}, \esp{un paisaje único}).
\item \cle{la playa} : la plage.
\item \cle{la reserva natural} : la réserve naturelle.
\item \cle{el parque nacional} : le parc national.
\end{itemize}
\end{definition}

\begin{attention}[Faux amis et pièges de genre]
\begin{itemize}
\item \esp{el paisage} est masculin : on dit \esp{un paisaje único}, jamais \esp{la paisaje}.
\item \esp{el camping} (masc.) désigne le terrain ; le verbe est \esp{acampar}. Ne confonds pas avec \esp{el campamento} (campement, colonie de vacances).
\item \esp{el viaje} = le voyage (nom d'action) ; \esp{el viajero} = le voyageur (personne). Erreur fréquente : dire \esp{*muchos viajes} pour « beaucoup de voyageurs ».
\end{itemize}
\end{attention}

\courssub{Champ 3 : le patrimoine et les activités}

\begin{definition}[El patrimonio y las actividades]
\begin{itemize}
\item \cle{el patrimonio} : le patrimoine (naturel, culturel, historique).
\item \cle{el monumento} : le monument.
\item \cle{el museo} : le musée.
\item \cle{la excursión} : l'excursion. \esp{Hacer una excursión} = faire une excursion.
\item \cle{visitar} : visiter (un lieu, un monument). \esp{Visitar un museo.}
\item \cle{recorrer} : parcourir, sillonner (un pays, des sentiers). \esp{Recorrer los senderos.}
\item \cle{conocer} : connaître, visiter pour la première fois. \esp{Conocer el monte Camerún.}
\end{itemize}
\end{definition}

\begin{exemplebox}[Exemples traduits]
\esp{El guía muestra el patrimonio a los turistas.} « Le guide montre le patrimoine aux touristes. »

\esp{La reserva natural de Waza atrae a muchos visitantes.} « La réserve naturelle de Waza attire beaucoup de visiteurs. »

\esp{Los turistas recorren el país en autobús.} « Les touristes parcourent le pays en car. »

\esp{Mañana vamos a hacer una excursión al monte Camerún.} « Demain nous allons faire une excursion au mont Cameroun. »

\esp{Quiero conocer las cataratas de la Lobé.} « Je veux visiter/découvrir les chutes de la Lobé. »
\end{exemplebox}

\courssub{Champ 4 : le tourisme durable}

\begin{definition}[El ecoturismo]
\cle{El ecoturismo} (l'écotourisme) : \esp{forma de turismo que respeta la naturaleza y la cultura local} — forme de tourisme qui respecte la nature et la culture locale. On parle aussi de \cle{turismo sostenible} (tourisme durable) ou \esp{turismo sustentable} (Amérique latine).
\end{definition}

\begin{definition}[Los verbos de la protección]
\begin{itemize}
\item \cle{proteger} : protéger. \esp{Proteger el medio ambiente.}
\item \cle{respetar} : respecter. \esp{Respetar la cultura local.}
\item \cle{conservar} : conserver, préserver. \esp{Conservar la naturaleza.}
\item \cle{el medio ambiente} : l'environnement (deux mots).
\end{itemize}
\end{definition}

\begin{attention}[Attention aux calques du français]
\begin{itemize}
\item « L'environnement » se dit \esp{el medio ambiente} (deux mots), jamais \esp{el ambiente} seul dans ce sens (\esp{ambiente} seul = ambiance, atmosphère).
\item \esp{conservar} = préserver, maintenir en bon état (nature, monuments). Pour « conserver un document » (archiver), on utilise \esp{guardar} ou \esp{archivar}.
\item \esp{sostenible} / \esp{sustentable} = durable. N'utilise pas \esp{duradero} (qui signifie « qui dure longtemps », ex. \esp{un material duradero}).
\end{itemize}
\end{attention}

\courssub{Expressions utiles pour parler du tourisme}

\begin{aretenir}[Expresiones útiles]
\begin{itemize}
\item \esp{hacer turismo} : faire du tourisme. \esp{Muchos europeos hacen turismo en África.}
\item \esp{atraer a los turistas} : attirer les touristes. \esp{Las playas atraen a los turistas.}
\item \esp{promover el turismo} / \esp{fomentar el turismo} : promouvoir le tourisme. \esp{El gobierno quiere promover el turismo sostenible.}
\item \esp{vale la pena visitar} : cela vaut la peine de visiter (impersonnel, 3\textsuperscript{e} pers. sg). \esp{Vale la pena visitar Waza.}
\item \esp{hay que proteger} : il faut protéger (impersonnel, infinitif). \esp{Hay que proteger el medio ambiente.}
\end{itemize}
\end{aretenir}

\begin{attention}[« Atraer » et la préposition « a » personnelle]
Le verbe \esp{atraer} (attirer) se construit avec la préposition \esp{a} devant la personne (COI animé) : \esp{atraer a los turistas}. C'est la \emph{a} personnelle espagnole, obligatoire devant un complément de personne, alors que le français n'a pas de préposition : « attirer les touristes ». Ne l'oublie pas. Conjugaison présent : \esp{atraigo, atraes, atrae, atraemos, atraéis, atraen} (diphthongaison \emph{ie} / \emph{i}).
\end{attention}

\courssub{Tableau récapitulatif : les quatre champs lexicaux}

\begin{center}
\begin{tabularx}{\linewidth}{|l|X|X|}
\hline
\rowcolor{popblueL} \textbf{Champ} & \textbf{Español} & \textbf{Français} \\
\hline
Acteurs & el/la turista ; el/la guía ; el/la viajero/a ; el/la recepcionista ; el visitante & le/la touriste ; le/la guide ; le/la voyageur(-euse) ; le/la réceptionniste ; le visiteur \\
\hline
Lieux et hébergement & el hotel ; el alojamiento ; el camping ; el paisaje ; la playa ; la reserva natural ; el parque nacional & l'hôtel ; l'hébergement ; le camping ; le paysage ; la plage ; la réserve naturelle ; le parc national \\
\hline
Patrimoine et activités & el patrimonio ; el monumento ; el museo ; la excursión ; visitar ; recorrer ; conocer & le patrimoine ; le monument ; le musée ; l'excursion ; visiter ; parcourir ; découvrir \\
\hline
Tourisme durable & el ecoturismo ; el turismo sostenible ; proteger ; respetar ; conservar ; el medio ambiente & l'écotourisme ; le tourisme durable ; protéger ; respecter ; préserver ; l'environnement \\
\hline
\end{tabularx}
\end{center}

\begin{popfigure}
\begin{center}\tcbox[colback=popblue,colframe=popblue,coltext=white,arc=9pt,boxsep=4pt,left=6pt,right=6pt]{\bfseries\small El turismo}\end{center}
\vspace{-2pt}
\begin{tcbraster}[raster columns=2,raster equal height=rows,raster column skip=6pt,raster row skip=6pt,enhanced,blankest]
\begin{tcolorbox}[enhanced,colback=white,colframe=poporange,boxrule=0.9pt,arc=3pt,title={Actores},fonttitle=\bfseries\scriptsize,coltitle=white,colbacktitle=poporange,left=4pt,right=4pt,top=2pt,bottom=2pt,boxsep=1pt,toptitle=1.5pt,bottomtitle=1.5pt,halign=left]
\begin{itemize}[nosep,leftmargin=*,itemsep=1pt,font=\scriptsize]
\item turista
\item guía
\end{itemize}
\end{tcolorbox}
\begin{tcolorbox}[enhanced,colback=white,colframe=popgreen,boxrule=0.9pt,arc=3pt,title={Lugares},fonttitle=\bfseries\scriptsize,coltitle=white,colbacktitle=popgreen,left=4pt,right=4pt,top=2pt,bottom=2pt,boxsep=1pt,toptitle=1.5pt,bottomtitle=1.5pt,halign=left]
\begin{itemize}[nosep,leftmargin=*,itemsep=1pt,font=\scriptsize]
\item hotel
\item playa
\end{itemize}
\end{tcolorbox}
\begin{tcolorbox}[enhanced,colback=white,colframe=poppurple,boxrule=0.9pt,arc=3pt,title={Patrimonio},fonttitle=\bfseries\scriptsize,coltitle=white,colbacktitle=poppurple,left=4pt,right=4pt,top=2pt,bottom=2pt,boxsep=1pt,toptitle=1.5pt,bottomtitle=1.5pt,halign=left]
\begin{itemize}[nosep,leftmargin=*,itemsep=1pt,font=\scriptsize]
\item museo
\item excursión
\end{itemize}
\end{tcolorbox}
\begin{tcolorbox}[enhanced,colback=white,colframe=poppink,boxrule=0.9pt,arc=3pt,title={Sostenible},fonttitle=\bfseries\scriptsize,coltitle=white,colbacktitle=poppink,left=4pt,right=4pt,top=2pt,bottom=2pt,boxsep=1pt,toptitle=1.5pt,bottomtitle=1.5pt,halign=left]
\begin{itemize}[nosep,leftmargin=*,itemsep=1pt,font=\scriptsize]
\item ecoturismo
\item proteger
\end{itemize}
\end{tcolorbox}
\end{tcbraster}

\legende{Carte mentale des quatre champs lexicaux du tourisme.}
\end{popfigure}

\courssub{Les sites touristiques du Cameroun en espagnol}

Pour présenter ton pays aux hispanophones, apprends par cœur ces noms propres (avec accents) :

\begin{aretenir}[Sitios turísticos de Camerún]
\begin{itemize}
\item \esp{las playas de Kribi} : les plages de Kribi.
\item \esp{el monte Camerún} : le mont Cameroun (accents sur \emph{ó} et \emph{ú}).
\item \esp{el parque nacional de Waza} : le parc national de Waza.
\item \esp{las cataratas de la Lobé} : les chutes de la Lobé (accent sur \emph{é}).
\end{itemize}
Remarque : le pays s'écrit \esp{Camerún} (accent sur le \emph{u}) et l'habitant est \esp{camerunés / camerunesa} (accent sur le \emph{é}).
\end{aretenir}

\begin{popfigure}
\begin{tikzpicture}[scale=1.0]
\draw[fill=popcream, draw=popdark, thick, rounded corners=6pt]
  (0,0) -- (0.4,1.6) -- (0.2,3.0) -- (1.0,4.2) -- (2.4,4.8) -- (3.8,4.4) -- (4.2,3.2) -- (3.6,2.0) -- (3.4,0.8) -- (2.2,-0.4) -- (1.0,-0.2) -- cycle;
\node[font=\bfseries\small, text=popdark] at (2.1,2.2) {CAMERÚN};
\fill[popblue] (0.7,0.5) circle (2.5pt);
\node[right, font=\small, text=popblue] at (0.85,0.5) {Kribi (playas, cataratas de la Lobé)};
\fill[poporange] (1.1,1.2) circle (2.5pt);
\node[right, font=\small, text=poporange] at (1.25,1.2) {Monte Camerún};
\fill[popgreen] (2.6,4.1) circle (2.5pt);
\node[left, font=\small, text=popgreen!60!black] at (2.45,4.55) {Parque nacional de Waza};
\draw[->, thick, popmuted] (-0.6,-0.6) -- (-0.6,0.2) node[above, font=\small]{N};
\node[font=\footnotesize\itshape, text=popmuted] at (0.4,-0.9) {Océano Atlántico};
\end{tikzpicture}
\legende{Carte simplifiée du Cameroun : Kribi, le mont Cameroun et le parc de Waza.}
\end{popfigure}

\begin{savaistu}[Las cataratas de la Lobé]
Les chutes de la Lobé, près de Kribi, sont parmi les rares cascades au monde qui se jettent directement dans l'océan. En espagnol : \esp{Las cataratas de la Lobé caen directamente en el océano Atlántico.} C'est un argument formidable pour promouvoir le tourisme camerounais auprès des visiteurs hispanophones d'Espagne, d'Amérique latine ou de Guinée équatoriale.
\end{savaistu}

\begin{exoresolu}[Complète avec le mot du lexique qui convient]
Énoncé. Complète chaque phrase avec le mot correct : \esp{guía}, \esp{paisaje}, \esp{reserva natural}, \esp{ecoturismo}, \esp{alojamiento}.

1. El \dots{} de Waza es magnífico: montañas, animales y sabana.

2. El parque nacional de Waza es una gran \dots{}.

3. El \dots{} muestra los monumentos a los turistas.

4. El \dots{} respeta la naturaleza y la cultura local.

5. Los hoteles de Kribi ofrecen un \dots{} cómodo.
\tcblower
Solution détaillée.

1. \esp{El paisaje de Waza es magnífico.} « Le paysage de Waza est magnifique. » (\esp{paisaje} est masculin : \esp{el paisaje}.)

2. \esp{El parque nacional de Waza es una gran reserva natural.} « Le parc national de Waza est une grande réserve naturelle. »

3. \esp{El guía muestra los monumentos a los turistas.} « Le guide montre les monuments aux touristes. »

4. \esp{El ecoturismo respeta la naturaleza y la cultura local.} « L'écotourisme respecte la nature et la culture locale. »

5. \esp{Los hoteles de Kribi ofrecen un alojamiento cómodo.} « Les hôtels de Kribi offrent un hébergement confortable. »
\end{exoresolu}

\begin{exercice}[À ton tour]
1. Traduis en espagnol :
   a) La plage de Kribi attire beaucoup de touristes.
   b) Il faut protéger l'environnement.
   c) Cela vaut la peine de visiter le mont Cameroun.

2. Classe ces mots dans les quatre champs lexicaux : \esp{museo}, \esp{turista}, \esp{proteger}, \esp{camping}, \esp{excursión}, \esp{recepcionista}, \esp{medio ambiente}, \esp{playa}.
\end{exercice}

\begin{corrige}[Exercice « À ton tour »]
1. a) \esp{La playa de Kribi atrae a muchos turistas.} (n'oublie pas le \esp{a} après \esp{atraer}) \\
   b) \esp{Hay que proteger el medio ambiente.} \\
   c) \esp{Vale la pena visitar el monte Camerún.}

2. Acteurs : \esp{turista}, \esp{recepcionista}. Lieux et hébergement : \esp{camping}, \esp{playa}. Patrimoine et activités : \esp{museo}, \esp{excursión}. Tourisme durable : \esp{proteger}, \esp{medio ambiente}.
\end{corrige}

\begin{aretenir}[L'essentiel de la section]
Le lexique du tourisme s'organise en quatre champs : les acteurs (\esp{turista}, \esp{guía}), les lieux et l'hébergement (\esp{hotel}, \esp{playa}, \esp{reserva natural}), le patrimoine et les activités (\esp{patrimonio}, \esp{excursión}, \esp{recorrer}, \esp{conocer}) et le tourisme durable (\esp{ecoturismo}, \esp{proteger}, \esp{medio ambiente}). Retiens les pièges : \esp{el/la guía} selon la personne, \esp{el paisage} masculin, \esp{atraer a} avec préposition personnelle, \esp{Camerún} avec accent, \esp{medio ambiente} en deux mots. Ce vocabulaire te servira dans les sections suivantes pour décrire un lieu avec \esp{hay}, \esp{ser} et \esp{estar}, puis pour présenter et promouvoir un site touristique.
\end{aretenir}

\coursec{Grammaire : hay, ser, estar pour décrire un lieu}

Dans la section précédente, nous avons découvert le \cle{lexique thématique du tourisme} : \esp{turismo}, \esp{turista}, \esp{hotel}, \esp{paisaje}, \esp{patrimonio}, \esp{reserva natural}, \esp{ecoturismo}, \esp{guía}. Ces mots nous permettent de nommer les choses. Mais pour \emph{décrire} un lieu touristique — dire ce qu'il y a, ce qu'il est, où il se trouve — il nous faut maintenant trois outils grammaticaux essentiels : \esp{hay}, \esp{ser} et \esp{estar}. C'est l'un des points les plus importants de l'année, car le français, lui, utilise un seul verbe « être » (et la tournure « il y a ») là où l'espagnol en distingue trois.

\courssub{Observation : trois phrases, trois verbes}

Reprenons trois phrases tirées de notre texte support et observons-les attentivement.

\begin{exemplebox}[Phrases à observer]
\esp{Hay playas magníficas en Kribi.} « Il y a des plages magnifiques à Kribi. »\par
\esp{Kribi es una ciudad turística.} « Kribi est une ville touristique. »\par
\esp{Waza está en el norte de Camerún.} « Waza est dans le nord du Cameroun. »
\end{exemplebox}

Posons-nous trois questions :

\begin{enumerate}
\item Dans la première phrase, que dit-on ? On affirme l'\cle{existence}, la présence de quelque chose (des plages). L'espagnol utilise \esp{hay}.
\item Dans la deuxième phrase, que dit-on ? On \cle{définit} Kribi, on donne sa nature, sa caractéristique permanente (c'est une ville touristique). L'espagnol utilise \esp{es} (verbe \esp{ser}).
\item Dans la troisième phrase, que dit-on ? On indique la \cle{localisation} de Waza (où il se trouve). L'espagnol utilise \esp{está} (verbe \esp{estar}).
\end{enumerate}

En français, les deux dernières phrases utilisent le même verbe « est ». En espagnol, ce serait une faute grave : il faut choisir entre \esp{ser} et \esp{estar}. Voyons chaque règle en détail.

\courssub{Règle 1 : hay, le verbe de l'existence}

\begin{propriete}[Hay : haber impersonnel]
\esp{Hay} est la forme impersonnelle du verbe \esp{haber} au présent. Elle exprime l'\cle{existence} ou la \cle{présence} de quelque chose, et correspond au français « il y a ».\par
\textbf{Forme :} \esp{hay} + nom (avec ou sans article, avec un nombre, avec \esp{mucho}, \esp{poco}, etc.)\par
\textbf{Particularité essentielle :} \esp{hay} est \cle{invariable}. Il reste identique même devant un nom pluriel.\par
\esp{Hay un hotel en la playa.} « Il y a un hôtel sur la plage. »\par
\esp{Hay diez hoteles en la playa.} « Il y a dix hôtels sur la plage. »\par
\esp{Hay muchos turistas en Kribi.} « Il y a beaucoup de touristes à Kribi. »
\end{propriete}

Remarque la logique : \esp{hay} ne désigne pas une personne, il n'a donc ni sujet ni conjugaison complète au présent. On ne dit jamais \esp{han} au présent pour « il y a ».

\begin{exemplebox}[Hay en contexte touristique]
\esp{En Waza hay elefantes y leones.} « À Waza il y a des éléphants et des lions. »\par
\esp{En el monte Camerún hay una reserva natural.} « Sur le mont Cameroun il y a une réserve naturelle. »\par
\esp{¿Hay un guía en el parque?} « Est-ce qu'il y a un guide dans le parc ? »\par
\esp{No hay muchos hoteles en esta región.} « Il n'y a pas beaucoup d'hôtels dans cette région. »
\end{exemplebox}

Pour la négation, on place simplement \esp{no} devant : \esp{no hay} « il n'y a pas ». Pour la question, l'intonation et les signes \esp{¿...?} suffisent : \esp{¿Hay...?}

\begin{attention}[Piège n° 1 : « il y a » n'est jamais « tiene » ni « es »]
Les francophones traduisent souvent mot à mot. Retiens bien :
\begin{itemize}
\item « il y a » = \esp{hay}, jamais \esp{tiene} (qui signifie « il/elle a », possession, verbe \esp{tener}) ;
\item « il y a » n'est jamais \esp{es} non plus.
\end{itemize}
Faux : \esp{Tiene muchos hoteles en Kribi.} (cela voudrait dire « Il a beaucoup d'hôtels à Kribi », comme si quelqu'un les possédait !)\par
Juste : \esp{Hay muchos hoteles en Kribi.}
\end{attention}

\courssub{Règle 2 : ser, le verbe de l'identité et de la caractéristique permanente}

\begin{propriete}[Ser : qui c'est, comment c'est (durablement)]
Le verbe \esp{ser} s'emploie pour exprimer :
\begin{itemize}
\item l'\cle{identité}, la définition, la nature : \esp{Kribi es una ciudad turística.} « Kribi est une ville touristique. »
\item la \cle{caractéristique permanente}, la qualité durable d'un lieu ou d'une chose : \esp{El paisaje es hermoso.} « Le paysage est magnifique. » \esp{El monte Camerún es alto.} « Le mont Cameroun est haut. »
\item l'\cle{origine} avec la préposition \esp{de} : \esp{El guía es de Douala.} « Le guide est de Douala. » \esp{Soy de Yaoundé.} « Je suis de Yaoundé. »
\item la \cle{nationalité} et la profession : \esp{Los turistas son franceses.} « Les touristes sont français. » \esp{Ella es guía.} « Elle est guide. »
\end{itemize}
\end{propriete}

L'idée directrice est simple : \esp{ser} répond à la question « \emph{qu'est-ce que c'est ? comment est-ce, par nature ?} ». Si tu peux remplacer le verbe par « est par nature », « est défini comme », alors c'est \esp{ser}.

\begin{exemplebox}[Ser dans la description d'un lieu]
\esp{El parque de Waza es muy grande.} « Le parc de Waza est très grand. »\par
\esp{Las playas de Kribi son magníficas.} « Les plages de Kribi sont magnifiques. »\par
\esp{El patrimonio de España es muy rico.} « Le patrimoine de l'Espagne est très riche. »\par
\esp{El ecoturismo es una actividad responsable.} « L'écotourisme est une activité responsable. »
\end{exemplebox}

\courssub{Règle 3 : estar, le verbe de la localisation et de l'état temporaire}

\begin{propriete}[Estar : où c'est, dans quel état c'est]
Le verbe \esp{estar} s'emploie pour exprimer :
\begin{itemize}
\item la \cle{localisation}, la position dans l'espace (souvent avec \esp{en}, \esp{cerca de}, \esp{lejos de}, \esp{al lado de}) : \esp{El hotel está cerca de la playa.} « L'hôtel est près de la plage. » \esp{Waza está en el norte de Camerún.} « Waza est dans le nord du Cameroun. »
\item l'\cle{état temporaire}, la situation du moment : \esp{El museo está cerrado hoy.} « Le musée est fermé aujourd'hui. » \esp{Los turistas están contentos.} « Les touristes sont contents (en ce moment). »
\end{itemize}
\end{propriete}

L'idée directrice : \esp{estar} répond à la question « \emph{où est-ce ? dans quel état est-ce maintenant ?} ». Pour un lieu, la localisation s'exprime \emph{toujours} avec \esp{estar}, même si elle est permanente : on dit \esp{El monte Camerún está en el suroeste} « Le mont Cameroun est dans le sud-ouest », jamais \esp{es en el suroeste}.

\begin{exemplebox}[Estar dans la description d'un lieu]
\esp{El parque está en el norte.} « Le parc est dans le nord. »\par
\esp{La reserva natural está lejos de la ciudad.} « La réserve naturelle est loin de la ville. »\par
\esp{El hotel está abierto.} « L'hôtel est ouvert. »\par
\esp{Estamos en Kribi con un grupo de turistas.} « Nous sommes à Kribi avec un groupe de touristes. »
\end{exemplebox}

\begin{attention}[Exception utile : la localisation d'un événement]
La règle « localisation = \esp{estar} » vaut pour les \emph{choses} et les \emph{personnes}. Mais pour un \emph{événement} (fête, match, concert, réunion), l'espagnol utilise \esp{ser} : \esp{El partido es en el estadio.} « Le match a lieu au stade. » \esp{La fiesta es en Kribi.} « La fête a lieu à Kribi. » Retiens : une chose \emph{se trouve} quelque part (\esp{estar}), un événement \emph{a lieu} quelque part (\esp{ser}).
\end{attention}

\courssub{Les formes au présent : ser, estar, hay}

\begin{aretenir}[Conjugaison au présent de l'indicatif]
\begin{center}
\begin{tabular}{|l|l|l|l|}
\hline
\rowcolor{popblueL} \textbf{Personne} & \textbf{ser} & \textbf{estar} & \textbf{hay} \\
\hline
yo & soy & estoy & \\
\hline
tú & eres & estás & \\
\hline
él / ella / usted & es & está & hay \\
\hline
nosotros / -as & somos & estamos & (invariable) \\
\hline
vosotros / -as & sois & estáis & \\
\hline
ellos / ellas / ustedes & son & están & \\
\hline
\end{tabular}
\end{center}
Attention aux accents : \esp{estás}, \esp{está}, \esp{estáis}, \esp{están} portent un accent écrit ; les formes de \esp{ser} n'en portent pas.
\end{aretenir}

\courssub{Comparaison français / espagnol}

\begin{center}
\begin{tabularx}{\linewidth}{|X|X|X|}
\hline
\rowcolor{popblueL} \textbf{Français} & \textbf{Espagnol} & \textbf{Remarque} \\
\hline
Il y a des plages. & Hay playas. & « il y a » = \esp{hay}, invariable \\
\hline
Kribi est une ville touristique. & Kribi es una ciudad turística. & identité $\rightarrow$ \esp{ser} \\
\hline
Le paysage est magnifique. & El paisaje es hermoso. & qualité durable $\rightarrow$ \esp{ser} \\
\hline
L'hôtel est près de la plage. & El hotel está cerca de la playa. & localisation $\rightarrow$ \esp{estar} \\
\hline
Le musée est fermé aujourd'hui. & El museo está cerrado hoy. & état temporaire $\rightarrow$ \esp{estar} \\
\hline
Le guide est de Douala. & El guía es de Douala. & origine : \esp{ser + de} \\
\hline
Il faut protéger la nature. & Hay que proteger la naturaleza. & obligation générale : \esp{hay que} + infinitif \\
\hline
\end{tabularx}
\end{center}

Le français dit « être » partout ; l'espagnol partage le travail entre \esp{ser} et \esp{estar}. C'est toi qui dois choisir, à chaque phrase, selon le sens.

\courssub{Cas particuliers utiles}

\begin{propriete}[Trois constructions à mémoriser]
\begin{enumerate}
\item \textbf{\esp{ser + de}} : origine, provenance. \esp{Este guía es de Guinea Ecuatorial.} « Ce guide est de Guinée équatoriale. »
\item \textbf{\esp{estar + en / cerca de / lejos de}} : localisation. \esp{La oficina de turismo está en el centro.} « L'office de tourisme est au centre. »
\item \textbf{\esp{hay que + infinitif}} : obligation générale, « il faut ». Très utile pour parler de tourisme durable : \esp{Hay que proteger la naturaleza.} « Il faut protéger la nature. » \esp{Hay que respetar el patrimonio.} « Il faut respecter le patrimoine. »
\end{enumerate}
\end{propriete}

\begin{savaistu}[Pourquoi deux verbes « être » ?]
\esp{Ser} descend du latin \emph{esse} (être par essence) et \esp{estar} du latin \emph{stare} (se tenir debout, se trouver quelque part). L'espagnol a gardé les deux verbes latins, alors que le français les a fondus en un seul : « être ». C'est pourquoi cette distinction n'existe pas dans ta langue et demande un vrai effort de réflexion. Bonne nouvelle : une fois la logique comprise (nature contre position/état), elle devient automatique.
\end{savaistu}

\courssub{Figure : l'organigramme de choix}

\begin{popfigure}
\begin{center}
\begin{tikzpicture}[
  node distance=0.9cm and 1.6cm,
  q/.style={diamond, aspect=2.4, draw=popblue, fill=popblueL, align=center, inner sep=1.5pt, font=\small},
  r/.style={rectangle, rounded corners, draw=poporange, fill=poporangeL, align=center, inner sep=4pt, font=\small},
  f/.style={-{Stealth[length=2.5mm]}, thick, popdark}
]
\node[q, align=center] (q1){Existence ?\\« il y a » ?};
\node[r, right=of q1, align=center] (r1){\textbf{hay}\\Hay elefantes.};
\node[q, below=of q1, align=center] (q2){Identité, qualité\\permanente ?};
\node[r, right=of q2, align=center] (r2){\textbf{ser}\\El parque es grande.};
\node[q, below=of q2, align=center] (q3){Localisation\\ou état ?};
\node[r, right=of q3, align=center] (r3){\textbf{estar}\\El parque está en el norte.};
\draw[f] (q1) -- node[above, font=\scriptsize]{oui} (r1);
\draw[f] (q1) -- node[left, font=\scriptsize]{non} (q2);
\draw[f] (q2) -- node[above, font=\scriptsize]{oui} (r2);
\draw[f] (q2) -- node[left, font=\scriptsize]{non} (q3);
\draw[f] (q3) -- node[above, font=\scriptsize]{oui} (r3);
\end{tikzpicture}
\end{center}
\legende{Organigramme de choix : existence $\rightarrow$ hay ; identité ou qualité durable $\rightarrow$ ser ; localisation ou état $\rightarrow$ estar.}
\end{popfigure}

\courssub{Méthode : comment choisir entre ser et estar}

\begin{methode}[La méthode des deux questions]
Pour chaque phrase, suis ces étapes dans l'ordre :
\begin{enumerate}
\item \textbf{Étape 1 :} la phrase exprime-t-elle « il y a » ? Si oui, utilise \esp{hay} (invariable, même devant un pluriel).
\item \textbf{Étape 2 :} sinon, pose la question : « est-ce une définition, une qualité durable, une origine ? » Si oui, utilise \esp{ser}. Exemple : \esp{El monte Camerún es alto} (le mont ne changera pas de taille demain).
\item \textbf{Étape 3 :} sinon, pose la question : « est-ce une position dans l'espace ou un état du moment ? » Si oui, utilise \esp{estar}. Exemple : \esp{El monte Camerún está en el suroeste} (position), \esp{El hotel está lleno} (état du moment : il est plein aujourd'hui).
\item \textbf{Étape 4 :} vérifie l'accord du verbe avec le sujet (\esp{es} / \esp{son} ; \esp{está} / \esp{están}) et l'accent écrit sur les formes de \esp{estar}.
\end{enumerate}
\end{methode}

\begin{exoresolu}[Choisir le bon verbe]
Énoncé : complète avec \esp{hay}, la bonne forme de \esp{ser} ou de \esp{estar}, puis traduis.\par
1. En Kribi \_\_\_\_\_\_ playas magníficas.\par
2. El paisaje del monte Camerún \_\_\_\_\_\_ hermoso.\par
3. El parque de Waza \_\_\_\_\_\_ en el norte del país.\par
4. Los turistas \_\_\_\_\_\_ contentos con el guía.\par
5. El guía \_\_\_\_\_\_ de Douala.\par
6. \_\_\_\_\_\_ que proteger las reservas naturales.
\tcblower
Solution détaillée :\par
1. \esp{En Kribi hay playas magníficas.} « À Kribi il y a des plages magnifiques. » — existence, « il y a » $\rightarrow$ \esp{hay} (invariable malgré le pluriel \esp{playas}).\par
2. \esp{El paisaje del monte Camerún es hermoso.} « Le paysage du mont Cameroun est magnifique. » — qualité durable $\rightarrow$ \esp{ser}, 3\up{e} personne du singulier : \esp{es}.\par
3. \esp{El parque de Waza está en el norte del país.} « Le parc de Waza est dans le nord du pays. » — localisation $\rightarrow$ \esp{estar} : \esp{está} (avec accent !).\par
4. \esp{Los turistas están contentos con el guía.} « Les touristes sont contents du guide. » — état du moment, sujet pluriel $\rightarrow$ \esp{están}.\par
5. \esp{El guía es de Douala.} « Le guide est de Douala. » — origine $\rightarrow$ \esp{ser + de} : \esp{es}.\par
6. \esp{Hay que proteger las reservas naturales.} « Il faut protéger les réserves naturelles. » — obligation générale $\rightarrow$ \esp{hay que} + infinitif.
\end{exoresolu}

\begin{attention}[Les trois erreurs classiques des francophones]
\begin{enumerate}
\item Traduire « il y a » par \esp{tiene} ou \esp{es}. Faux : \esp{Tiene muchos hoteles.} Juste : \esp{Hay muchos hoteles.}
\item Utiliser \esp{ser} pour la localisation d'une chose. Faux : \esp{El hotel es cerca de la playa.} Juste : \esp{El hotel está cerca de la playa.} En espagnol, un lieu \emph{se trouve} quelque part (\esp{estar}) ; \esp{ser} ne sert pas à situer une chose.
\item Faire varier \esp{hay} devant un pluriel. Faux : \esp{Han muchos turistas.} Juste : \esp{Hay muchos turistas.} \esp{Hay} ne change jamais au présent.
\end{enumerate}
\end{attention}

\begin{exercice}[À toi de jouer]
Complète avec \esp{hay}, \esp{ser} ou \esp{estar} (conjugué au présent), puis traduis chaque phrase en français.\par
1. En la región del Norte \_\_\_\_\_\_ un parque nacional famoso.\par
2. El parque \_\_\_\_\_\_ muy grande y hermoso.\par
3. El parque \_\_\_\_\_\_ lejos de la capital.\par
4. Los hoteles de Kribi \_\_\_\_\_\_ modernos.\par
5. Mi familia y yo \_\_\_\_\_\_ en un hotel cerca de la playa.\par
6. \_\_\_\_\_\_ que cuidar el medio ambiente.
\end{exercice}

\begin{corrige}[Exercice « À toi de jouer »]
1. \esp{hay} — « Dans la région du Nord il y a un parc national célèbre. » (existence)\par
2. \esp{es} — « Le parc est très grand et magnifique. » (qualité durable)\par
3. \esp{está} — « Le parc est loin de la capitale. » (localisation)\par
4. \esp{son} — « Les hôtels de Kribi sont modernes. » (caractéristique, pluriel)\par
5. \esp{estamos} — « Ma famille et moi sommes dans un hôtel près de la plage. » (localisation, 1\up{re} personne du pluriel)\par
6. \esp{Hay} — « Il faut prendre soin de l'environnement. » (\esp{hay que} + infinitif)
\end{corrige}

\begin{exercice}[Thème : traduis en espagnol]
1. Il y a beaucoup de touristes à Kribi en décembre.\par
2. Le mont Cameroun est très haut et il est dans le sud-ouest du pays.\par
3. Les plages de Kribi sont magnifiques.\par
4. L'office de tourisme est près de l'hôtel.\par
5. Il faut respecter la nature et le patrimoine.
\end{exercice}

\begin{corrige}[Exercice « Thème »]
1. \esp{Hay muchos turistas en Kribi en diciembre.} (« il y a » $\rightarrow$ \esp{hay}, invariable ; « en décembre » = \esp{en diciembre}, sans majuscule au nom du mois.)\par
2. \esp{El monte Camerún es muy alto y está en el suroeste del país.} (qualité durable $\rightarrow$ \esp{es} ; localisation $\rightarrow$ \esp{está} : les deux verbes dans la même phrase !)\par
3. \esp{Las playas de Kribi son magníficas.} (caractéristique, sujet pluriel $\rightarrow$ \esp{son}.)\par
4. \esp{La oficina de turismo está cerca del hotel.} (localisation $\rightarrow$ \esp{está} ; \esp{de + el} se contracte en \esp{del}.)\par
5. \esp{Hay que respetar la naturaleza y el patrimonio.} (« il faut » $\rightarrow$ \esp{hay que} + infinitif.)
\end{corrige}

\begin{aretenir}[L'essentiel de la section]
\begin{itemize}
\item \esp{Hay} = « il y a » : existence, forme invariable, même devant un pluriel.
\item \esp{Ser} (soy, eres, es, somos, sois, son) : identité, qualité permanente, origine (\esp{ser + de}), nationalité, profession ; aussi la localisation d'un \emph{événement}.
\item \esp{Estar} (estoy, estás, está, estamos, estáis, están) : localisation des choses et des personnes (\esp{estar + en / cerca de / lejos de}) et état temporaire.
\item \esp{Hay que + infinitif} = « il faut » : parfait pour promouvoir le tourisme durable.
\item Le français dit « être » partout ; toi, tu choisis selon le sens : nature $\rightarrow$ \esp{ser}, position ou état $\rightarrow$ \esp{estar}.
\end{itemize}
\end{aretenir}

Tu maîtrises maintenant les trois verbes de la description. Dans la section suivante, nous enrichirons ces phrases grâce aux \cle{adjectifs}, pour passer d'une description correcte à une description vivante et attractive d'un site touristique.

\coursec{Les adjectifs dans la description}

Dans la section précédente, nous avons appris à décrire un lieu avec \esp{hay}, \esp{ser} et \esp{estar} : \esp{En Kribi hay playas magníficas. El paisaje es hermoso. El parque está cerca.} Mais regardez bien ces phrases : pour rendre la description vivante et attirante, il ne suffit pas de dire qu'un lieu \emph{existe} ou \emph{où il est} ; il faut le \emph{qualifier}. C'est le rôle des \cle{adjectifs} : \esp{magníficas}, \esp{hermoso}, \esp{grande}... Or l'adjectif espagnol obéit à deux règles que le francophone oublie souvent : il \textbf{s'accorde} avec le nom et il se place en général \textbf{après} lui. C'est tout l'objet de cette section.

\courssub{Observation : repérons les adjectifs dans un texte}

\begin{textebox}[Un folleto turístico : « Descubre el Camerún »]
El Camerún es un país turístico extraordinario. En el sur, la ciudad de Kribi tiene playas hermosas y tranquilas, con un mar azul y una arena blanca. Es un lugar ideal para el ecoturismo.

En el oeste, el monte Camerún es un volcán impresionante : es la montaña más alta del África occidental. Los paisajes verdes de la región son magníficos y la gente local es muy amable.

En el norte, el parque nacional de Waza es una reserva natural famosa. Allí hay animales salvajes : elefantes grandes, jirafas elegantes y leones impresionantes. Los guías son jóvenes y competentes.

El Camerún tiene también un patrimonio cultural rico : mercados coloridos, música tradicional y una gastronomía deliciosa. Visitar el Camerún es una experiencia inolvidable. ¡Ven a descubrir este gran país !
\end{textebox}

\textbf{Glossaire :} \esp{la arena} = le sable ; \esp{el volcán} = le volcan ; \esp{la gente} = les gens ; \esp{amable} = gentil(le) ; \esp{la reserva natural} = la réserve naturelle ; \esp{salvaje} = sauvage ; \esp{la jirafa} = la girafe ; \esp{colorido} = coloré ; \esp{inolvidable} = inoubliable.

\textbf{Questions de compréhension :}
\begin{enumerate}
\item ¿Qué playas tiene Kribi ? (Kribi a des plages magnifiques et tranquilles.)
\item ¿Dónde está el parque nacional de Waza ? (Il est dans le nord du Cameroun.)
\item ¿Qué animales hay en Waza ? (Il y a des éléphants, des girafes et des lions.)
\end{enumerate}

\textbf{Observons.} Soulignons les adjectifs du texte : \esp{playas \textbf{hermosas} y \textbf{tranquilas}}, \esp{un mar \textbf{azul}}, \esp{una arena \textbf{blanca}}, \esp{un volcán \textbf{impresionante}}, \esp{paisajes \textbf{verdes}}, \esp{elefantes \textbf{grandes}}, \esp{un \textbf{gran} país}. Deux constatations s'imposent :
\begin{enumerate}
\item L'adjectif change de forme selon le nom : \esp{mar azul} mais \esp{arena blanca} ; \esp{volcán impresionante} mais \esp{playas hermosas}.
\item L'adjectif se place presque toujours \textbf{après} le nom : on dit \esp{playas hermosas} et non \esp{hermosas playas}.
\end{enumerate}

\courssub{L'accord de l'adjectif : genre et nombre}

\begin{propriete}[Règle d'accord de l'adjectif]
En espagnol, l'adjectif \textbf{s'accorde en genre et en nombre} avec le nom qu'il qualifie, exactement comme en français, mais avec des formes écrites plus visibles :
\begin{itemize}
\item masculin singulier : \esp{un paisaje \textbf{hermoso}} (un beau paysage) ;
\item féminin singulier : \esp{una playa \textbf{hermosa}} (une belle plage) ;
\item masculin pluriel : \esp{unos paisajes \textbf{hermosos}} (de beaux paysages) ;
\item féminin pluriel : \esp{unas playas \textbf{hermosas}} (de belles plages).
\end{itemize}
Les adjectifs terminés en \textbf{-o} forment le féminin en \textbf{-a} et le pluriel en \textbf{-os / -as}. Les adjectifs terminés en \textbf{-e} ou en consonne (\esp{verde}, \esp{grande}, \esp{azul}, \esp{importante}, \esp{sostenible}) sont \textbf{invariables en genre} ; ils prennent seulement la marque du pluriel : \esp{verde} $\rightarrow$ \esp{verdes}, \esp{azul} $\rightarrow$ \esp{azules}.
\end{propriete}

\begin{center}
\begin{tabularx}{\linewidth}{|l|X|X|X|}
\hline
\rowcolor{popblueL} Adjectif & Masc. sing. & Fém. sing. & Pluriel (m./f.) \\
\hline
hermoso (beau) & un hotel hermoso & una playa hermosa & hoteles hermosos / playas hermosas \\
\hline
verde (vert) & un bosque verde & una montaña verde & bosques verdes / montañas verdes \\
\hline
grande (grand) & un parque grande & una reserva grande & parques grandes / reservas grandes \\
\hline
azul (bleu) & un mar azul & una laguna azul & mares azules / lagunas azules \\
\hline
importante & un museo importante & una ciudad importante & museos / ciudades importantes \\
\hline
\end{tabularx}
\end{center}

\begin{exemplebox}[Exemples traduits]
\esp{Kribi tiene playas hermosas y tranquilas.} « Kribi a des plages magnifiques et tranquilles. »

\esp{El monte Camerún es un volcán impresionante.} « Le mont Cameroun est un volcan impressionnant. »

\esp{En Waza hay animales salvajes.} « À Waza, il y a des animaux sauvages. »

\esp{Madrid es una ciudad famosa y moderna.} « Madrid est une ville célèbre et moderne. »

\esp{Los hoteles nuevos son cómodos.} « Les hôtels nouveaux sont confortables. »
\end{exemplebox}

\courssub{La place de l'adjectif : après le nom}

\begin{propriete}[Place normale de l'adjectif]
En espagnol, l'adjectif qualificatif se place normalement \textbf{APRÈS} le nom, contrairement au français où il est souvent placé avant :

\esp{un hotel moderno} = un hôtel moderne ; \esp{una playa magnífica} = une plage magnifique ; \esp{unos paisajes impresionantes} = des paysages impressionnants.

Placé avant le nom, l'adjectif prend une valeur littéraire, affective ou exclamative : \esp{la hermosa ciudad de mi infancia} (la belle ville de mon enfance, style poétique). Au niveau de la Première, retenez la règle simple : \textbf{nom + adjectif}.
\end{propriete}

\begin{popfigure}
\begin{tikzpicture}[
  boxfr/.style={draw=poporange, fill=poporangeL, rounded corners=2pt, minimum height=0.8cm, align=center, font=\small},
  boxes/.style={draw=popblue, fill=popblueL, rounded corners=2pt, minimum height=0.8cm, align=center, font=\small},
  lab/.style={font=\bfseries\small}
]
\node[lab, text=poporange] at (-4,1.6) {FRANÇAIS};
\node[boxfr, align=center] (a1) at (-6,0.6){adjectif\\ magnifique};
\node[boxfr, align=center] (a2) at (-3,0.6){nom\\ plage};
\draw[-{Stealth}, thick, poporange] (a1) -- (a2);
\node[lab, text=popblue] at (3.5,1.6) {ESPAÑOL};
\node[boxes, align=center] (b1) at (1.5,0.6){nom\\ playa};
\node[boxes, align=center] (b2) at (5.5,0.6){adjectif\\ magnífica};
\draw[-{Stealth}, thick, popblue] (b1) -- (b2);
\node[align=center, font=\small, text=popdark] at (0,-0.6) {une plage magnifique $\longrightarrow$ una playa magnífica};
\end{tikzpicture}
\legende{Schéma comparatif : en français l'adjectif précède souvent le nom ; en espagnol il le suit.}
\end{popfigure}

\begin{center}
\begin{tabularx}{\linewidth}{|X|X|X|}
\hline
\rowcolor{popblueL} Français & Español & Remarque \\
\hline
une plage magnifique & una playa magnífica & inversion de place \\
\hline
un hôtel moderne & un hotel moderno & nom + adjectif \\
\hline
un parc célèbre & un parque famoso & nom + adjectif \\
\hline
des paysages verts & paisajes verdes & accord au pluriel \\
\hline
une ville tranquille & una ciudad tranquila & adjectif en -o $\rightarrow$ -a \\
\hline
\end{tabularx}
\end{center}

\begin{attention}[Piège de place]
Ne pas écrire \esp{una magnífica playa} en imitation du français « une magnifique plage ». La place normale est \esp{una playa magnífica}. De même : \esp{un hotel moderno} et non \esp{un moderno hotel}. La place avant le nom est réservée au style littéraire ou à quelques adjectifs spéciaux (voir l'apocope ci-dessous).
\end{attention}

\courssub{L'apocope : gran, buen, mal}

\begin{propriete}[L'apocope]
Certains adjectifs perdent leur dernière syllabe (\emph{apocope}) devant un nom singulier :
\begin{itemize}
\item \esp{grande} $\rightarrow$ \esp{\textbf{gran}} devant tout nom singulier (masculin ou féminin) : \esp{un gran hotel}, \esp{una gran ciudad}. Placé avant le nom, \esp{gran} signifie « grand » au sens figuré (remarquable, important) ; placé après, \esp{grande} garde le sens de « grand » en taille : \esp{un hotel grande} = un hôtel de grande taille.
\item \esp{bueno} $\rightarrow$ \esp{\textbf{buen}} et \esp{malo} $\rightarrow$ \esp{\textbf{mal}} devant un nom \textbf{masculin singulier} uniquement : \esp{un buen guía} (un bon guide), \esp{un mal día} (un mauvais jour), mais \esp{una buena playa}, \esp{una mala experiencia}.
\end{itemize}
\end{propriete}

\begin{exemplebox}[Exemples traduits]
\esp{Waza es un gran parque nacional.} « Waza est un grand parc national. »

\esp{El Camerún es un gran país.} « Le Cameroun est un grand pays. »

\esp{El guía es un buen profesional.} « Le guide est un bon professionnel. »

\esp{Kribi es una gran ciudad turística.} « Kribi est une grande ville touristique. »
\end{exemplebox}

\begin{attention}[Piège de sens]
\esp{Un hotel grande} (adjectif après) = un hôtel de grande taille. \esp{Un gran hotel} (adjectif avant, apocope) = un hôtel remarquable, prestigieux. Attention aussi au faux ami : \esp{grande} ne signifie jamais « grand » au sens de « grand homme politique » sans nuance de taille si on le place après ; et \esp{largo} = long, pas « large » (\esp{ancho}).
\end{attention}

\courssub{Adjectifs utiles pour décrire un site touristique}

\begin{definition}[Vocabulaire descriptif du tourisme]
\esp{hermoso/a} = beau/belle ; \esp{magnífico/a} = magnifique ; \esp{impresionante} = impressionnant ; \esp{famoso/a} = célèbre ; \esp{tranquilo/a} = tranquille ; \esp{salvaje} = sauvage ; \esp{turístico/a} = touristique ; \esp{limpio/a} = propre ; \esp{acogedor/a} = accueillant ; \esp{exótico/a} = exotique ; \esp{sostenible} = durable ; \esp{inolvidable} = inoubliable.
\end{definition}

\begin{exoresolu}[Accorder et placer les adjectifs]
\textbf{Énoncé.} Complétez avec la bonne forme de l'adjectif entre parenthèses et à la bonne place, puis traduisez :
\begin{enumerate}
\item Kribi tiene unas playas \_\_\_\_\_\_\_\_ (hermoso).
\item El parque de Waza es una reserva \_\_\_\_\_\_\_\_ (natural) muy \_\_\_\_\_\_\_\_ (famoso).
\item El monte Camerún es un volcán \_\_\_\_\_\_\_\_ (impresionante).
\item Madrid es una ciudad \_\_\_\_\_\_\_\_ (grande) y \_\_\_\_\_\_\_\_ (moderno).
\end{enumerate}
\tcblower
\textbf{Solution détaillée.}
\begin{enumerate}
\item \esp{Kribi tiene unas playas \textbf{hermosas}.} — \esp{playas} est féminin pluriel, donc \esp{hermoso} $\rightarrow$ \esp{hermosas}. « Kribi a de belles plages. »
\item \esp{El parque de Waza es una reserva \textbf{natural} muy \textbf{famosa}.} — \esp{natural} est invariable en genre ; \esp{famoso} s'accorde avec \esp{reserva} (féminin singulier) $\rightarrow$ \esp{famosa}. « Le parc de Waza est une réserve naturelle très célèbre. »
\item \esp{El monte Camerún es un volcán \textbf{impresionante}.} — adjectif en -e, invariable en genre ; masculin singulier, pas de marque. « Le mont Cameroun est un volcan impressionnant. »
\item \esp{Madrid es una ciudad \textbf{grande} y \textbf{moderna}.} — \esp{grande} invariable en genre ; \esp{moderno} $\rightarrow$ \esp{moderna} (féminin singulier). « Madrid est une grande ville moderne. »
\end{enumerate}
\end{exoresolu}

\begin{exercice}[À vous de décrire]
Traduisez en espagnol, en plaçant correctement les adjectifs :
\begin{enumerate}
\item Kribi est une ville tranquille et accueillante.
\item Le mont Cameroun est une montagne impressionnante.
\item Waza est un grand parc avec des animaux sauvages.
\item Les plages du Cameroun sont magnifiques et propres.
\item Le tourisme durable est important pour le pays.
\end{enumerate}
\end{exercice}

\begin{corrige}[Exercice : À vous de décrire]
\begin{enumerate}
\item \esp{Kribi es una ciudad tranquila y acogedora.} (féminin singulier : \esp{tranquila}, \esp{acogedora})
\item \esp{El monte Camerún es una montaña impresionante.} (invariable en genre)
\item \esp{Waza es un gran parque con animales salvajes.} (apocope de \esp{grande} devant nom singulier)
\item \esp{Las playas del Camerún son magníficas y limpias.} (féminin pluriel)
\item \esp{El turismo sostenible es importante para el país.} (adjectifs invariables en genre)
\end{enumerate}
\end{corrige}

\begin{aretenir}[L'essentiel]
\begin{itemize}
\item L'adjectif espagnol \textbf{s'accorde} en genre et en nombre avec le nom : \esp{hermoso / hermosa / hermosos / hermosas}.
\item Il se place normalement \textbf{après} le nom : \esp{una playa magnífica} (une plage magnifique).
\item Les adjectifs en \textbf{-e} ou en consonne (\esp{verde}, \esp{grande}, \esp{azul}, \esp{importante}, \esp{sostenible}) sont invariables en genre.
\item Apocope : \esp{gran} devant tout nom singulier (\esp{un gran país}) ; \esp{buen} et \esp{mal} devant un nom masculin singulier (\esp{un buen guía}).
\end{itemize}
\end{aretenir}

\coursec{Méthode du commentaire et commentaire modèle}

Dans la section précédente, nous avons appris à enrichir une description grâce aux adjectifs : \esp{una playa magnífica}, \esp{un paisaje impresionante}, \esp{un patrimonio único}. Ces adjectifs ne servent pas seulement à décrire : ils révèlent aussi le \cle{point de vue} de l'auteur. C'est exactement ce que l'on fait dans un \cle{commentaire de texte} : on ne se contente pas de lire, on \emph{analyse} comment l'auteur construit son idée. Cette compétence est au cœur de l'épreuve du Baccalauréat. Voyons d'abord de quoi il s'agit.

\courssub{Qu'est-ce qu'un commentaire de texte ?}

\begin{definition}[Le commentaire de texte]
Le \cle{commentaire de texte} (\esp{comentario de texto}) est un exercice écrit qui consiste à :
\begin{itemize}
\item \textbf{présenter} le texte (thème, type, source) ;
\item \textbf{dégager la thèse} de l'auteur, c'est-à-dire l'idée principale qu'il défend ;
\item \textbf{analyser} les arguments et les procédés (exemples, connecteurs, adjectifs, chiffres) ;
\item \textbf{conclure} par un bilan et une ouverture, en donnant son avis argumenté.
\end{itemize}
\end{definition}

\begin{attention}[Commentaire n'est pas résumé]
Au Baccalauréat, la faute la plus grave est de \emph{résumer} le texte au lieu de le \emph{commenter}. Un résumé répète ce que dit le texte ; un commentaire \textbf{explique comment} le texte le dit et \textbf{pourquoi} c'est efficace. Chaque affirmation doit être justifiée par une \cle{citation courte} du texte, entre guillemets. Exemple : \esp{El autor afirma que el turismo «crea empleos», lo que demuestra su valor económico.}
\end{attention}

\courssub{La structure du commentaire : le plan en trois parties}

\begin{propriete}[Le plan du commentaire]
Un commentaire de texte suit toujours la même architecture, en trois parties :
\begin{enumerate}
\item \textbf{Introducción} : présenter le texte (thème, source, type) et annoncer la thèse.
\item \textbf{Desarrollo} : deux ou trois paragraphes, chacun consacré à une idée principale, avec citations et analyse.
\item \textbf{Conclusión} : bilan de l'analyse, avis personnel argumenté, ouverture.
\end{enumerate}
\end{propriete}

\begin{popfigure}
\begin{tikzpicture}[
node distance=0.9cm,
etape/.style={rectangle, rounded corners, draw=popblue, fill=popblueL, thick, align=center, minimum width=6.5cm, minimum height=1.1cm},
fleche/.style={-{Stealth[length=3mm]}, thick, popdark}
]
\node[etape, align=center] (intro){\textbf{Introducción}\\ thème + thèse de l'auteur};
\node[etape, below=of intro, fill=poporangeL, draw=poporange, align=center] (des){\textbf{Desarrollo}\\ paragraphe 1 : arguments économiques\\ paragraphe 2 : arguments écologiques};
\node[etape, below=of des, fill=popgreenL, draw=popgreen, align=center] (conc){\textbf{Conclusión}\\ bilan + avis personnel + ouverture};
\draw[fleche] (intro) -- (des);
\draw[fleche] (des) -- (conc);
\end{tikzpicture}
\legende{Organigramme du plan de commentaire : Introducción, Desarrollo, Conclusión.}
\end{popfigure}

\begin{methode}[Commenter un texte en quatre étapes]
\begin{enumerate}
\item \textbf{Identifier le thème et la thèse} : de quoi parle le texte ? Quelle idée l'auteur défend-il ? Dans notre texte, la thèse est : \emph{le tourisme est un moteur de développement, mais il doit être durable}.
\item \textbf{Dégager deux ou trois idées principales} : ici, les arguments économiques (emplois, devises, infrastructures) et les arguments écologiques (protection du paysage, du patrimoine, des réserves naturelles).
\item \textbf{Citer le texte pour justifier} : chaque idée s'appuie sur une citation courte, analysée en une ou deux phrases.
\item \textbf{Donner son avis argumenté en conclusion} : êtes-vous d'accord ? Pourquoi ? Ouvrez sur une question plus large (l'écotourisme au Cameroun, par exemple).
\end{enumerate}
\end{methode}

\courssub{Repérer la thèse et les arguments du texte}

Relisons la fin du texte support \esp{El turismo, motor de desarrollo} et observons les deux familles d'arguments.

\textbf{1. Les arguments économiques.} L'auteur montre que le tourisme rapporte de l'argent au pays :

\begin{itemize}
\item \esp{El turismo crea empleos en los hoteles y los restaurantes.} « Le tourisme crée des emplois dans les hôtels et les restaurants. »
\item \esp{Los turistas traen divisas al país.} « Les touristes apportent des devises au pays. »
\item \esp{Gracias al turismo, se construyen carreteras y aeropuertos.} « Grâce au tourisme, on construit des routes et des aéroports. »
\end{itemize}

\textbf{2. Les arguments écologiques.} L'auteur ajoute une condition : le développement ne doit pas détruire la nature.

\begin{itemize}
\item \esp{Hay que proteger el paisaje y el patrimonio.} « Il faut protéger le paysage et le patrimoine. »
\item \esp{Las reservas naturales, como el parque de Waza, necesitan un turismo responsable.} « Les réserves naturelles, comme le parc de Waza, ont besoin d'un tourisme responsable. »
\item \esp{El ecoturismo respeta la naturaleza y ayuda a las comunidades locales.} « L'écotourisme respecte la nature et aide les communautés locales. »
\end{itemize}

Remarquez le rôle des \cle{connecteurs} : \esp{en primer lugar} annonce le premier argument, \esp{además} ajoute une idée, \esp{sin embargo} introduit la nuance écologique, \esp{en conclusión} ferme le raisonnement. Ce sont les mêmes connecteurs que vous utiliserez dans votre propre commentaire.

\begin{exemplebox}[Connecteurs et phrases utiles, traduits]
\begin{itemize}
\item \esp{En primer lugar, el turismo crea empleos.} « D'abord, le tourisme crée des emplois. »
\item \esp{Además, atrae divisas y mejora las infraestructuras.} « De plus, il attire des devises et améliore les infrastructures. »
\item \esp{Sin embargo, hay que proteger el medio ambiente.} « Cependant, il faut protéger l'environnement. »
\item \esp{En conclusión, el turismo es una riqueza si es sostenible.} « En conclusion, le tourisme est une richesse s'il est durable. »
\end{itemize}
\end{exemplebox}

\begin{center}
\begin{tabularx}{\linewidth}{|X|X|X|}
\hline
\rowcolor{popblueL} \textbf{Español} & \textbf{Français} & \textbf{Fonction} \\
\hline
en primer lugar & d'abord, en premier lieu & annoncer le premier argument \\
\hline
además & de plus, en outre & ajouter une idée \\
\hline
por ejemplo & par exemple & illustrer \\
\hline
sin embargo & cependant, pourtant & introduire une opposition \\
\hline
por eso / por lo tanto & c'est pourquoi / donc & exprimer la conséquence \\
\hline
en conclusión & en conclusion & conclure \\
\hline
\end{tabularx}
\end{center}

\begin{attention}[Pièges de traduction des connecteurs]
\begin{itemize}
\item \esp{sin embargo} signifie « cependant », jamais « sans embargo ». C'est un \cle{faux ami} total.
\item \esp{en primer lugar} se traduit « d'abord » ou « en premier lieu » ; ne dites pas \esp{en primer sitio} dans ce sens.
\item \esp{además} prend toujours un accent sur le \emph{á} ; sans accent, \esp{ademas} est une faute.
\end{itemize}
\end{attention}

\courssub{Commentaire modèle}

Voici un commentaire court (huit à dix lignes), rédigé selon le plan en trois parties. Lisez-le attentivement : repérez l'introduction, les deux idées du développement et la conclusion.

\begin{textebox}[Comentario modelo : El turismo, motor de desarrollo]
Este texto trata del turismo como motor de desarrollo económico. El autor defiende una idea clara: el turismo es una riqueza, pero debe ser sostenible.

En primer lugar, el texto muestra los beneficios económicos del turismo. El autor afirma que el turismo «crea empleos» y «atrae divisas», por ejemplo en Kribi y en Waza. Además, gracias a los turistas, se construyen carreteras y hoteles, lo que mejora la vida de la población.

Sin embargo, el autor insiste en la necesidad del ecoturismo para proteger el paisaje y el patrimonio. La expresión «turismo responsable» resume bien esta idea: la naturaleza es un tesoro frágil.

En conclusión, el turismo es una riqueza si respeta la naturaleza. En mi opinión, el Camerún tiene un potencial turístico enorme, y el ecoturismo es el camino del futuro.
\end{textebox}

\textbf{Glossaire} : \esp{trata de} = il s'agit de ; \esp{defiende} = il défend ; \esp{sostenible} = durable ; \esp{beneficios} = bénéfices, avantages ; \esp{lo que} = ce qui ; \esp{tesoro} = trésor ; \esp{en mi opinión} = à mon avis ; \esp{el camino} = le chemin, la voie.

\textbf{Analyse du modèle.} Observez comment ce commentaire applique la méthode :
\begin{itemize}
\item \textbf{Introduction} : la première phrase présente le thème (\esp{trata del turismo como motor de desarrollo}) et annonce la thèse (\esp{una riqueza, pero debe ser sostenible}).
\item \textbf{Développement} : deux paragraphes, un par famille d'arguments, chacun avec une citation courte entre guillemets (\esp{«crea empleos»}, \esp{«turismo responsable»}).
\item \textbf{Conclusion} : bilan (\esp{el turismo es una riqueza si respeta la naturaleza}), avis personnel (\esp{en mi opinión}) et ouverture (\esp{el camino del futuro}).
\end{itemize}

\begin{aretenir}[Les trois qualités d'un bon commentaire]
\begin{enumerate}
\item \textbf{Structuré} : introduction, développement en deux ou trois paragraphes, conclusion.
\item \textbf{Justifié} : chaque idée s'appuie sur une citation courte du texte.
\item \textbf{Personnel} : la conclusion donne un avis argumenté, avec \esp{en mi opinión}, \esp{creo que}, \esp{pienso que}.
\end{enumerate}
\end{aretenir}

\courssub{Entraînement guidé}

\begin{exoresolu}[Rédiger l'introduction d'un commentaire]
\textbf{Énoncé.} À partir du texte \esp{El turismo, motor de desarrollo}, rédigez l'introduction d'un commentaire (trois ou quatre phrases) : présentez le thème, le type de texte et la thèse de l'auteur. Utilisez au moins deux des expressions suivantes : \esp{trata de}, \esp{el autor defiende}, \esp{según el texto}.
\tcblower
\textbf{Solution détaillée.} Une bonne introduction répond à trois questions : de quoi parle le texte ? quel type de texte est-ce ? quelle idée l'auteur défend-il ?

\esp{Este texto es un artículo de opinión sobre el turismo. Trata del turismo como motor de desarrollo económico en el Camerún y en los países hispanohablantes. Según el texto, el turismo crea empleos y atrae divisas, por ejemplo en Kribi y en Waza. Sin embargo, el autor defiende una idea esencial: el turismo debe ser sostenible para proteger el paisaje y el patrimonio.}

Traduction : « Ce texte est un article d'opinion sur le tourisme. Il traite du tourisme comme moteur de développement économique au Cameroun et dans les pays hispanophones. Selon le texte, le tourisme crée des emplois et attire des devises, par exemple à Kribi et à Waza. Cependant, l'auteur défend une idée essentielle : le tourisme doit être durable pour protéger le paysage et le patrimoine. »

Remarquez la progression : type de texte, thème, arguments, thèse. Le connecteur \esp{sin embargo} prépare déjà l'analyse du développement.
\end{exoresolu}

\begin{exercice}[À vous de rédiger]
Rédigez la \textbf{conclusion} d'un commentaire sur le même texte (quatre ou cinq phrases). Votre conclusion doit contenir : un bilan commençant par \esp{en conclusión}, votre avis personnel avec \esp{en mi opinión} ou \esp{creo que}, et une ouverture sur l'avenir du tourisme durable au Cameroun.
\end{exercice}

\begin{corrige}[Exercice : la conclusion]
\textbf{Corrigé de méthode.} Il n'existe pas une seule bonne réponse, mais toute bonne conclusion suit le schéma bilan, avis, ouverture. Modèle possible :

\esp{En conclusión, este texto demuestra que el turismo es una gran riqueza para un país: crea empleos, atrae divisas y mejora las infraestructuras. Sin embargo, esta riqueza es frágil si no respetamos la naturaleza. En mi opinión, el ecoturismo es la mejor solución, porque protege el paisaje y ayuda a las comunidades locales. El futuro del turismo en el Camerún depende de nuestra capacidad de unir desarrollo y respeto del medio ambiente.}

Traduction : « En conclusion, ce texte démontre que le tourisme est une grande richesse pour un pays : il crée des emplois, attire des devises et améliore les infrastructures. Cependant, cette richesse est fragile si nous ne respectons pas la nature. À mon avis, l'écotourisme est la meilleure solution, parce qu'il protège le paysage et aide les communautés locales. L'avenir du tourisme au Cameroun dépend de notre capacité à unir développement et respect de l'environnement. »

Barème indicatif : bilan avec \esp{en conclusión} (1 point), avis personnel (1 point), ouverture (1 point), correction grammaticale et connecteurs (1 point).
\end{corrige}

\begin{savaistu}[Le commentaire de texte, une tradition hispanique]
En Espagne, le \esp{comentario de texto} est l'épreuve reine de la \esp{Selectividad}, l'équivalent du Baccalauréat : chaque candidat analyse un texte littéraire ou journalistique en une heure environ. En Guinée équatoriale, seul pays hispanophone d'Afrique, cet exercice est aussi au programme du lycée. En maîtrisant le commentaire en Première, vous préparez donc à la fois le Baccalauréat camerounais et une compétence universelle du monde hispanophone.
\end{savaistu}

Nous savons désormais analyser le texte comme de véritables critiques : thèse repérée, arguments classés, citations choisies, avis personnel argumenté. Il ne reste plus qu'à passer de l'analyse à la création : dans la dernière section, vous deviendrez vous-mêmes guides et ambassadeurs, en présentant et en promouvant un site touristique du Cameroun.

\coursec{El turismo y el desarrollo económico}

\courssub{Texte support : El turismo, motor de desarrollo}

\begin{textebox}[El turismo, motor de desarrollo]
El turismo es un sector clave para la economía de muchos países. En España, por ejemplo, representa más del 12\,\% del PIB y genera millones de empleos. En América Latina, países como México, Costa Rica o Perú viven en gran parte de la llegada de visitantes internacionales.\par
Camerún, por su parte, posee un potencial turístico excepcional: playas paradisíacas en Kribi y Limbé, el imponente monte Camerún, la sabana y la fauna del parque de Waza, sin olvidar la riqueza cultural de los reinos del Oeste como Foumban. Sin embargo, el sector sigue poco desarrollado comparado con sus posibilidades.\par
Para que el turismo sea un verdadero motor de desarrollo, debe ser \textbf{sostenible}. Esto significa proteger el medio ambiente (\esp{proteger el medio ambiente}), respetar las culturas locales (\esp{respetar las culturas locales}) y asegurar que los beneficios económicos lleguen a la población (\esp{beneficiar a la población local}). El \textbf{ecoturismo} y el \textbf{turismo comunitario} son modelos prometedores: el turista descubre la naturaleza y la cultura, y la gente local gestiona los recursos y gana ingresos dignos.\par
Los gobiernos deben invertir en infraestructuras (carreteras, hoteles, formación de \esp{guías}) y promover una imagen positiva del país en el extranjero. El reto es grande, pero la recompensa lo vale: un turismo bien gestionado crea empleo, conserva el \textbf{patrimonio} natural y cultural y refuerza el orgullo nacional.
\end{textebox}

\textbf{Glossaire} : \esp{el PIB} = le PIB (Produit Intérieur Brut) ; \esp{la llegada} = l'arrivée ; \esp{el potencial} = le potentiel ; \esp{la sabana} = la savane ; \esp{la fauna} = la faune ; \esp{el reino} = le royaume ; \esp{sostenible} = durable ; \esp{el medio ambiente} = l'environnement ; \esp{el ecoturismo} = l'écotourisme ; \esp{el turismo comunitario} = le tourisme communautaire ; \esp{gestionar} = gérer ; \esp{los ingresos} = les revenus ; \esp{la infraestructura} = l'infrastructure ; \esp{el patrimonio} = le patrimoine ; \esp{el orgullo} = la fierté.

\textbf{Questions de compréhension} :
\begin{enumerate}
\item ¿Qué porcentaje del PIB representa el turismo en España? (Quel pourcentage du PIB le tourisme représente-t-il en Espagne ?)
\item Cita tres atractivos turísticos de Camerún. (Cite trois attraits touristiques du Cameroun.)
\item ¿Qué significa « turismo sostenible »? (Que signifie « tourisme durable » ?)
\item ¿Cuáles son los dos modelos prometedores citados en el texto? (Quels sont les deux modèles prometteurs cités ?)
\item ¿Qué deben hacer los gobiernos para desarrollar el sector? (Que doivent faire les gouvernements pour développer le secteur ?)
\end{enumerate}

\courssub{Lexique thématique du tourisme}

\begin{definition}[Vocabulaire essentiel : el turismo]
Le tableau suivant regroupe les mots indispensables pour parler du tourisme, classés par catégories. Apprends-les avec leur genre et leur orthographe exacte (accents, \~n).
\end{definition}

\begin{center}
\begin{tabularx}{\linewidth}{|l|X|X|}\hline
\rowcolor{popblueL} \textbf{Categoría} & \textbf{Español} & \textbf{Français} \\ \hline
Personnes & \esp{el turista, la turista} (m/f invariable), \esp{el guía, la guía}, \esp{el viajero, la viajera} & le touriste, le guide, le voyageur \\ \hline
Lieux & \esp{el hotel}, \esp{la playa}, \esp{la montaña}, \esp{el paisaje}, \esp{la reserva natural}, \esp{el parque nacional}, \esp{el patrimonio}, \esp{la cascada}, \esp{la selva}, \esp{la sabana}, \esp{el volcán} & l'hôtel, la plage, la montagne, le paysage, la réserve naturelle, le parc national, le patrimoine, la cascade, la jungle, la savane, le volcan \\ \hline
Activités & \esp{el turismo}, \esp{el ecoturismo}, \esp{el safari}, \esp{la excursión}, \esp{la visita guiada}, \esp{el senderismo}, \esp{la fotografía} & le tourisme, l'écotourisme, le safari, l'excursion, la visite guidée, la randonnée, la photo \\ \hline
Actions & \esp{visitar}, \esp{recorrer}, \esp{disfrutar}, \esp{conocer}, \esp{proteger}, \esp{conservar}, \esp{respetar}, \esp{promover}, \esp{generar empleo} & visiter, parcourir, profiter, connaître, protéger, conserver, respecter, promouvoir, créer de l'emploi \\ \hline
Adjectifs & \esp{impresionante}, \esp{único}, \esp{magnífico}, \esp{paradisíaco}, \esp{salvaje}, \esp{auténtico}, \esp{sostenible}, \esp{local} & impressionnant, unique, magnifique, paradisiaque, sauvage, authentique, durable, local \\ \hline
\end{tabularx}
\end{center}

\begin{attention}[Faux amis et pièges orthographiques]
\begin{itemize}
\item \esp{El hotel} (masculin) $\neq$ \esp{el hostal} (auberge de jeunesse, plus simple). Pas de \og h \fg{} aspiré en espagnol.
\item \esp{La guía} (le guide, personne ou livre) porte l'accent sur le \emph{i} pour marquer la diphtongue \og uí \fg{} (\emph{gu-í-a}), sinon cela se lit \og guia \fg{} (une syllabe).
\item \esp{La montaña} (montagne) $\rightarrow$ pluriel \esp{las monta\~nas} (tilde sur le \~n).
\item \esp{El patrimonio} (patrimoine) : accent sur le \emph{i} (\emph{pa-tri-mo-nio}).
\item \esp{Único} (unique / seul) : accent sur le \emph{ú} pour distinguer de \esp{unico} (forme verbale inexistante ici).
\end{itemize}
\end{attention}

\courssub{Grammaire : \texorpdfstring{\esp{hay}, \esp{ser}, \esp{estar}}{hay, ser, estar} pour décrire un lieu}

\begin{propriete}[Choisir entre \esp{hay}, \esp{ser} et \esp{estar}]
Ces trois outils sont la base de toute description de lieu. Ne les confonds pas !
\end{propriete}

\begin{center}
\begin{tabularx}{\linewidth}{|l|X|X|l|}\hline
\rowcolor{popgreenL} \textbf{Outil} & \textbf{Emploi principal} & \textbf{Exemple traduit} & \textbf{Clé} \\ \hline
\esp{Hay} (verbe \esp{haber} impersonnel) & \textbf{Existence / Inventaire} : « Il y a », « Il existe ». Invariable (toujours \esp{hay}). Ne s'accorde pas avec le nom qui suit. & \esp{Hay una playa.} (Il y a une plage.)\par \esp{Hay muchos hoteles.} (Il y a beaucoup d'hôtels.) & \cle{Il y a = Hay} \\ \hline
\esp{Estar} & \textbf{Localisation / Position} : Où se trouve le lieu ? (Géographie, adresse). & \esp{Kribi está en el sur.} (Kribi est au sud.)\par \esp{El hotel está cerca del mar.} (L'hôtel est près de la mer.) & \cle{Où ? = Estar} \\ \hline
\esp{Ser} & \textbf{Identité / Caractéristique essentielle} : C'est quoi ? Quelle est sa nature, son caractère permanent ? & \esp{Es un volcán impresionante.} (C'est un volcan impressionnant.)\par \esp{Las playas son paradisíacas.} (Les plages sont paradisiaques.) & \cle{C'est quoi ? = Ser} \\ \hline
\end{tabularx}
\end{center}

\begin{exemplebox}[Le trio en action : Les chutes de la Lobé]
\esp{Las cataratas de la Lobé \textbf{están} cerca de Kribi.} (Localisation : \esp{estar})\par
\esp{Son \textbf{un espectáculo natural único}.} (Caractéristique : \esp{ser})\par
\esp{\textbf{Hay} un hotel y un restaurante cerca.} (Existence : \esp{hay})\par
\esp{Los turistas \textbf{pueden} nadar en la piscina natural.} (Activité : \esp{poder} + infinitif)
\end{exemplebox}

\begin{attention}[Erreurs fatales des francophones]
\begin{itemize}
\item \textbf{« Il y a » $\neq$ \esp{es} / \esp{está}.} Jamais ! \esp{Hay} est un verbe à part entière (forme figée de \esp{haber}).
\item \textbf{« C'est » $\neq$ \esp{está}.} \esp{Es un parque} (C'est un parc = identification). \esp{Está en el norte} (Il est au nord = position).
\item \textbf{Accord avec \esp{hay} :} \esp{Hay} ne change jamais. \esp{Hay un problema} / \esp{Hay muchos problemas}.
\item \textbf{Prépositions de lieu avec \esp{estar} :} \esp{en} (dans/pays/région), \esp{a} (à/ville), \esp{cerca de} (près de), \esp{lejos de} (loin de), \esp{al lado de} (à côté de), \esp{entre} (entre).
\end{itemize}
\end{attention}

\courssub{Les adjectifs dans la description}

\begin{propriete}[Place et accord de l'adjectif qualificatif]
En espagnol, l'adjectif qualificatif se place \textbf{généralement après le nom} (contrairement au français où il est souvent avant). Il s'accorde en \textbf{genre} (masculin/féminin) et en \textbf{nombre} (singulier/pluriel) avec le nom qu'il qualifie.
\end{propriete}

\begin{center}
\begin{tabular}{|l|l|l|l|}\hline
\rowcolor{poppurpleL} & \textbf{Masculin Singulier} & \textbf{Féminin Singulier} & \textbf{Pluriel (m/f)} \\ \hline
Terminaison en \emph{-o} & \esp{impresionant\textbf{o}} & \esp{impresionant\textbf{a}} & \esp{impresionant\textbf{os}} / \esp{impresionant\textbf{as}} \\ \hline
Terminaison en \emph{-e} ou consonne & \esp{grande} / \esp{azul} & \esp{grande} / \esp{azul} & \esp{grand\textbf{es}} / \esp{azul\textbf{es}} \\ \hline
Terminaison en \emph{-or}, \emph{-ón}, \emph{-ín} & \esp{trabajador} / \esp{bonit\textbf{ón}} & \esp{trabajad\textbf{ora}} / \esp{bonit\textbf{ona}} & \esp{trabajad\textbf{ores}} / \esp{bonit\textbf{ones}} \\ \hline
Invariables (\emph{-ista}, \emph{-e} nationalité) & \esp{turista} / \esp{vasco} & \esp{turista} / \esp{vasca} & \esp{turistas} / \esp{vascos} \\ \hline
\end{tabular}
\end{center}

\begin{exemplebox}[Exemples d'accord et de place]
\esp{Un paisaje \textbf{impresionante}} (masc. sg.) -- \esp{Una playa \textbf{impresionante}} (fém. sg.)\par
\esp{Unos paisajes \textbf{impresionantes}} -- \esp{Unas playas \textbf{impresionantes}}\par
\esp{Un hotel \textbf{cómodo y barato}} (deux adjectifs après le nom)\par
\emph{Exception (place avant le nom) :} adjectifs subjectifs, poétiques ou changant de sens : \esp{un \textbf{gran} hotel} (un grand hôtel = célèbre) vs \esp{un hotel \textbf{grande}} (un hôtel grand = de grande taille) ; \esp{un \textbf{buen} guía} (apocope de \esp{bueno}).
\end{exemplebox}

\courssub{Méthode du commentaire et commentaire modèle}

\begin{methode}[Commenter un texte de presse ou d'actualité (Écrit / Oral)]
\begin{enumerate}
\item \textbf{Lecture active} : Lis deux fois. Souligne les mots-clés (thème, chiffres, noms propres, connecteurs logiques : \esp{por eso, sin embargo, además}).
\item \textbf{Identification} : Nature du document (article, rapport, blog), source, date, auteur, thème général.
\item \textbf{Structure} : Repère les paragraphes = les grandes idées (Introduction, Arguments, Conclusion).
\item \textbf{Synthèse} : Résume l'idée principale en 2 phrases (français ou espagnol selon consigne).
\item \textbf{Analyse lexicale} : Releve 5--7 mots du champ lexical étudié (tourisme, développement, nature). Explique-en un ou deux dans le contexte.
\item \textbf{Point de vue / Réaction personnelle} : Donne ton avis argumenté : \esp{En mi opinión...}, \esp{Estoy de acuerdo porque...}, \esp{Me parece que...}.
\item \textbf{Rédaction finale} : Structure : Introduction (présentation doc + thèse) -- Développement (2--3 arguments appuyés sur le texte) -- Conclusion (ouverture / avis personnel).
\end{enumerate}
\end{methode}

\begin{exemplebox}[Commentaire modèle (niveau A2+/B1) -- Sur le texte support « El turismo, motor de desarrollo »]
\textbf{Introducción :} El texto « El turismo, motor de desarrollo » presenta el turismo como un sector económico fundamental para países como España o México, y destaca el potencial inexplotado de Camerún. El autor defiende la idea de un turismo sostenible.\par
\textbf{Desarrollo :} En primer lugar, el texto da datos económicos: el turismo supera el 12\,\% del PIB en España. En segundo lugar, enumera los atractivos de Camerún: playas (Kribi), montaña (Monte Camerún), fauna (Waza) y cultura (Foumban). En tercer lugar, propone soluciones: el ecoturismo y el turismo comunitario para proteger el medio ambiente (\esp{proteger el medio ambiente}) y beneficiar a la gente local (\esp{beneficiar a la población local}).\par
\textbf{Conclusión :} En mi opinión, Camerún tiene todo para ser una potencia turística africana. \esp{Hay que} invertir en formación de guías e infraestructuras, pero \esp{hay que} hacerlo con respeto a la naturaleza. \esp{¡El turismo sí, pero sostenible!}
\end{exemplebox}

\courssub{Production guidée : présenter et promouvoir un site}

Dans la section précédente, nous avons appris à \cle{commenter un texte} : lire, repérer l'idée principale, relever le lexique et rédiger un commentaire structuré. Nous allons maintenant passer de la lecture à la \cle{production} : c'est à ton tour de parler et d'écrire ! L'objectif de cette dernière étape est de savoir \textbf{présenter un site touristique} (le mont Cameroun, Kribi, le parc de Waza, une plage d'Espagne...) et de \textbf{promouvoir le tourisme durable}. C'est une compétence très concrète : un guide touristique, un animateur d'office de tourisme ou un élève qui rédige un article pour le journal du lycée utilisent exactement cette démarche.

\courssub{Le plan de la présentation : cinq étapes en escalier}

Pour présenter un lieu, on ne dit pas les choses dans le désordre. On suit un chemin logique, comme un escalier qu'on monte marche par marche :

\begin{enumerate}
\item \textbf{Nombre y localización} : on dit d'abord le nom du lieu et où il se trouve, avec \esp{estar} (la position dans l'espace).
\item \textbf{Descripción} : on décrit le lieu, sa nature, son caractère, avec \esp{ser} + adjectifs.
\item \textbf{Qué hay} : on énumère les richesses du site avec \esp{hay} (il y a).
\item \textbf{Actividades} : on dit ce que les touristes peuvent faire (\esp{los turistas pueden} + infinitif).
\item \textbf{Promoción del turismo sostenible} : on termine par un appel à visiter et à protéger le lieu (\esp{hay que} + infinitif).
\end{enumerate}

\begin{popfigure}
\begin{tikzpicture}[
etape/.style={rectangle, rounded corners=4pt, draw=popblue, fill=popblueL, text width=3.8cm, align=center, minimum height=1.1cm, font=\small, inner sep=3pt},
num/.style={circle, draw=popdark, fill=poporange, text=white, font=\bfseries\small, inner sep=2pt, minimum width=0.6cm}
]
\node[etape, align=center] (e1){1. \textbf{Nombre y localización}\\\emph{estar + lugar}};
\node[num] at (-2.2,0) {1};
\node[etape, align=center] (e2) at (0, -1.6){2. \textbf{Descripción}\\\emph{ser + adjetivos}};
\node[num] at (-2.2,-1.6) {2};
\node[etape, align=center] (e3) at (0, -3.2){3. \textbf{Qué hay}\\\emph{hay + sustantivos}};
\node[num] at (-2.2,-3.2) {3};
\node[etape, align=center] (e4) at (0, -4.8){4. \textbf{Actividades}\\\emph{poder + infinitivo}};
\node[num] at (-2.2,-4.8) {4};
\node[etape, align=center] (e5) at (0, -6.4){5. \textbf{Turismo sostenible}\\\emph{hay que + infinitivo}};
\node[num] at (-2.2,-6.4) {5};
\draw[-{Stealth[length=3mm]}, thick, popdark] (e1.south) -- (e2.north);
\draw[-{Stealth[length=3mm]}, thick, popdark] (e2.south) -- (e3.north);
\draw[-{Stealth[length=3mm]}, thick, popdark] (e3.south) -- (e4.north);
\draw[-{Stealth[length=3mm]}, thick, popdark] (e4.south) -- (e5.north);
\end{tikzpicture}
\legende{Les cinq étapes de la présentation d'un site touristique : on monte l'escalier marche par marche, de la localisation jusqu'à l'appel à protéger.}
\end{popfigure}

\begin{methode}[Présenter un lieu touristique en cinq temps]
\begin{enumerate}
\item \textbf{Commencer par la localisation} avec \esp{estar} : \esp{Kribi está en el sur de Camerún, en la costa atlántica.} (Kribi est dans le sud du Cameroun, sur la côte atlantique.)
\item \textbf{Continuer par les caractéristiques} avec \esp{ser} + adjectifs : \esp{Es una ciudad pequeña y tranquila.} (C'est une ville petite et tranquille.)
\item \textbf{Énumérer les richesses} avec \esp{hay} : \esp{Hay playas magníficas y una cascada famosa.} (Il y a des plages magnifiques et une cascade célèbre.)
\item \textbf{Dire ce qu'on peut faire} : \esp{Los turistas pueden nadar, pescar y comer pescado fresco.} (Les touristes peuvent nager, pêcher et manger du poisson frais.)
\item \textbf{Finir par un appel à visiter et à protéger} : \esp{¡Ven a Kribi! Hay que proteger sus playas.} (Viens à Kribi ! Il faut protéger ses plages.)
\end{enumerate}
\end{methode}

\courssub{La trame guidée : des phrases toutes faites pour parler}

Pour t'aider, voici une \cle{trame} (un squelette de texte) avec des trous à compléter. Chaque ligne correspond à une marche de l'escalier :

\begin{center}
\begin{tabularx}{\linewidth}{|l|X|l|}\hline
\rowcolor{popblueL} \textbf{Étape} & \textbf{Trame espagnole} & \textbf{Français} \\ \hline
1 & \esp{... está en ...} & ... est (se trouve) à/dans ... \\ \hline
2 & \esp{Es un lugar ...} & C'est un lieu ... \\ \hline
3 & \esp{Hay ...} & Il y a ... \\ \hline
4 & \esp{Los turistas pueden ...} & Les touristes peuvent ... \\ \hline
5 & \esp{Hay que proteger ...} & Il faut protéger ... \\ \hline
\end{tabularx}
\end{center}

Observons maintenant un exemple complet, appliqué au mont Cameroun :

\begin{exemplebox}[Un modèle de présentation : le mont Cameroun]
\esp{El monte Camerún está en el suroeste del país. Es un volcán impresionante. Hay senderos para los turistas y una naturaleza única. Los turistas pueden subir a la cumbre y observar animales raros. Hay que respetar el medio ambiente.}\par
\medskip
« Le mont Cameroun est dans le sud-ouest du pays. C'est un volcan impressionnant. Il y a des sentiers pour les touristes et une nature unique. Les touristes peuvent monter au sommet et observer des animaux rares. Il faut respecter l'environnement. »
\end{exemplebox}

Remarque comme chaque phrase utilise un outil de grammaire déjà étudié : \esp{está} (position), \esp{es} (caractéristique permanente), \esp{hay} (existence), \esp{pueden} + infinitif (possibilité), \esp{hay que} + infinitif (obligation générale).

\begin{attention}[« Il faut » = \esp{hay que} + infinitif]
En français, « il faut » est très fréquent. En espagnol, on le traduit par \esp{hay que} suivi d'un \textbf{infinitif}, jamais par une forme de \esp{ser} !\par
\esp{Hay que conservar las reservas naturales.} = Il faut conserver les réserves naturelles.\par
\esp{Hay que respetar la cultura de los pueblos.} = Il faut respecter la culture des peuples.\par
\textbf{Erreur typique du francophone} : écrire \esp{es necesario que protege} (subjonctif mal maîtrisé) ou calquer « il faut » mot à mot. Retiens : \esp{hay que} est invariable, et le verbe qui suit reste à l'infinitif.
\end{attention}

\courssub{Promouvoir le tourisme durable : les arguments}

Présenter un site ne suffit pas : un bon promoteur sait aussi \textbf{défendre le tourisme durable} (\esp{el turismo sostenible}), c'est-à-dire un tourisme qui enrichit sans détruire. Voici les trois grands arguments, avec leur vocabulaire :

\begin{center}
\begin{tabularx}{\linewidth}{|X|X|}\hline
\rowcolor{popgreenL} \textbf{Argumento en español} & \textbf{Traduction française} \\ \hline
\esp{proteger la naturaleza y el paisaje} & protéger la nature et le paysage \\ \hline
\esp{respetar la cultura y las tradiciones} & respecter la culture et les traditions \\ \hline
\esp{crear empleos para la gente local} & créer des emplois pour la population locale \\ \hline
\esp{conservar las reservas naturales} & conserver les réserves naturelles \\ \hline
\esp{no tirar basura en las playas} & ne pas jeter de déchets sur les plages \\ \hline
\esp{comprar productos locales} & acheter des produits locaux \\ \hline
\end{tabularx}
\end{center}

Un \cle{slogan} est une phrase courte, rythmée et facile à retenir. En voici deux modèles :

\begin{exemplebox}[Slogans pour promouvoir le Cameroun]
\esp{¡Visita Camerún, África en miniatura!} = Visite le Cameroun, l'Afrique en miniature !\par
\esp{Turismo sí, pero sostenible.} = Tourisme oui, mais durable.\par
\esp{Kribi te espera: sol, mar y naturaleza.} = Kribi t'attend : soleil, mer et nature.
\end{exemplebox}

\begin{savaistu}[África en miniatura]
Le Cameroun est surnommé \esp{África en miniatura} (« l'Afrique en miniature ») parce qu'il réunit sur un seul territoire presque tous les paysages du continent : les plages de Kribi et Limbé, les montagnes de l'Ouest, la savane et les éléphants du parc de Waza, la forêt dense du Sud et même un volcan actif, le mont Cameroun. C'est un argument idéal pour un slogan touristique : visiter un seul pays, c'est découvrir toute l'Afrique !
\end{savaistu}

\courssub{Un texte modèle complet}

Lisons maintenant une présentation rédigée par un élève de Première. Observe comment il suit l'escalier des cinq étapes.

\begin{textebox}[El parque de Waza, un tesoro del norte]
El parque nacional de Waza está en el extremo norte de Camerún, cerca de la frontera con Nigeria y el Chad. Es un lugar extraordinario y muy famoso en toda África.\par
En el parque hay una sabana inmensa, muchos animales salvajes y paisajes maravillosos. Hay elefantes, leones, jirafas, antílopes y cientos de pájaros de colores. También hay guías cameruneses que conocen muy bien el parque.\par
Los turistas pueden hacer un safari en coche, fotografiar a los animales y dormir cerca del parque. Pueden también hablar con la gente local y descubrir su cultura.\par
El turismo crea empleos para los habitantes de la región: guías, cocineros, artesanos. Pero hay que proteger este tesoro. Hay que respetar a los animales, no tirar basura y conservar la reserva natural para los niños del futuro.\par
¡Visita Waza, el corazón salvaje de Camerún!
\end{textebox}

\textbf{Glossaire} : \esp{la sabana} = la savane ; \esp{salvaje} = sauvage ; \esp{la jirafa} = la girafe ; \esp{el pájaro} = l'oiseau ; \esp{hacer un safari} = faire un safari ; \esp{la gente local} = la population locale ; \esp{el artesano} = l'artisan ; \esp{tirar basura} = jeter des déchets ; \esp{el corazón} = le cœur.

\textbf{Questions de compréhension} :
\begin{enumerate}
\item ¿Dónde está el parque de Waza? (Où se trouve le parc de Waza ?)
\item ¿Qué animales hay en el parque? (Quels animaux y a-t-il dans le parc ?)
\item ¿Qué pueden hacer los turistas? (Que peuvent faire les touristes ?)
\item ¿Qué hay que proteger? (Que faut-il protéger ?)
\end{enumerate}

\courssub{Version guidée : du français vers l'espagnol}

\begin{exoresolu}[Version : un site camerounais]
\textbf{Énoncé.} Traduis en espagnol les quatre phrases suivantes :\par
1. Les chutes de la Lobé sont près de Kribi.\par
2. C'est un lieu magnifique et unique.\par
3. Il y a une plage et des restaurants de poisson.\par
4. Il faut protéger ce paysage extraordinaire.
\tcblower
\textbf{Solution détaillée.}\par
1. « sont près de » exprime une position : verbe \esp{estar}. \esp{Las cataratas de la Lobé están cerca de Kribi.}\par
2. « c'est » exprime une caractéristique : verbe \esp{ser}. Attention à l'accord des adjectifs avec \esp{un lugar} (masculin singulier) : \esp{Es un lugar magnífico y único.}\par
3. « il y a » = \esp{hay}, invariable : \esp{Hay una playa y restaurantes de pescado.}\par
4. « il faut » = \esp{hay que} + infinitif : \esp{Hay que proteger este paisaje extraordinario.}
\end{exoresolu}

\courssub{Thème guidé : de l'espagnol vers le français}

\begin{exoresolu}[Thème : phrases tirées du texte support]
\textbf{Énoncé.} Traduis en français les trois phrases suivantes :\par
1. \esp{El turismo es un motor de desarrollo económico para muchos países.}\par
2. \esp{En Camerún hay playas, montañas y reservas naturales.}\par
3. \esp{Hay que promover un turismo sostenible que respete la cultura local.}
\tcblower
\textbf{Solution détaillée.}\par
1. Le tourisme est un moteur de développement économique pour beaucoup de pays. (\esp{motor} = moteur, masculin malgré son aspect.)\par
2. Au Cameroun, il y a des plages, des montagnes et des réserves naturelles. (\esp{hay} = il y a ; \esp{montañas} avec le tilde sur le \~n.)\par
3. Il faut promouvoir un tourisme durable qui respecte la culture locale. (\esp{hay que} + infinitif = il faut ; \esp{respete} est un subjonctif présent après un antécédent indéfini \esp{un turismo... que}, on le traduit simplement par l'indicatif français.)
\end{exoresolu}

\begin{attention}[Pièges de traduction à éviter]
\begin{itemize}
\item Ne traduis jamais « il y a » par \esp{es} ou \esp{está} : c'est toujours \esp{hay}.
\item \esp{único} signifie « unique » mais aussi « seul » : \esp{una naturaleza única} = une nature unique (pas « une nature seule »).
\item Attention à l'ordre des adjectifs : en espagnol, l'adjectif qualificatif se place en général \textbf{après} le nom : \esp{un volcán impresionante}, \esp{playas magníficas}.
\item N'oublie pas les accents : \esp{único}, \esp{magnífico}, \esp{económico}, \esp{región}, \esp{también}, \esp{país}.
\end{itemize}
\end{attention}

\courssub{À toi de produire : activité guidée}

\begin{experience}[Soyons guides touristiques !]
Par groupes de quatre, choisissez un site du Cameroun (Kribi, le mont Cameroun, Waza, les chutes de la Lobé, Foumban et son palais royal...) ou d'un pays hispanophone (les plages d'Espagne, le Machu Picchu, Malabo en Guinée équatoriale, Carthagène des Indes en Colombie).\par
\textbf{Étape 1} : rédigez une présentation de 8 à 10 phrases en suivant l'escalier des cinq étapes et la trame guidée.\par
\textbf{Étape 2} : inventez un slogan avec un point d'exclamation : \esp{¡Visita ...!}\par
\textbf{Étape 3} : un membre du groupe joue le rôle du guide et présente le site à la classe ; les autres groupes posent deux questions en espagnol : \esp{¿Qué hay en ...? ¿Qué pueden hacer los turistas?}
\end{experience}

\courssub{Grille d'évaluation}

Ta production sera évaluée sur 10 points selon la grille suivante. Lis-la bien avant de rédiger : elle te dit exactement ce que le correcteur attend.

\begin{center}
\begin{tabularx}{\linewidth}{|X|c|X|}\hline
\rowcolor{popgoldL} \textbf{Critère} & \textbf{Points} & \textbf{Ce qui est attendu} \\ \hline
Lexique du tourisme & 3 pts & au moins 6 mots du champ lexical : \esp{turismo, hotel, paisaje, patrimonio, reserva natural, guía, ecoturismo, playa, montaña...} \\ \hline
Emploi de hay / ser / estar & 3 pts & \esp{estar} pour la localisation, \esp{ser} pour les caractéristiques, \esp{hay} pour l'existence (distinction claire) \\ \hline
Adjectifs & 2 pts & adjectifs corrects, placés après le nom et accordés en genre/nombre \\ \hline
Cohérence & 2 pts & les 5 étapes respectées, un slogan final, un appel à protéger (\esp{hay que} + inf.) \\ \hline
\end{tabularx}
\end{center}

\begin{aretenir}[L'essentiel de la production guidée]
Pour présenter et promouvoir un site touristique, monte les cinq marches de l'escalier : \textbf{localiser} (\esp{está en}), \textbf{décrire} (\esp{es un lugar} + adjectifs), \textbf{énumérer} (\esp{hay}), \textbf{proposer des activités} (\esp{los turistas pueden} + infinitif), et \textbf{conclure par un appel} (\esp{hay que proteger} + un slogan). « Il faut » se dit toujours \esp{hay que} + infinitif. Avec ce plan, tu peux présenter n'importe quel lieu du Cameroun ou du monde hispanophone, à l'écrit comme à l'oral.
\end{aretenir}

\newpage

\coursec{Activité d'intégration}

\coursec{El turismo y el desarrollo económico}

\courssub{Objectifs de l'activité}

À l'issue de cette activité d'intégration, l'élève sera capable de :
\begin{itemize}
\item comprendre et exploiter un modèle de texte promotionnel touristique en espagnol ;
\item présenter oralement et par écrit un site touristique camerounais en mobilisant le lexique du tourisme (\esp{turismo, turista, hotel, paisaje, patrimonio, reserva natural, ecoturismo, guía}) ;
\item décrire un lieu en utilisant correctement \esp{hay}, \esp{ser} et \esp{estar} ainsi qu'au moins trois adjectifs bien accordés ;
\item promouvoir le tourisme durable à l'aide de la structure impersonnelle \esp{hay que + infinitivo} ;
\item travailler en groupe, présenter une production en deux minutes et répondre aux questions d'un public.
\end{itemize}

\courssub{Situation}

L'Office du tourisme de ta région organise un concours scolaire intitulé \esp{« Descubre Camerún »}. Une délégation de touristes espagnols visitera bientôt le lycée : ils veulent découvrir les richesses touristiques du Cameroun, mais ils ne parlent pas français. Ton groupe a été choisi pour créer un dépliant touristique (\esp{folleto turístico}) en espagnol sur l'un de ces trois sites : \textbf{Kribi} (les plages et les chutes de la Lobé), \textbf{le mont Cameroun} (le volcan et la randonnée) ou \textbf{le parc national de Waza} (la faune et la savane). Le meilleur dépliant sera envoyé à l'agence de voyage partenaire à Madrid.

Voici le modèle de dépliant publié par l'Office du tourisme pour inspirer les participants :

\begin{textebox}[Modelo de folleto — Oficina de Turismo]
\textbf{DESCUBRE KRIBI, LA PERLA DEL ATLÁNTICO}

Kribi está en el sur de Camerún, a orillas del océano Atlántico. Es una ciudad pequeña pero muy acogedora.

En Kribi hay playas magníficas de arena blanca, hoteles cómodos y restaurantes de pescado fresco. También hay una maravilla natural única: las cataratas del río Lobé, que caen directamente en el mar.

¿Qué puedes hacer? Puedes nadar, pasear en piragua con un guía local y visitar los pueblos de pescadores.

¡Atención, turista! Hay que respetar la naturaleza: no tires basura en la playa y protege las tortugas marinas.

\textbf{Kribi te espera: ¡el paraíso está aquí!}
\end{textebox}

\emph{Glossaire :} \esp{a orillas de} = au bord de ; \esp{acogedora} = accueillante ; \esp{cataratas} = chutes ; \esp{piragua} = pirogue ; \esp{basura} = ordures ; \esp{tortugas marinas} = tortues marines.

\courssub{Consignes}

Par groupes de 3 ou 4 élèves, réalisez les tâches suivantes dans l'ordre :

\begin{enumerate}
\item \textbf{Compréhension.} Lisez le modèle de dépliant ci-dessus. Soulignez en rouge toutes les formes de \esp{hay}, en bleu les formes de \esp{ser} et en vert les formes de \esp{estar}. Relevez aussi trois adjectifs qualificatifs.
\item \textbf{Choix du site.} Choisissez un site : Kribi, le mont Cameroun ou le parc de Waza. Complétez une fiche d'identité du lieu : nom, région, richesses naturelles ou culturelles, activités possibles.
\item \textbf{Localisation.} Rédigez une ou deux phrases avec \esp{estar} pour situer le site : \esp{El monte Camerún está en el suroeste, cerca de Buea.}
\item \textbf{Description.} Rédigez deux ou trois phrases avec \esp{ser} et au moins \textbf{trois adjectifs différents} bien accordés en genre et en nombre : \esp{Es un volcán impresionante y peligroso.}
\item \textbf{Richesses.} Rédigez deux phrases avec \esp{hay} pour présenter ce qu'on trouve sur place : paysages, animaux, hôtels, monuments, guides...
\item \textbf{Activités.} Proposez \textbf{deux activités} pour les touristes avec \esp{se puede + infinitivo} ou \esp{puedes + infinitivo} : randonner, observer les éléphants, se baigner, goûter la cuisine locale...
\item \textbf{Tourisme durable.} Rédigez un message écologique avec \textbf{deux phrases} en \esp{hay que + infinitivo} : \esp{Hay que proteger a los animales. Hay que respetar a la población local.}
\item \textbf{Slogan.} Inventez un slogan court, rythmé et positif en espagnol : \esp{¡Waza, la aventura salvaje!}
\item \textbf{Mise en page.} Disposez tous les éléments sur une feuille ou une affiche : titre, paragraphes, slogan, dessins ou symboles (un arbre, un éléphant, une vague).
\item \textbf{Présentation orale.} Chaque groupe présente son dépliant en \textbf{2 minutes} : chaque membre parle au moins 20 secondes. La classe joue le rôle des touristes espagnols et pose au moins deux questions par groupe (\esp{¿Hay hoteles cerca? ¿Qué animales hay? ¿Cómo está el clima?}).
\end{enumerate}

\begin{attention}[Critères d'évaluation]
Ton dépliant sera noté sur : la correction de \esp{hay/ser/estar} (4 points), l'accord des adjectifs (4 points), le lexique du tourisme (4 points), le message écologique et le slogan (4 points), la présentation orale et la prononciation (4 points).
\end{attention}

\courssub{Ressources à mobiliser}

\begin{aretenir}[Les trois outils de la description d'un lieu]
\begin{itemize}
\item \esp{hay + nom} : il y a (existence, invariable) — \esp{En Waza hay leones y jirafas.} « À Waza, il y a des lions et des girafes. »
\item \esp{ser + adjectif} : caractéristique permanente, identité, nature — \esp{El paisaje es hermoso.} « Le paysage est magnifique. »
\item \esp{estar + lieu / adjectif d'état} : localisation géographique ou état passager — \esp{Kribi está en el sur.} « Kribi est dans le Sud. » / \esp{La playa está limpia.} « La plage est propre (en ce moment). »
\end{itemize}
\end{aretenir}

\begin{propriete}[Structures impersonnelles de promotion touristique]
\begin{itemize}
\item \esp{hay que + infinitivo} : « il faut + infinitif » (obligation générale, impersonnel, pas de sujet). \esp{Hay que respetar la naturaleza.}
\item \esp{se puede + infinitivo} : « on peut + infinitif » (possibilité générale, voix passive réflexe impersonnelle). \esp{Se puede nadar en el río.}
\item \esp{puedes + infinitivo} : « tu peux / vous pouvez + infinitif » (adresse directe au touriste, 2e personne). \esp{Puedes visitar el parque.}
\end{itemize}
\end{propriete}

\begin{definition}[Lexique thématique du tourisme et adjectifs utiles]
\begin{itemize}
\item \textbf{Noms :} \esp{el turismo} (tourisme), \esp{el turista} (touriste), \esp{el hotel} (hôtel), \esp{el paisaje} (paysage), \esp{el patrimonio} (patrimoine), \esp{la reserva natural} (réserve naturelle), \esp{el ecoturismo} (écotourisme), \esp{el/la guía} (guide), \esp{la playa} (plage), \esp{la montaña} (montagne), \esp{la fauna} (faune), \esp{la selva} (forêt tropicale), \esp{el volcán} (volcan), \esp{la sabana} (savane), \esp{la catarata} (chute d'eau).
\item \textbf{Adjectifs (accord genre/nombre) :} \esp{hermoso/a} (magnifique), \esp{impresionante} (impressionnant), \esp{salvaje} (sauvage), \esp{acogedor/acogedora} (accueillant), \esp{único/a} (unique), \esp{famoso/a} (célèbre), \esp{tranquilo/a} (tranquille), \esp{caluroso/a} (chaud), \esp{peligroso/a} (dangereux), \esp{verde} (vert), \esp{azul} (bleu).
\item \textbf{Verbes d'action :} \esp{nadar} (nager), \esp{pasear} (se promener), \esp{observar} (observer), \esp{fotografiar} (photographier), \esp{proteger} (protéger), \esp{respetar} (respecter), \esp{descubrir} (découvrir), \esp{disfrutar} (profiter).
\item \textbf{Prononciation :} \esp{guía} [« gui-a »] (hiatus, deux syllabes, accent sur le \emph{i}) ; \esp{paisaje} [« pa-ï-sa-ré »] (diphthongue \emph{ai} + hiatus \emph{je}) ; \esp{catarata} [« ca-ta-ra-ta »] ; \esp{ecoturismo} [« é-co-tou-ris-mo »].
\end{itemize}
\end{definition}

\courssub{Méthode : Rédiger un dépliant touristique}

\begin{methode}[Étapes de la rédaction d'un \esp{folleto turístico}]
\begin{enumerate}
\item \textbf{Analyser le modèle} : repérer la structure (titre, localisation, description, richesses, activités, message écologique, slogan) et les outils grammaticaux (\esp{estar, ser, hay, hay que, se puede}).
\item \textbf{Choisir et documenter le site} : remplir une fiche d'identité (nom, région, atouts, activités).
\item \textbf{Rédiger la localisation} : \esp{estar} + complément de lieu (prépositions : \esp{en, cerca de, al norte de, a orillas de}).
\item \textbf{Rédiger la description} : \esp{ser} + adjectifs accordés (qualité, nature, taille, beauté). Varier les adjectifs.
\item \textbf{Lister les richesses} : \esp{hay} + noms (paysages, faune, flore, infrastructures, guides). \esp{Hay} est invariable.
\item \textbf{Proposer des activités} : \esp{se puede / puedes / podéis} + infinitif. Utiliser des verbes variés.
\item \textbf{Rédiger l'engagement durable} : \esp{hay que} + infinitif (obligation morale, protection environnementale et culturelle). Minimum deux phrases.
\item \textbf{Créer le slogan} : court, impératif (\esp{ven, descubre, visita, vive, disfruta}) + image forte.
\item \textbf{Soigner la mise en page} : titre visible, paragraphes aérés, illustrations, orthographe vérifiée (accents, \esp{ñ}, \esp{¿ ¡}).
\item \textbf{Répéter l'oral} : répartir le texte, articuler, regarder le public, anticiper les questions (\esp{¿Hay...? ¿Dónde está...? ¿Es...?}).
\end{enumerate}
\end{methode}

\courssub{Proposition de production et corrigé commenté}

\begin{exoresolu}[Folleto modelo : « Descubre el Parque Nacional de Waza »]
Rédige le dépliant complet d'un groupe qui a choisi le parc de Waza, en respectant toutes les consignes.
\tcblower
\textbf{Production modèle :}

\esp{\textbf{DESCUBRE WAZA, EL CORAZÓN SALVAJE DE CAMERÚN}}

\esp{El Parque Nacional de Waza está en el extremo norte de Camerún, cerca de la frontera con Nigeria y Chad. Es un parque muy famoso y un gran patrimonio natural del país.}

\esp{En Waza hay muchos animales salvajes: elefantes, leones, jirafas, antílopes y aves de colores. También hay un paisaje impresionante de sabana y guías muy amables que conocen bien el parque.}

\esp{¿Qué puedes hacer? Puedes observar a los elefantes al amanecer y hacer fotos magníficas con un guía local. También se puede visitar los pueblos de la región y descubrir la cultura del norte.}

\esp{Turista responsable: hay que respetar a los animales y no darles comida. Hay que proteger la naturaleza y no dejar basura en el parque. El ecoturismo es el futuro.}

\esp{\textbf{Waza te espera: ¡ven y vive la aventura salvaje!}}

\medskip
\textbf{Commentaire des choix de langue :}
\begin{itemize}
\item \esp{está en el extremo norte} : \esp{estar} pour la localisation géographique — correct ; on n'utilise jamais \esp{ser} pour situer un lieu sur la carte.
\item \esp{Es un parque muy famoso} : \esp{ser} pour la caractéristique permanente (identité, renommée) ; adjectif \esp{famoso} accordé au masculin singulier avec \esp{parque}.
\item \esp{hay muchos animales salvajes} : \esp{hay} invariable devant un pluriel ; \esp{salvajes} accordé au pluriel (adjectif en \emph{-e} : même forme masc./fém.).
\item \esp{guías muy amables} : \esp{guía} est épicène (\esp{el guía / la guía}) ; \esp{amables} invariable en genre, prend \emph{-s} au pluriel.
\item \esp{Puedes observar / se puede visitar} : deux variantes correctes pour proposer des activités (adresse directe vs impersonnel).
\item \esp{al amanecer} : « au lever du soleil » ; \esp{al} = \esp{a + el} devant le nom masculin \esp{amanecer} (pas un infinitif ici).
\item \esp{hay que respetar / hay que proteger} : structure impersonnelle d'obligation générale ; le verbe reste à l'infinitif (\esp{*hay que respetas} est une faute grave).
\item \esp{no darles comida} : pronom \esp{les} (COI, « à eux ») rattaché à l'infinitif après la négation — « ne leur donne pas à manger ».
\item \esp{no dejar basura} : infinitif après \esp{hay que} ; \esp{basura} (ordures) est singulier collectif.
\item Slogan final : impératifs \esp{ven} (verbe \esp{venir}, 2e pers. sg.) + \esp{vive} (verbe \esp{vivir}, 2e pers. sg.) ; vocabulaire positif (\esp{aventura salvaje}) ; court, rythmé, mémorable.
\end{itemize}
\end{exoresolu}

\begin{experience}[Variantes et prolongements]
\begin{itemize}
\item \textbf{Jeu de rôle :} un élève joue le guide touristique, un autre le touriste espagnol exigeant qui pose des questions difficiles (\esp{¿Es peligroso el volcán? ¿Hay agua potable en Waza? ¿Cuánto cuesta la entrada?}).
\item \textbf{Concours de slogans :} la classe vote pour le meilleur slogan ; les slogans gagnants sont affichés au mur de la salle avec une illustration.
\item \textbf{Ouverture hispanophone :} comparer avec un site d'un pays hispanophone en deux phrases (ex. \esp{El Parque de Doñana, en España, está en Andalucía. Es una reserva natural única con flamencos y linces.}) ; citer la Guinée équatoriale : \esp{Malabo, en Guinea Ecuatorial, está en la isla de Bioko. Es una ciudad acogedora con playas vírgenes.}
\item \textbf{Version numérique :} réaliser le dépliant sur Canva ou PowerPoint avec photos libres de droits, QR code vers une vidéo de présentation orale.
\end{itemize}
\end{experience}

\courssub{Bilan de l'activité}

\begin{aretenir}[Ce que j'ai appris à faire]
\begin{itemize}
\item Je sais \textbf{lire et exploiter} un texte promotionnel touristique en espagnol.
\item Je sais \textbf{localiser} un lieu avec \esp{estar}, le \textbf{décrire} avec \esp{ser} + adjectifs accordés, et présenter ses richesses avec \esp{hay}.
\item Je sais \textbf{proposer des activités} avec \esp{se puede / puedes + infinitivo}.
\item Je sais \textbf{promouvoir le tourisme durable} avec \esp{hay que + infinitivo} et rédiger un slogan percutant à l'impératif.
\item Je sais \textbf{présenter un projet à l'oral} en deux minutes et répondre aux questions d'un public hispanophone.
\end{itemize}
\end{aretenir}

Cette activité t'a permis de réinvestir l'ensemble des acquis de la leçon dans une tâche concrète et authentique : informer de vrais visiteurs. Le tourisme durable n'est pas seulement un sujet d'examen : c'est un enjeu réel pour le Cameroun, dont les paysages, la faune et le patrimoine culturel constituent une richesse à protéger et à faire connaître. En espagnol, tu es désormais capable d'en être l'ambassadeur : \esp{¡El turismo responsable empieza contigo!} « Le tourisme responsable commence avec toi ! »

\newpage

\coursec{Méthodes et exercices résolus}

\courssub{Méthode 1 : Lire et comprendre un texte sur le tourisme}

\begin{methode}[Comprendre un texte touristique en 5 étapes]
\begin{enumerate}
\item \textbf{Lire le titre et la source.} Le titre (\esp{El turismo, motor de desarrollo}) annonce le thème et souvent la thèse de l'auteur. Souligne les mots-clés : \esp{turismo}, \esp{desarrollo}, \esp{empleo}, \esp{medio ambiente}.
\item \textbf{Repérer la structure.} Un texte journalistique suit souvent le plan : présentation du phénomène $\rightarrow$ avantages économiques $\rightarrow$ dangers $\rightarrow$ solutions (tourisme durable).
\item \textbf{Exploiter le lexique connu.} Les mots transparents aident : \esp{patrimonio} = patrimoine, \esp{reserva natural} = réserve naturelle, \esp{ecoturismo} = écotourisme. Méfie-toi des faux amis : \esp{actual} = actuel (pas « actuel » au sens de « réel »), \esp{asistir a} = assister à (pas « aider »).
\item \textbf{Répondre aux questions en reformulant.} Ne recopie jamais la phrase du texte telle quelle : reprends l'idée avec tes mots, en phrases complètes, avec sujet + verbe conjugué.
\item \textbf{Vérifier la cohérence.} Relis ta réponse : est-elle en espagnol correct ? Répond-elle vraiment à la question posée (\esp{¿Por qué?} demande une cause avec \esp{porque}) ?
\end{enumerate}
\end{methode}

\begin{exoresolu}[Compréhension de l'écrit — type Bac]
\textbf{Texte :} \esp{El turismo es una de las actividades económicas más importantes del mundo. En Camerún, lugares como Kribi, con sus playas magníficas, o el parque de Waza, con sus animales salvajes, atraen cada año a muchos visitantes. Sin embargo, el turismo puede dañar el medio ambiente si no se controla. Por eso, el ecoturismo propone viajar respetando la naturaleza y ayudando a la población local.}

\textbf{Questions :} 1) \esp{¿Qué lugares turísticos de Camerún menciona el texto?} 2) \esp{¿Qué peligro tiene el turismo?} 3) \esp{¿Qué propone el ecoturismo?}
\tcblower
\textbf{Réflexion :} Pour la question 1, je cherche les noms propres de lieux : \esp{Kribi} et \esp{el parque de Waza}. Pour la question 2, le mot \esp{peligro} me renvoie à \esp{dañar el medio ambiente}. Pour la question 3, je repère \esp{el ecoturismo propone...}.

\textbf{Réponses rédigées :}

1) \esp{El texto menciona Kribi, con sus playas, y el parque de Waza, con sus animales salvajes.}

2) \esp{El turismo puede dañar el medio ambiente si no se controla.}

3) \esp{El ecoturismo propone viajar respetando la naturaleza y ayudando a la población local.}

\textbf{Conclusion méthodologique :} chaque réponse est une phrase complète (sujet + verbe), qui reformule légèrement le texte sans le copier mot à mot, et qui répond exactement au mot interrogatif (\esp{¿qué?} $\rightarrow$ une chose).
\end{exoresolu}

\courssub{Méthode 2 : Présenter un site touristique}

\begin{methode}[Décrire un lieu avec hay, ser, estar]
\begin{enumerate}
\item \textbf{Situer le lieu} avec \esp{estar} + lieu : \esp{Kribi está en el sur de Camerún, cerca del mar.} (« Kribi se trouve dans le sud du Cameroun, près de la mer. »)
\item \textbf{Dire ce qu'il y a} avec \esp{hay} (invariable) : \esp{En Kribi hay playas hermosas, hoteles modernos y restaurantes de pescado.} (« Il y a de belles plages, des hôtels modernes et des restaurants de poisson. »)
\item \textbf{Caractériser} avec \esp{ser} + adjectif (caractéristique permanente) : \esp{El paisaje es impresionante y la gente es muy acogedora.} (« Le paysage est impressionnant et les gens sont très accueillants. »)
\item \textbf{Ajouter les activités} avec \esp{se puede / se pueden} + infinitif : \esp{Se puede nadar, pescar y visitar las cascadas de Lobé.} (« On peut nager, pêcher et visiter les chutes de la Lobé. »)
\item \textbf{Conclure par une invitation} : \esp{¡Venga a descubrir Kribi, no se arrepentirá!} (« Venez découvrir Kribi, vous ne le regretterez pas ! »)
\end{enumerate}
\textbf{Réflexe :} \esp{ser} = identité/caractéristique ; \esp{estar} = localisation et état temporaire ; \esp{hay} = existence (jamais d'accord : \esp{hay muchos hoteles}, jamais \esp{*han muchos hoteles}).
\end{methode}

\begin{exoresolu}[Rédaction guidée — présentation d'un site]
\textbf{Énoncé :} \esp{Redacta un texto de 8 a 10 líneas para presentar el monte Camerún a los turistas extranjeros. Utiliza hay, ser, estar y expresiones de promoción.}
\tcblower
\textbf{Plan de travail :} 1) localisation (\esp{estar}) ; 2) ce qu'on y trouve (\esp{hay}) ; 3) caractéristiques (\esp{ser}) ; 4) activités (\esp{se puede}) ; 5) invitation finale.

\textbf{Production modèle :}

\esp{El monte Camerún está en el suroeste del país, cerca de la ciudad de Buea. Es el volcán más alto de África occidental y su paisaje es realmente espectacular. En la región hay bosques tropicales, senderos para caminar y pueblos muy acogedores. Los guías locales son amables y conocen bien la montaña. Se puede hacer senderismo, observar aves raras y admirar la vista sobre el océano. Además, el clima es fresco y agradable. ¡Visite el monte Camerún: es una aventura inolvidable en el corazón de África!}

« Le mont Cameroun se trouve dans le sud-ouest du pays, près de la ville de Buéa. C'est le volcan le plus haut d'Afrique occidentale et son paysage est vraiment spectaculaire. Dans la région, il y a des forêts tropicales, des sentiers de randonnée et des villages très accueillants. Les guides locaux sont gentils et connaissent bien la montagne. On peut faire de la randonnée, observer des oiseaux rares et admirer la vue sur l'océan. De plus, le climat est frais et agréable. Visitez le mont Cameroun : c'est une aventure inoubliable au cœur de l'Afrique ! »

\textbf{Justifications grammaticales :} \esp{está} pour la localisation ; \esp{es} pour la caractéristique permanente (\esp{el volcán más alto}) ; \esp{hay} invariable devant un pluriel ; \esp{se puede} + infinitif ; superlatif \esp{el más alto de} (et non \esp{*el más alto en}).
\end{exoresolu}

\courssub{Méthode 3 : Promouvoir le tourisme durable}

\begin{methode}[Argumenter pour un tourisme responsable]
\begin{enumerate}
\item \textbf{Annoncer le problème} : \esp{El turismo masivo contamina las playas y destruye la naturaleza.} (« Le tourisme de masse pollue les plages et détruit la nature. »)
\item \textbf{Proposer des solutions} avec \esp{hay que} + infinitif (il faut), \esp{debemos} + infinitif (nous devons), \esp{es importante} + infinitif : \esp{Hay que proteger las reservas naturales.} / \esp{Debemos respetar la cultura local.}
\item \textbf{Mettre en valeur les bénéfices} avec connecteurs : \esp{además} (de plus), \esp{por eso} (c'est pourquoi), \esp{sin embargo} (cependant), \esp{gracias a} (grâce à).
\item \textbf{Appeler à l'action} à l'impératif : \esp{¡Viaje de forma responsable!} / \esp{¡Protejamos nuestro patrimonio!} (« Protégeons notre patrimoine ! »)
\end{enumerate}
\textbf{Structures à retenir :} \esp{hay que + inf.} (obligation générale), \esp{es necesario + inf.}, \esp{para + inf.} (but : \esp{para preservar el medio ambiente}).
\end{methode}

\begin{exoresolu}[Expression écrite — promotion du tourisme durable]
\textbf{Énoncé :} \esp{Escribe un pequeño artículo (10 líneas) para el periódico de tu colegio: « Por un turismo sostenible en Camerún ».}
\tcblower
\textbf{Démarche :} titre accrocheur $\rightarrow$ problème $\rightarrow$ solutions avec \esp{hay que / debemos} $\rightarrow$ bénéfices $\rightarrow$ appel final.

\textbf{Production modèle :}

\esp{Por un turismo sostenible en Camerún. Nuestro país tiene un patrimonio natural extraordinario: las playas de Kribi, el parque de Waza, el monte Camerún. Sin embargo, el turismo mal organizado puede destruir esta riqueza. Hay que limitar la basura en los sitios turísticos y debemos respetar a los animales salvajes. Es importante también que el dinero del turismo ayude a la población local, por ejemplo comprando productos artesanales. El ecoturismo crea empleos y protege la naturaleza al mismo tiempo. Por eso, ¡viajemos de manera responsable para preservar nuestro país para las generaciones futuras!}

« Pour un tourisme durable au Cameroun. Notre pays possède un patrimoine naturel extraordinaire : les plages de Kribi, le parc de Waza, le mont Cameroun. Cependant, le tourisme mal organisé peut détruire cette richesse. Il faut limiter les déchets dans les sites touristiques et nous devons respecter les animaux sauvages. Il est important aussi que l'argent du tourisme aide la population locale, par exemple en achetant des produits artisanaux. L'écotourisme crée des emplois et protège la nature en même temps. C'est pourquoi, voyageons de manière responsable pour préserver notre pays pour les générations futures ! »

\textbf{Points de langue vérifiés :} \esp{hay que} + infinitif ; \esp{debemos respetar} ; gérondif \esp{comprando} ; subjonctif après \esp{es importante que} : \esp{ayude} (et non \esp{*ayuda}) — structure B1 valorisée au Bac.
\end{exoresolu}

\courssub{Méthode 4 : Thème (français $\rightarrow$ espagnol)}

\begin{methode}[Traduire du français vers l'espagnol]
\begin{enumerate}
\item \textbf{Identifier le temps français} : présent $\rightarrow$ présent ; passé composé $\rightarrow$ \esp{pretérito perfecto} ou \esp{indefinido} ; imparfait $\rightarrow$ \esp{imperfecto}.
\item \textbf{Repérer les pièges} : « il y a » = \esp{hay} ; « se trouver » = \esp{estar} ; « c'est » (caractéristique) = \esp{es} ; « pour » + but = \esp{para} + infinitif.
\item \textbf{Conjuguer sans faute} : vérifier sujet, terminaison, accents (\esp{está}, \esp{visitará}, \esp{atracción}).
\item \textbf{Accorder les adjectifs} : genre et nombre avec le nom (\esp{las playas magníficas}).
\item \textbf{Relire} en chassant les calques : « visiter un lieu » = \esp{visitar un lugar} (jamais \esp{*visitar a} un lieu).
\end{enumerate}
\end{methode}

\begin{exoresolu}[Thème — phrases à traduire]
\textbf{Énoncé :} Traduire en espagnol : 1) « Le parc de Waza se trouve dans le nord du Cameroun. » 2) « Il y a beaucoup d'animaux sauvages dans la réserve. » 3) « Les touristes doivent respecter la nature pour protéger l'environnement. »
\tcblower
\textbf{Analyse :} 1) « se trouve » = localisation $\rightarrow$ \esp{estar}. 2) « il y a » $\rightarrow$ \esp{hay} invariable. 3) « doivent » = \esp{deben} ; « pour protéger » = but $\rightarrow$ \esp{para} + infinitif.

\textbf{Traductions :}

1) \esp{El parque de Waza está en el norte de Camerún.} — \esp{está} (localisation), accent obligatoire ; \esp{norte} sans article après \esp{en el}.

2) \esp{Hay muchos animales salvajes en la reserva.} — \esp{hay} ne s'accorde jamais ; \esp{salvajes} accordé au pluriel masculin.

3) \esp{Los turistas deben respetar la naturaleza para proteger el medio ambiente.} — \esp{deben} + infinitif ; \esp{para proteger} exprime le but ; « environnement » se dit \esp{medio ambiente} (faux ami : \esp{*ambiente} seul signifie « atmosphère, ambiance »).
\end{exoresolu}

\courssub{Méthode 5 : Version (espagnol $\rightarrow$ français)}

\begin{methode}[Traduire de l'espagnol vers le français]
\begin{enumerate}
\item \textbf{Lire la phrase entière} avant de traduire : repérer sujet, verbe, compléments.
\item \textbf{Identifier les temps espagnols} : \esp{atrae} = présent ; \esp{visitó} = passé simple $\rightarrow$ passé composé ou passé simple français ; \esp{era} = imparfait.
\item \textbf{Traduire les mots de liaison} : \esp{sin embargo} = cependant ; \esp{por eso} = c'est pourquoi ; \esp{además} = de plus.
\item \textbf{Écrire un français naturel} : ne pas calquer l'ordre espagnol mot à mot (\esp{un hotel grande} $\rightarrow$ « un grand hôtel »).
\item \textbf{Vérifier les faux amis} : \esp{actual} = actuel ; \esp{una larga estancia} = un long séjour ; \esp{realizar} = réaliser/effectuer.
\end{enumerate}
\end{methode}

\begin{exoresolu}[Version — texte à traduire]
\textbf{Énoncé :} Traduire en français : \esp{El ecoturismo atrae cada año a más viajeros que quieren descubrir la naturaleza sin destruirla. En Guinea Ecuatorial, los turistas pueden visitar playas vírgenes y bosques tropicales. Sin embargo, es necesario construir más hoteles y formar guías locales.}
\tcblower
\textbf{Analyse phrase par phrase :}

Phrase 1 : \esp{atrae} = présent « attire » ; \esp{cada año} = « chaque année » ; \esp{que quieren} = relatif « qui veulent » ; \esp{sin destruirla} = « sans la détruire » (\esp{la} = la nature, pronom féminin).

Phrase 2 : \esp{pueden visitar} = « peuvent visiter » ; \esp{playas vírgenes} = « plages vierges » (adjectif placé après le nom en espagnol, avant en français).

Phrase 3 : \esp{es necesario construir} = « il est nécessaire de construire » ; \esp{formar guías} = « former des guides » (faux ami : \esp{formar} = former, pas « formuler »).

\textbf{Traduction :} « L'écotourisme attire chaque année davantage de voyageurs qui veulent découvrir la nature sans la détruire. En Guinée équatoriale, les touristes peuvent visiter des plages vierges et des forêts tropicales. Cependant, il est nécessaire de construire plus d'hôtels et de former des guides locaux. »

\textbf{Conclusion :} la traduction respecte le sens, les temps et produit un français idiomatique (adjectifs replacés, tournures naturelles).
\end{exoresolu}

\begin{attention}[Les 5 erreurs qui coûtent des points au Bac]
\begin{enumerate}
\item \textbf{Oublier l'accent sur \esp{está}.} \esp{está} (il/elle est) sans accent devient \esp{esta} (ceci/cette). Écris toujours \esp{Kribi está en el sur}.
\item \textbf{Accorder \esp{hay}.} \esp{Hay} est invariable : \esp{hay muchos hoteles}, jamais \esp{*han muchos hoteles} ni \esp{*hay muchos hotel} (le nom, lui, s'accorde).
\item \textbf{Confondre \esp{ser} et \esp{estar}.} Localisation $\rightarrow$ \esp{estar} : \esp{El hotel está cerca de la playa}. Caractéristique $\rightarrow$ \esp{ser} : \esp{El hotel es grande}. Ne dis jamais \esp{*el hotel es en la playa}.
\item \textbf{Calquer le français.} « Visiter un lieu » = \esp{visitar un lugar} (sans \esp{a}) ; « environnement » = \esp{medio ambiente} ; « il faut » = \esp{hay que} + infinitif (pas \esp{*es necesario que hay}).
\item \textbf{Négliger l'accord des adjectifs.} \esp{las playas hermosas}, \esp{los paisajes impresionantes} : l'adjectif prend le genre et le nombre du nom, même placé après.
\end{enumerate}
\end{attention}

\newpage

\coursec{Exercices}

\courssub{Je vérifie mes connaissances}

\begin{exercice}[QCM : le tourisme et son vocabulaire]
Entoure la bonne réponse.
\begin{enumerate}
\item Un \esp{guía} es una persona que :
\begin{enumerate}
\item[a)] construye hoteles
\item[b)] acompaña y explica a los turistas
\item[c)] vende billetes de avión
\end{enumerate}
\item El \esp{ecoturismo} es un turismo que :
\begin{enumerate}
\item[a)] destruye la naturaleza
\item[b)] respeta el medio ambiente y las culturas locales
\item[c)] se practica solo en las ciudades
\end{enumerate}
\item Una \esp{reserva natural} es :
\begin{enumerate}
\item[a)] un espacio protegido donde viven animales y plantas
\item[b)] un restaurante de lujo
\item[c)] una playa privada
\end{enumerate}
\item El \esp{patrimonio} de un país es :
\begin{enumerate}
\item[a)] su dinero en el banco
\item[b)] el conjunto de sus riquezas culturales e históricas
\item[c)] su ejército
\end{enumerate}
\end{enumerate}
\end{exercice}

\begin{exercice}[Vrai ou faux ? Justifie]
Dis si chaque affirmation est vraie (V) ou fausse (F) et justifie ta réponse en une phrase en français.
\begin{enumerate}
\item \esp{Hay} s'emploie pour dire qu'une chose existe, sans préciser où elle se trouve.
\item On utilise \esp{estar} pour donner la localisation d'un lieu.
\item \esp{Ser} sert à décrire le caractère permanent d'un lieu.
\item En espagnol, l'adjectif se place toujours avant le nom, comme en français.
\end{enumerate}
\end{exercice}

\begin{exercice}[Vocabulaire : trouve le mot]
Complète chaque phrase avec le mot espagnol qui convient : \esp{turistas, hotel, paisaje, patrimonio, reserva natural, ecoturismo, guía, turismo}.
\begin{enumerate}
\item Los \dots\dots\dots\ visitan las playas de Kribi cada año.
\item Dormimos en un \dots\dots\dots\ de cuatro estrellas.
\item El \dots\dots\dots\ del monte Camerún es impresionante.
\item El \dots\dots\dots\ nos explicó la historia del monumento.
\item En el parque de Waza hay una \dots\dots\dots\ con elefantes y leones.
\item El \dots\dots\dots\ es una actividad económica muy importante.
\end{enumerate}
\end{exercice}

\begin{exercice}[Grammaire : hay, es ou está]
Complète avec \esp{hay}, \esp{es} ou \esp{está}.
\begin{enumerate}
\item En Douala \dots\dots\dots\ muchos hoteles modernos.
\item El monte Camerún \dots\dots\dots\ muy alto.
\item Kribi \dots\dots\dots\ en el sur del país.
\item La ciudad \dots\dots\dots\ tranquila y acogedora.
\item En la playa \dots\dots\dots\ palmeras y restaurantes.
\end{enumerate}
\end{exercice}

\courssub{Je m'entraîne}

\begin{exercice}[Thème : phrases à traduire]
Traduis en espagnol.
\begin{enumerate}
\item Il y a beaucoup de touristes à Kribi en décembre.
\item Le mont Cameroun est dans l'ouest du pays.
\item Le parc de Waza est magnifique et très grand.
\item Le guide est devant l'hôtel avec les touristes.
\item L'écotourisme est important pour le développement économique.
\end{enumerate}
\end{exercice}

\begin{exercice}[Transformation de phrases]
Transforme chaque phrase selon la consigne entre parenthèses.
\begin{enumerate}
\item \esp{El hotel es grande.} (mets au pluriel)
\item \esp{Hay una playa bonita.} (mets au pluriel)
\item \esp{El museo está en el centro.} (pose la question avec \esp{¿Dónde?})
\item \esp{La ciudad es turística.} (mets à la forme négative)
\item \esp{Hay muchos animales en la reserva.} (transforme en question avec \esp{¿Qué?})
\end{enumerate}
\end{exercice}

\begin{exercice}[Texte à trous : Bienvenidos a Kribi]
Complète le texte avec les mots suivants : \esp{hay, está, es, paisaje, turistas, playas, guía, patrimonio}.

\begin{textebox}[Bienvenidos a Kribi]
Kribi \dots\dots\dots\ en el sur de Camerún, a orillas del océano Atlántico. \dots\dots\dots\ una ciudad pequeña pero muy turística. En Kribi \dots\dots\dots\ playas magníficas de arena blanca. Las \dots\dots\dots\ más famosas son largas y tranquilas. Muchos \dots\dots\dots\ extranjeros visitan también las cataratas de la Lobé : el \dots\dots\dots\ es espectacular. Un \dots\dots\dots\ local puede explicar la historia de la región. Todo esto forma parte del \dots\dots\dots\ natural de Camerún.
\end{textebox}
\end{exercice}

\begin{exercice}[Appariements et réemploi]
\textbf{A.} Associe chaque mot espagnol à sa définition française.

\begin{center}
\begin{tabularx}{\linewidth}{|l|X|}
\hline
\rowcolor{popblueL} Español & Français \\
\hline
1. \esp{el paisaje} & a. personne qui accompagne les visiteurs \\
\hline
2. \esp{el guía} & b. espace naturel protégé \\
\hline
3. \esp{la reserva natural} & c. vue, aspect d'une région \\
\hline
4. \esp{el patrimonio} & d. tourisme respectueux de la nature \\
\hline
5. \esp{el ecoturismo} & e. richesses culturelles et historiques d'un pays \\
\hline
6. \esp{el turista} & f. personne qui voyage pour le plaisir \\
\hline
\end{tabularx}
\end{center}

\textbf{B.} Réemploie quatre de ces mots dans quatre phrases complètes en espagnol sur le Cameroun.
\end{exercice}

\courssub{Je me prépare au Bac}

\begin{exercice}[Type Bac — Compréhension écrite et grammaire en contexte]
Lis le texte suivant, puis réponds aux questions.

\begin{textebox}[El ecoturismo en Costa Rica]
Costa Rica es un pequeño país de América Central famoso por su naturaleza extraordinaria. Allí hay volcanes, selvas tropicales y playas magníficas. El país está entre el océano Pacífico y el mar Caribe. El ecoturismo es muy importante para su economía : miles de turistas visitan cada año las reservas naturales, donde viven monos, tucanes y perezosos. Los guías locales explican a los visitantes la riqueza de las plantas y de los animales. El gobierno protege los parques nacionales porque son el patrimonio de todos. Gracias al turismo responsable, muchas familias tienen trabajo y los pueblos se desarrollan. Sin embargo, es necesario controlar el número de visitantes para no destruir la naturaleza. El ecoturismo demuestra que el turismo y la protección del medio ambiente pueden ir juntos.
\end{textebox}

\textbf{I. Comprensión} (réponds en espagnol, avec des phrases complètes)
\begin{enumerate}
\item Di si es verdadero o falso y justifica con una cita del texto :
\begin{enumerate}
\item[a)] Costa Rica es un país muy grande.
\item[b)] El ecoturismo da trabajo a muchas familias.
\item[c)] El gobierno no protege los parques nacionales.
\end{enumerate}
\item ¿Dónde está Costa Rica ?
\item ¿Qué animales viven en las reservas naturales ?
\item ¿Por qué es necesario controlar el número de visitantes ?
\end{enumerate}

\textbf{II. Lengua}
\begin{enumerate}
\item Busca en el texto dos ejemplos de \esp{hay}, dos de \esp{ser} y un ejemplo de \esp{estar}, y explica cada uso.
\item Pon en plural : \esp{El guía local explica la riqueza de la planta.}
\end{enumerate}
\end{exercice}

\begin{exercice}[Type Bac — Grammaire en contexte : Kribi]
Complète le texte suivant avec \esp{hay}, \esp{es}, \esp{está} ou \esp{están} (10 occurrences).

\begin{textebox}[Kribi, ciudad turística]
Kribi (1) \dots\dots\dots\ una ciudad del sur de Camerún. (2) \dots\dots\dots\ a unos ciento cincuenta kilómetros de Duala. En Kribi (3) \dots\dots\dots\ playas de arena blanca y hoteles cómodos. La ciudad (4) \dots\dots\dots\ tranquila y acogedora. Las cataratas de la Lobé (5) \dots\dots\dots\ a pocos kilómetros del centro : (6) \dots\dots\dots\ un espectáculo natural único porque el río cae directamente en el mar. En el puerto (7) \dots\dots\dots\ barcos de pescadores. La gastronomía local (8) \dots\dots\dots\ deliciosa : (9) \dots\dots\dots\ pescado fresco en todos los restaurantes. Kribi (10) \dots\dots\dots\ un destino ideal para el turismo responsable.
\end{textebox}
\end{exercice}

\begin{attention}[Le piège du pluriel avec estar]
Quand le sujet est pluriel, \esp{estar} devient \esp{están} : \esp{Las cataratas \textbf{están} a pocos kilómetros.} « Les chutes sont à quelques kilomètres. » En revanche, \esp{hay} est \textbf{invariable} : \esp{Hay una playa} / \esp{Hay muchas playas}.
\end{attention}

\begin{exercice}[Type Bac — Thème, version et expression écrite]
\textbf{A. Thème.} Traduis en espagnol les phrases suivantes.
\begin{enumerate}
\item À Waza, il y a des éléphants et des lions.
\item Le parc est dans le nord du Cameroun.
\item C'est un lieu magnifique qu'il faut protéger.
\item Les touristes visitent le parc avec un guide.
\item Le tourisme durable est important pour l'économie du pays.
\end{enumerate}

\textbf{B. Version.} Traduis en français le paragraphe suivant.

\begin{textebox}[Turismo sostenible]
El turismo sostenible es una forma de viajar que respeta la naturaleza y ayuda a la población local. Los turistas duermen en pequeños hoteles de la región, compran productos artesanales y no destruyen el medio ambiente. Así, el dinero del turismo queda en el pueblo y las familias pueden vivir mejor. Este modelo de turismo es el futuro para muchos países.
\end{textebox}

\textbf{C. Expression écrite.} Tu écris pour un blog de voyage espagnol. Rédige un texte de \textbf{10 à 12 lignes} (environ 100 mots) dans lequel tu :
\begin{itemize}
\item présentes un site touristique de ta région (nom, situation, description avec \esp{hay}, \esp{ser}, \esp{estar} et des adjectifs) ;
\item expliques ce que les touristes peuvent y voir et y faire ;
\item termines par un appel au tourisme responsable.
\end{itemize}
\end{exercice}

\begin{methode}[Réussir la présentation d'un site touristique]
\begin{enumerate}
\item \textbf{Nommer et situer} le lieu avec \esp{estar} : \esp{Kribi está en el sur de Camerún.}
\item \textbf{Décrire} son caractère avec \esp{ser} + adjectifs : \esp{Es una ciudad pequeña y acogedora.}
\item \textbf{Énumérer} ses richesses avec \esp{hay} : \esp{Hay playas de arena blanca y hoteles cómodos.}
\item \textbf{Dire ce qu'on peut faire} : \esp{Los turistas pueden nadar, pescar y visitar las cataratas.}
\item \textbf{Conclure} par un appel au tourisme durable : \esp{¡Ven a Kribi y respeta su patrimonio natural!}
\end{enumerate}
\end{methode}

\newpage

\coursec{Corrigés des exercices}

\coursec{Corrigés des exercices — Leçon EA13 : El turismo y el desarrollo económico}

\begin{corrige}[Exercice 1 — QCM : le tourisme et son vocabulaire]
\begin{enumerate}
\item \textbf{b)} \esp{Un guía acompaña y explica a los turistas.} « Un guide accompagne les touristes et leur donne des explications. »
\item \textbf{b)} \esp{El ecoturismo respeta el medio ambiente y las culturas locales.} « L'écotourisme respecte l'environnement et les cultures locales. »
\item \textbf{a)} \esp{Una reserva natural es un espacio protegido donde viven animales y plantas.} « Une réserve naturelle est un espace protégé où vivent des animaux et des plantes. »
\item \textbf{b)} \esp{El patrimonio es el conjunto de las riquezas culturales e históricas de un país.} « Le patrimoine est l'ensemble des richesses culturelles et historiques d'un pays. »
\end{enumerate}
\textbf{Points de vigilance :} attention au faux ami \esp{guía} : il désigne la personne (le guide) mais aussi le livre (le guide touristique) ; il est masculin même pour une femme : \esp{la guía}. Ne pas confondre \esp{turismo} (l'activité) et \esp{turista} (la personne).
\end{corrige}

\begin{corrige}[Exercice 2 — Vrai ou faux ? Justifie]
\begin{enumerate}
\item \textbf{Vrai.} \esp{Hay} exprime l'existence (il y a), sans préciser la localisation : \esp{En Kribi hay playas magníficas.} « À Kribi, il y a des plages magnifiques. »
\item \textbf{Vrai.} \esp{Estar} indique la localisation d'un lieu : \esp{Kribi está en el sur del país.} « Kribi est dans le sud du pays. »
\item \textbf{Vrai.} \esp{Ser} décrit une caractéristique permanente : \esp{La ciudad es tranquila.} « La ville est tranquille. »
\item \textbf{Faux.} En espagnol, l'adjectif qualificatif se place le plus souvent \emph{après} le nom dans la description neutre : \esp{una playa bonita} « une belle plage ». La place antéposée (\esp{una bonita playa}) est possible dans un contexte subjectif, littéraire ou affectif, mais ce n'est pas la norme descriptive standard.
\end{enumerate}
\textbf{Points de vigilance :} retenir le triptyque de la description d'un lieu : \esp{hay} = existence ; \esp{estar} = localisation ; \esp{ser} = caractéristique permanente.
\end{corrige}

\begin{corrige}[Exercice 3 — Vocabulaire : trouve le mot]
\begin{enumerate}
\item \esp{Los \textbf{turistas} visitan las playas de Kribi cada año.} « Les touristes visitent les plages de Kribi chaque année. »
\item \esp{Dormimos en un \textbf{hotel} de cuatro estrellas.} « Nous dormons dans un hôtel quatre étoiles. »
\item \esp{El \textbf{paisaje} del monte Camerún es impresionante.} « Le paysage du mont Cameroun est impressionnant. »
\item \esp{El \textbf{guía} nos explicó la historia del monumento.} « Le guide nous a expliqué l'histoire du monument. »
\item \esp{En el parque de Waza hay una \textbf{reserva natural} con elefantes y leones.} « Dans le parc de Waza, il y a une réserve naturelle avec des éléphants et des lions. »
\item \esp{El \textbf{turismo} es una actividad económica muy importante.} « Le tourisme est une activité économique très importante. »
\end{enumerate}
\textbf{Points de vigilance :} \esp{hotel} s'écrit sans accent (mot oxytonne terminé par \emph{l} : l'accent tonique est sur la dernière syllabe, pas de signe graphique) ; se prononce « o-tèl ». \esp{paisaje} se prononce « pa-ï-sa-ré » (diphthongue \emph{ai} + \emph{je} géodésique).
\end{corrige}

\begin{corrige}[Exercice 4 — Grammaire : hay, es ou está]
\begin{enumerate}
\item \esp{En Douala \textbf{hay} muchos hoteles modernos.} — existence (il y a).
\item \esp{El monte Camerún \textbf{es} muy alto.} — caractéristique permanente.
\item \esp{Kribi \textbf{está} en el sur del país.} — localisation.
\item \esp{La ciudad \textbf{es} tranquila y acogedora.} — caractéristique permanente.
\item \esp{En la playa \textbf{hay} palmeras y restaurantes.} — existence.
\end{enumerate}
\textbf{Points de vigilance :} devant un complément de lieu (\esp{en el sur}, \esp{a orillas del océano}), choisir \esp{estar} ; devant un adjectif qualificatif, choisir \esp{ser} ; devant un nom indéfini (\esp{muchos hoteles}, \esp{palmeras}), choisir \esp{hay}.
\end{corrige}

\begin{corrige}[Exercice 5 — Thème : phrases à traduire]
\begin{enumerate}
\item \esp{En diciembre hay muchos turistas en Kribi.} — \esp{hay} + nom pluriel indéfini ; les mois ne prennent pas de majuscule en espagnol.
\item \esp{El monte Camerún está en el oeste del país.} — localisation : \esp{estar} ; « dans l'ouest de » = \esp{en el oeste de}.
\item \esp{El parque de Waza es magnífico y muy grande.} — caractéristique : \esp{ser} ; adjectifs masculins singuliers, placés après le nom.
\item \esp{El guía está delante del hotel con los turistas.} — localisation d'une personne : \esp{estar} ; « devant » = \esp{delante de} + article contracté \esp{del} (\esp{de} + \esp{el}).
\item \esp{El ecoturismo es importante para el desarrollo económico.} — « pour » de destination = \esp{para}.
\end{enumerate}
\textbf{Points de vigilance :} ne pas traduire « il y a » par \esp{haber} conjugué autrement qu'à la 3\textsuperscript{e} personne du singulier : \esp{hay}, jamais \esp{han}. « C'est » en début de phrase descriptive se rend par \esp{es}.
\end{corrige}

\begin{corrige}[Exercice 6 — Transformation de phrases]
\begin{enumerate}
\item \esp{Los hoteles son grandes.} — article, nom, verbe et adjectif passent tous au pluriel ; \esp{ser} devient \esp{son}.
\item \esp{Hay unas playas bonitas.} — \esp{hay} reste invariable ; \esp{una} devient \esp{unas}, \esp{bonita} devient \esp{bonitas}.
\item \esp{¿Dónde está el museo?} — question d'interrogation totale sur le lieu : \esp{¿Dónde?} + verbe + sujet ; accent sur \esp{dónde} et point d'interrogation ouvrant \esp{¿}.
\item \esp{La ciudad no es turística.} — la négation se place simplement devant le verbe : \esp{no es}.
\item \esp{¿Qué hay en la reserva?} — \esp{¿Qué hay...?} « Qu'y a-t-il... ? » ; accent sur \esp{qué}.
\end{enumerate}
\textbf{Points de vigilance :} en espagnol, toute question s'ouvre par \esp{¿} et se ferme par \esp{?} ; les mots interrogatifs prennent toujours un accent : \esp{dónde, qué, cuándo, cómo}.
\end{corrige}

\begin{corrige}[Exercice 7 — Texte à trous : Bienvenidos a Kribi]
\esp{Kribi \textbf{está} en el sur de Camerún, a orillas del océano Atlántico. \textbf{Es} una ciudad pequeña pero muy turística. En Kribi \textbf{hay} playas magníficas de arena blanca. Las \textbf{playas} más famosas son largas y tranquilas. Muchos \textbf{turistas} extranjeros visitan también las cataratas de la Lobé : el \textbf{paisaje} es espectacular. Un \textbf{guía} local puede explicar la historia de la región. Todo esto forma parte del \textbf{patrimonio} natural de Camerún.}

Traduction : « Kribi est dans le sud du Cameroun, au bord de l'océan Atlantique. C'est une petite ville mais très touristique. À Kribi, il y a de magnifiques plages de sable blanc. Les plages les plus célèbres sont longues et tranquilles. Beaucoup de touristes étrangers visitent aussi les chutes de la Lobé : le paysage est spectaculaire. Un guide local peut expliquer l'histoire de la région. Tout cela fait partie du patrimoine naturel du Cameroun. »

\textbf{Points de vigilance :} \esp{está} (localisation de Kribi) ; \esp{Es} (identification : c'est une ville...) ; \esp{hay} (existence des plages) ; attention à l'accord : \esp{muchos turistas}, \esp{las playas más famosas}.
\end{corrige}

\begin{corrige}[Exercice 8 — Appariements et réemploi]
\textbf{A.} 1-c ; 2-a ; 3-b ; 4-e ; 5-d ; 6-f.

\textbf{B.} Exemples de phrases possibles (toute phrase correcte et en rapport avec le Cameroun est acceptée) :
\begin{itemize}
\item \esp{El paisaje del monte Camerún es impresionante.} « Le paysage du mont Cameroun est impressionnant. »
\item \esp{El guía explica la historia del monumento a los turistas.} « Le guide explique l'histoire du monument aux touristes. »
\item \esp{En el parque de Waza hay una reserva natural con elefantes.} « Dans le parc de Waza, il y a une réserve naturelle avec des éléphants. »
\item \esp{El ecoturismo protege el patrimonio natural de Camerún.} « L'écotourisme protège le patrimoine naturel du Cameroun. »
\end{itemize}
\textbf{Points de vigilance :} vérifier l'accord de l'adjectif avec le nom (\esp{patrimonio natural}, \esp{reserva natural}) et la présence des articles.
\end{corrige}

\begin{corrige}[Exercice 9 — Type Bac : Compréhension écrite et grammaire en contexte]
\textbf{I. Comprensión}
\begin{enumerate}
\item a) \esp{Falso. El texto dice : « Costa Rica es un pequeño país de América Central ».} \\
b) \esp{Verdadero. El texto dice : « Gracias al turismo responsable, muchas familias tienen trabajo ».} \\
c) \esp{Falso. El texto dice : « El gobierno protege los parques nacionales ».}
\item \esp{Costa Rica está en América Central, entre el océano Pacífico y el mar Caribe.} « Le Costa Rica se trouve en Amérique centrale, entre l'océan Pacifique et la mer des Caraïbes. »
\item \esp{En las reservas naturales viven monos, tucanes y perezosos.} « Dans les réserves naturelles vivent des singes, des toucans et des paresseux. »
\item \esp{Es necesario controlar el número de visitantes para no destruir la naturaleza.} « Il est nécessaire de contrôler le nombre de visiteurs pour ne pas détruire la nature. »
\end{enumerate}

\textbf{II. Lengua}
\begin{enumerate}
\item \esp{Hay} : \esp{« Allí hay volcanes, selvas tropicales y playas magníficas »} (existence). \esp{Ser} : \esp{« Costa Rica es un pequeño país »} (identification) ; \esp{« El ecoturismo es muy importante »} (caractéristique permanente). \esp{Estar} : \esp{« El país está entre el océano Pacífico y el mar Caribe »} (localisation).
\item \esp{Los guías locales explican la riqueza de las plantas.} — sujet, verbe et compléments au pluriel.
\end{enumerate}

\textbf{Barème indicatif (sur 20) :} Compréhension 10 pts (V/F justifié : 3 x 2 pts ; questions 2, 3, 4 : 2 + 1 + 1 pts) ; Lengua 6 pts (exemples et explications : 4 pts ; pluriel : 2 pts) ; qualité de la langue 4 pts.

\textbf{Points de vigilance :} pour le vrai/faux, la justification doit être une \emph{citation exacte} du texte, introduite par \esp{El texto dice : « ... »} ; répondre par des phrases complètes, pas par des mots isolés.
\end{corrige}

\begin{corrige}[Exercice 10 — Type Bac : Grammaire en contexte : Kribi]
(1) \esp{es} — identification (Kribi = une ville du sud). \\
(2) \esp{Está} — localisation (à 150 km de Douala). \\
(3) \esp{hay} — existence (des plages et des hôtels). \\
(4) \esp{es} — caractéristique permanente (tranquille et accueillante). \\
(5) \esp{están} — localisation, sujet pluriel : \esp{las cataratas}. \\
(6) \esp{es} — identification (un spectacle naturel unique). \\
(7) \esp{hay} — existence (des bateaux de pêcheurs). \\
(8) \esp{es} — caractéristique (délicieuse). \\
(9) \esp{hay} — existence (du poisson frais). \\
(10) \esp{es} — identification (une destination idéale).

\textbf{Barème indicatif (sur 10) :} 1 pt par réponse correcte.

\textbf{Points de vigilance :} le seul piège réel est le (5) : sujet pluriel \esp{las cataratas}, donc \esp{están} et non \esp{está}. Rappel : \esp{hay} est invariable, même devant un pluriel (\esp{hay barcos}).
\end{corrige}

\begin{corrige}[Exercice 11 — Type Bac : Thème, version et expression écrite]
\textbf{A. Thème.}
\begin{enumerate}
\item \esp{En Waza hay elefantes y leones.} — existence : \esp{hay} ; pas d'article après \esp{hay} devant un nom indéfini pluriel.
\item \esp{El parque está en el norte de Camerún.} — localisation : \esp{estar} ; « dans le nord de » = \esp{en el norte de}.
\item \esp{Es un lugar magnífico que hay que proteger.} — « c'est » = \esp{es} ; « il faut » = \esp{hay que} + infinitif.
\item \esp{Los turistas visitan el parque con un guía.} — présent de \esp{visitar}, 3\textsuperscript{e} personne du pluriel : \esp{visitan}.
\item \esp{El turismo sostenible es importante para la economía del país.} — « pour » de destination = \esp{para} ; \esp{del} = \esp{de} + \esp{el}.
\end{enumerate}

\textbf{B. Version.}
Texte original (espagnol) :
\esp{El turismo sostenible es una forma de viajar que respeta la naturaleza y ayuda a la población local. Los turistas duermen en pequeños hoteles de la región, compran productos artesanales y no destruyen el medio ambiente. Así, el dinero del turismo se queda en el pueblo y las familias pueden vivir mejor. Este modelo de turismo es el futuro de muchos países.}

Traduction française : « Le tourisme durable est une façon de voyager qui respecte la nature et aide la population locale. Les touristes dorment dans de petits hôtels de la région, achètent des produits artisanaux et ne détruisent pas l'environnement. Ainsi, l'argent du tourisme reste dans le village et les familles peuvent mieux vivre. Ce modèle de tourisme est l'avenir de beaucoup de pays. »

\textbf{C. Expression écrite — proposition de rédaction modèle.}

\esp{El monte Camerún está en el oeste del país, cerca de la ciudad de Buéa. Es un volcán alto y majestuoso, el punto más alto de África occidental. En la región hay senderos bonitos, cascadas impresionantes y pueblos acogedores. Los turistas pueden subir a la cumbre con un guía, observar aves raras y descubrir la cultura local. La gente de la región es muy amable con los visitantes. El paisaje es espectacular y el aire es puro. Este lugar forma parte del patrimonio natural de Camerún. ¡Ven al monte Camerún, disfruta de su belleza y respeta su naturaleza! El turismo responsable es el futuro de nuestro país.}

Traduction : « Le mont Cameroun se trouve dans l'ouest du pays, près de la ville de Buéa. C'est un volcan haut et majestueux, le point le plus haut d'Afrique occidentale. Dans la région, il y a de beaux sentiers, des cascades impressionnantes et des villages accueillants. Les touristes peuvent monter au sommet avec un guide, observer des oiseaux rares et découvrir la culture locale. Les gens de la région sont très gentils avec les visiteurs. Le paysage est spectaculaire et l'air est pur. Ce lieu fait partie du patrimoine naturel du Cameroun. Viens au mont Cameroun, profite de sa beauté et respecte sa nature ! Le tourisme responsable est l'avenir de notre pays. »

\textbf{Barème indicatif (sur 20) :} Thème 5 pts (1 pt par phrase) ; Version 5 pts (fidélité 3 pts, qualité du français 2 pts) ; Expression écrite 10 pts (respect des consignes et contenu 4 pts, correction grammaticale — \esp{hay}/\esp{ser}/\esp{estar}, accords — 4 pts, richesse du vocabulaire 2 pts).

\textbf{Points de vigilance :}
\begin{itemize}
\item Employer \esp{estar} pour situer, \esp{ser} pour décrire, \esp{hay} pour énumérer : c'est le cœur du barème.
\item Accorder les adjectifs : \esp{pueblos acogedores}, \esp{cascadas impresionantes}, \esp{aves raras}.
\item Placer l'adjectif après le nom : \esp{un volcán alto}, et non \esp{un alto volcán}.
\item Terminer par un impératif d'appel : \esp{¡Ven...!}, \esp{¡Visita...!}, avec les signes \esp{¡} et \esp{!}.
\item Soigner l'orthographe des noms propres : \esp{Buéa} (accent sur le \emph{e} pour marquer la diphtongue), \esp{Camerún}, \esp{Atlántico}.
\end{itemize}
\end{corrige}

\newpage

\coursec{Fiche bilan}

\begin{popfigure}
\begin{center}\tcbox[colback=popblue,colframe=popblue,coltext=white,arc=9pt,boxsep=4pt,left=6pt,right=6pt]{\bfseries\small El turismo y el desarrollo econ\'omico}\end{center}
\vspace{-2pt}
\begin{tcbraster}[raster columns=3,raster equal height=rows,raster column skip=6pt,raster row skip=6pt,enhanced,blankest]
\begin{tcolorbox}[enhanced,colback=popgreen,colframe=popgreen,coltext=white,boxrule=0.9pt,arc=3pt,left=4pt,right=4pt,top=3pt,bottom=3pt,boxsep=1pt,halign=left]\bfseries\scriptsize L\'exico: turismo, gu\'ia, patrimonio\end{tcolorbox}
\begin{tcolorbox}[enhanced,colback=poporange,colframe=poporange,coltext=white,boxrule=0.9pt,arc=3pt,left=4pt,right=4pt,top=3pt,bottom=3pt,boxsep=1pt,halign=left]\bfseries\scriptsize Describir: hay / ser / estar\end{tcolorbox}
\begin{tcolorbox}[enhanced,colback=poppurple,colframe=poppurple,coltext=white,boxrule=0.9pt,arc=3pt,left=4pt,right=4pt,top=3pt,bottom=3pt,boxsep=1pt,halign=left]\bfseries\scriptsize Adjetivos: g\'enero, n\'umero, posici\'on\end{tcolorbox}
\begin{tcolorbox}[enhanced,colback=poppink,colframe=poppink,coltext=white,boxrule=0.9pt,arc=3pt,left=4pt,right=4pt,top=3pt,bottom=3pt,boxsep=1pt,halign=left]\bfseries\scriptsize Potencial: Kribi, Waza, mont Cameroun\end{tcolorbox}
\begin{tcolorbox}[enhanced,colback=popgold,colframe=popgold,coltext=white,boxrule=0.9pt,arc=3pt,left=4pt,right=4pt,top=3pt,bottom=3pt,boxsep=1pt,halign=left]\bfseries\scriptsize Turismo sostenible: ecoturismo, reservas\end{tcolorbox}
\begin{tcolorbox}[enhanced,colback=popblue!70,colframe=popblue!70,coltext=white,boxrule=0.9pt,arc=3pt,left=4pt,right=4pt,top=3pt,bottom=3pt,boxsep=1pt,halign=left]\bfseries\scriptsize M\'etodo: comentario de texto\end{tcolorbox}
\begin{tcolorbox}[enhanced,colback=popgreen!70!black,colframe=popgreen!70!black,coltext=white,boxrule=0.9pt,arc=3pt,left=4pt,right=4pt,top=3pt,bottom=3pt,boxsep=1pt,halign=left]\bfseries\scriptsize Producci\'on: promover un sitio\end{tcolorbox}
\end{tcbraster}

\legende{Carte mentale de la leçon : le tourisme et le développement économique}
\end{popfigure}

\begin{bilan}
\textbf{1. Lexique clé (à connaître par cœur)}
\begin{itemize}
\item \esp{el turismo} : le tourisme ; \esp{el/la turista} : le/la touriste ; \esp{turístico/a} : touristique
\item \esp{el hotel} : l'hôtel ; \esp{el alojamiento} : l'hébergement ; \esp{la playa} : la plage
\item \esp{el paisaje} : le paysage ; \esp{el patrimonio} : le patrimoine ; \esp{el monumento} : le monument
\item \esp{la reserva natural} : la réserve naturelle ; \esp{el parque nacional} : le parc national
\item \esp{el ecoturismo} : l'écotourisme ; \esp{el turismo sostenible} : le tourisme durable
\item \esp{el/la guía} : le/la guide ; \esp{la visita guiada} : la visite guidée ; \esp{el folleto} : la brochure
\item \esp{visitar} : visiter ; \esp{descubrir} : découvrir ; \esp{proteger} : protéger ; \esp{promover} : promouvoir
\item \esp{la economía} : l'économie ; \esp{el desarrollo} : le développement ; \esp{el empleo} : l'emploi ; \esp{los ingresos} : les revenus
\end{itemize}

\textbf{2. Grammaire : décrire un lieu}
\begin{center}
\begin{tabularx}{\linewidth}{|l|X|X|}
\hline
\rowcolor{popblueL} \textbf{Outil} & \textbf{Emploi} & \textbf{Exemple} \\
\hline
\esp{hay} & existence (il y a), invariable & \esp{En Kribi hay playas magníficas.} (À Kribi, il y a des plages magnifiques.) \\
\hline
\esp{ser} & identité, nature, caractéristique permanente & \esp{El monte Camerún es un volcán.} (Le mont Cameroun est un volcan.) \\
\hline
\esp{estar} & localisation, état temporaire & \esp{Waza está en el norte del país.} (Waza est dans le nord du pays.) \\
\hline
adjectifs & accord en genre et en nombre avec le nom & \esp{un paisaje hermoso / unas playas hermosas} \\
\hline
\end{tabularx}
\end{center}

\textbf{3. Règles à retenir}
\begin{itemize}
\item \esp{hay} + nom : jamais d'article défini après (\esp{hay hoteles}, pas \esp{hay los hoteles}).
\item Localisation d'un lieu : toujours \esp{estar} (\esp{Kribi está en el sur}), jamais \esp{ser}.
\item L'adjectif s'accorde : \esp{el parque es grande} / \esp{la reserva es grande} / \esp{los paisajes son grandes}.
\item Adjectifs fréquents de description : \esp{hermoso, impresionante, famoso, tranquilo, salvaje, único, diverso}.
\item Pour promouvoir : impératif (\esp{¡Visite Camerún! ¡Descubra Kribi!}) et tournures \esp{se puede} + infinitif (\esp{Se puede observar elefantes en Waza}).
\end{itemize}

\textbf{4. Méthode express : commenter un texte sur le tourisme}
\begin{enumerate}
\item \textbf{Lire} deux fois ; souligner le thème et la thèse de l'auteur.
\item \textbf{Présenter} : type de texte, source, sujet, problématique (\esp{¿El turismo ayuda al desarrollo?}).
\item \textbf{Analyser} : idées principales, arguments, exemples (chiffres, lieux cités).
\item \textbf{Conclure} : bilan + opinion personnelle justifiée (\esp{En mi opinión... / Pienso que...}).
\end{enumerate}

\textbf{5. Méthode express : présenter et promouvoir un site}
\begin{enumerate}
\item Situer le lieu avec \esp{estar} ; le décrire avec \esp{ser} + adjectifs ; énumérer avec \esp{hay}.
\item Donner les activités possibles : \esp{se puede / se pueden} + infinitif.
\item Ajouter un argument de tourisme durable : \esp{Es importante proteger la naturaleza y respetar la cultura local.}
\item Terminer par une invitation à l'impératif : \esp{¡Venga a descubrir este paraíso!}
\end{enumerate}
\end{bilan}

\begin{aretenir}[Checklist avant le Bac]
\begin{itemize}
\item[$\square$] Je connais le lexique du tourisme : \esp{turista, paisaje, patrimonio, reserva natural, ecoturismo, guía}.
\item[$\square$] J'utilise \esp{hay} pour dire « il y a » sans jamais mettre d'article défini après.
\item[$\square$] Je choisis correctement entre \esp{ser} (nature permanente) et \esp{estar} (localisation, état).
\item[$\square$] Je sais situer un lieu : \esp{El parque de Waza está en el Extremo Norte de Camerún.}
\item[$\square$] J'accorde les adjectifs en genre et en nombre : \esp{playas hermosas, paisajes impresionantes}.
\item[$\square$] Je place correctement les accents : \esp{turístico, economía, región, Camerún}.
\item[$\square$] Je sais citer le potentiel touristique du Cameroun : Kribi, le mont Cameroun, Waza, les chutes de la Lobé.
\item[$\square$] Je connais des sites hispanophones : Machu Picchu (Perú), las pirámides de México, las playas de España, Malabo (Guinea Ecuatorial).
\item[$\square$] Je sais former l'impératif de politesse pour promouvoir : \esp{¡Visite! ¡Descubra! ¡Venga!}
\item[$\square$] J'utilise \esp{se puede} + infinitif pour présenter les activités d'un site.
\item[$\square$] Je sais structurer un commentaire : présentation, analyse, conclusion avec opinion.
\item[$\square$] Je peux rédiger un court texte de promotion du tourisme durable en 8 à 10 phrases.
\end{itemize}
\end{aretenir}

\findecours

\end{document}
